\documentclass[11pt]{amsart}

\usepackage[T1]{fontenc}
\usepackage[utf8]{inputenc}
\usepackage{lmodern}
\usepackage[a4paper,margin=1in]{geometry}
\usepackage{amsmath,amssymb,amsthm,mathtools,mathrsfs}
\usepackage{array,booktabs,longtable}
\usepackage{enumitem}
\usepackage{microtype}
\usepackage{xcolor}
\usepackage[colorlinks=true,linkcolor=blue!45!black,citecolor=blue!45!black,urlcolor=blue!45!black]{hyperref}
\usepackage[nameinlink,noabbrev]{cleveref}

\newtheorem{theorem}{Theorem}[section]
\newtheorem{proposition}[theorem]{Proposition}
\newtheorem{lemma}[theorem]{Lemma}
\newtheorem{corollary}[theorem]{Corollary}
\theoremstyle{definition}
\newtheorem{definition}[theorem]{Definition}
\theoremstyle{remark}
\newtheorem{remark}[theorem]{Remark}

\newcommand{\C}{\mathbf C}
\newcommand{\R}{\mathbf R}
\newcommand{\Q}{\mathbf Q}
\newcommand{\F}{\mathbf F}
\newcommand{\cO}{\mathcal O}
\newcommand{\p}{\mathfrak p}
\newcommand{\Tr}{\operatorname{Tr}}
\newcommand{\Nm}{\operatorname{N}}
\newcommand{\Gal}{\operatorname{Gal}}
\newcommand{\ph}{\operatorname{ph}}
\newcommand{\ord}{\operatorname{ord}}
\newcommand{\cond}{\operatorname{cond}}
\newcommand{\Res}{\operatorname{Res}}
\newcommand{\abs}[1]{\lvert #1\rvert}
\newcommand{\set}[1]{\left\{#1\right\}}
\newcommand{\Err}{\mathcal E}
\newcommand{\Aphase}{\mathcal A}
\newcommand{\HerbPhi}{\varphi^{\mathrm H}}
\newcommand{\HerbPsi}{\psi^{\mathrm H}}
\newcommand{\D}{\mathfrak D}
\newcommand{\eps}{\varepsilon}
\newcommand{\ceil}[1]{\left\lceil #1\right\rceil}
\newcommand{\floor}[1]{\left\lfloor #1\right\rfloor}

\title[The First Main Lemma for Local Epsilon Factors]
{A Local Proof of Langlands's First Main Lemma\\
for Local Epsilon Factors over Nonarchimedean Local Fields}
\author{Fukuhiro Ueda}
\address{Research Institute for Mathematical Sciences, Kyoto University, Kyoto 606--8502, Japan}
\email{fueda@kurims.kyoto-u.ac.jp}
\date{August 2026}
\subjclass[2020]{11S37, 11S40, 11R42}
\keywords{local epsilon factors, Artin $L$-functions, local constants,
Gauss sums, ramification, Lamprecht formula, cyclic extensions,
equal characteristic}
\hypersetup{
  pdftitle={A Local Proof of Langlands's First Main Lemma for Local Epsilon Factors over Nonarchimedean Local Fields},
  pdfauthor={Fukuhiro Ueda},
  pdfsubject={A local proof of Langlands's First Main Lemma for local epsilon factors in mixed and equal characteristic},
  pdfkeywords={local epsilon factors, Artin L-functions, Gauss sums, ramification, First Main Lemma}
}

\setcounter{tocdepth}{2}

\begin{document}
\begin{abstract}
We give a complete local proof of Langlands's First Main Lemma for
normalized one-dimensional local epsilon factors over every
nonarchimedean local field, in both mixed and equal characteristic.  The
proof reconstructs and completes the local method initiated by Dwork and
organized by Langlands; its decisive $p$-adic calculations and guiding
ramification-theoretic and stationary-phase ideas are due to them.  The
present paper unifies their normalizations and supplies omitted endpoint and
non-stable wild calculations. Exact norm--trace correction formulas yield the equal-characteristic
extension.  

The author is also completing a local
proof of Langlands's Second Main Lemma and will verify it in Lean.
Together with the present result and Langlands's remaining reductions,
that work will give a complete proof of the local construction of
epsilon factors envisioned by Langlands, without global functional
equations.  

Open AI ChatGPT assisted the drafting and English-language editing of this manuscript;
the author audited the complete text, checked the proof against the
completed Lean formalization, and assumes full responsibility.
\end{abstract}
\maketitle

\clearpage
\tableofcontents
\clearpage

\section{Introduction}\label{sec:introduction}

Let $F$ be a nonarchimedean local field with finite residue field and
Weil group $W_F$.  Langlands's local epsilon factors are determined by
their values on one-dimensional representations, multiplicativity, and
compatibility with induction.  Brauer induction reduces their
construction to local constants of quasi-characters, but it also creates
an independence problem: distinct Brauer decompositions must give the
same result.  In his notes on the functional equation of Artin
$L$-functions \cite{Langlands}, Langlands reduced this problem to four
explicit local identities.  The present paper gives a complete local
proof of the identity he called the First Main Lemma.

Let $K/F$ be cyclic of prime degree $\ell$.  For a continuous
quasi-character $\chi_F:F^\times\to\C^\times$ and a nontrivial
continuous additive character $\psi_F:F\to\C^\times$, put
\[
 \chi_K=\chi_F\circ N_{K/F},
 \qquad
 \psi_K=\psi_F\circ\Tr_{K/F}.
\]
Let $S(K/F)$ be the full group of continuous characters of $F^\times$
that are trivial on $N_{K/F}(K^\times)$. 

\begin{theorem}[First Main Lemma]\label{thm:first-main}
For every cyclic extension $K/F$ of prime degree, every continuous
quasi-character $\chi_F$ of $F^\times$, and every nontrivial continuous
additive character $\psi_F$ of $F$,
\begin{equation}\label{eq:first-main}
 \Delta_K(\chi_K,\psi_K)
 \prod_{\mu\in S(K/F)}\Delta_F(\mu,\psi_F)
 =
 \prod_{\mu\in S(K/F)}\Delta_F(\mu\chi_F,\psi_F).
\end{equation}
\end{theorem}

Here $\Delta_E(\theta,\psi)$ is Langlands's normalized local Gauss
factor, recalled in Section~\ref{sec:local-notation}.  This is the
identity quoted as Lemma 10.1 in Langlands, equation (12.1), p.~144, and
used again on p.~257; see \cite[pp.~144, 257]{Langlands}.

\subsection*{Historical background}
Dwork's 1954 thesis \cite{Dwork} contains the decisive $p$-adic
calculations behind the cyclic-prime relation, while the formulation of
the First Main Lemma and the ramification-theoretic and stationary-phase
organization are due to Langlands.  Dwork worked in mixed
characteristic, in different normalizations, and Langlands's local
manuscript was never completed: Chapter~10, which contained the
conclusion of the First Main Lemma, is missing.  After Deligne found a
shorter proof using global functional equations, Langlands stopped
preparing a complete version of the local argument; this history is
described in the electronic edition of his notes \cite{Langlands}.

\subsection*{The Second Main Lemma and the complete local construction}
The present theorem settles the First Main Lemma, but Langlands's
original local construction also requires the Second Main Lemma.  The
author is presently completing a companion local proof of the Second
Main Lemma and will formalize and verify the finished proof in Lean~4.
Once this is done, the Lean-verified First and Second Main Lemmas,
together with Langlands's remaining reductions and the standard
one-dimensional theory, will furnish a complete proof of the local
construction and well-definedness of epsilon factors envisioned in
\cite{Langlands}.  In particular, the proof of the construction will no
longer depend on a global functional equation or a globalization
argument.

\subsection*{Contributions of the present paper}
The proof is a reconstruction and completion of Dwork and Langlands
rather than an independent construction.  Nevertheless, several
ingredients required for a complete theorem are supplied here rather
than copied from either source.  We first translate Dwork's
mixed-characteristic calculations and Langlands's formulation into a
common system of normalizations and notation.  Langlands's Haar-integral
definition is converted into an exact finite sum, and its nonvanishing
is proved by Fourier duality, independently of Lamprecht's formula.

We then formulate stationary parameters as quotient classes, prove
their transport under norm before choosing representatives, and track
the phase change caused by every change of representative.  This yields
explicit high- and low-conductor tables in which the source
representatives, common denominators, and residual coordinates are
retained until their use.  The paper supplies the missing endpoint and
intermediate calculations, including the non-stable wild odd-prime
range and the complete wild quadratic calculation, and organizes the
unramified, tame, endpoint, stable, and non-stable cases into an
exhaustive dispatch.

Among the additional mathematical bridges are the critical norm
polynomial, the stable-to-critical trace comparison, the
affine-translation vanishing needed to pass from actual source rows to
normalized rows, and the branch matrix for characteristic-two critical
refinements.  Finally, exact norm--trace correction formulas isolate
the mixed-characteristic terms and show directly that they vanish in
equal characteristic.  Thus the equal-characteristic theorem is a
proved extension of the local calculation, not a formal assertion by
analogy.  The proof uses explicit local ramification and norm theory,
stationary phase, and finite-field harmonic analysis.  It uses neither
local class field theory nor a global functional equation.

\subsection*{Logical organization of the proof}
The argument is most transparent when the common reductions are made
before the individual cases.

First, Langlands's Haar-integral definition is converted into an exact
finite sum on $U_E^0/U_E^{m_E(\theta)}$.  The finite sum is proved
nonzero by Fourier duality, independently of Lamprecht's formula.  This
makes the local constant, the inverse-pairing identity, and the error
quotient available without a later nonvanishing dependency.

Second, the stationary parameter attached to a character is constructed
as a quotient class.  Representatives are chosen only when a later
formula needs one, and every representative change is accompanied by
the corresponding translation of the critical function.  This prevents
a stationary phase from being replaced silently by the phase of an
unrelated representative.

Third, the abstract transport of stationary classes under norm is
proved before the explicit high- and low-conductor tables.  The tables
then retain the actual source representatives, common denominators, and
residual coordinates until their final consumers.  A common
factor-separation identity isolates the admissible, elementary,
critical, residual, and correction terms in the First Main quotient.
The odd-prime and quadratic calculations differ only after this common
stage.

Finally, twist invariance permits the base character to be chosen with
minimal conductor in its $S(K/F)$-orbit.  The complete case tree is
\[
\begin{array}{c|l}
K/F & \text{case used in the proof}\\ \hline
\text{unramified} & \text{unramified comparison}\\
\text{ramified tame},\ \ell=2 & \text{tame quadratic calculation}\\
\text{ramified tame},\ \ell\text{ odd} & \text{tame odd calculation}\\
\text{wild},\ \ell=2 & \text{wild quadratic calculation}\\
\text{wild},\ \ell\text{ odd},\ m\le1
  & \text{endpoint conductors}\\
\text{wild},\ \ell\text{ odd},\ 2T\le m
  & \text{stable range}\\
\text{wild},\ \ell\text{ odd},\ m>1,\ m<2T
  & \text{non-stable odd-prime calculation}.
\end{array}
\]
Here $t$ is the unique lower ramification break in the ramified wild
case, $T=t+1$, and $m$ is the conductor of the chosen minimal-orbit
representative.   

\subsection*{Equal characteristic}
The local calculation is sufficiently exact to show where mixed
characteristic is actually used.  In the odd-prime non-stable case the
last correction has the closed form
\[
 X=\frac{pB}{1-n^{p-1}}.
\]
The denominator is proved nonzero.  In mixed characteristic the
numerator is controlled by a ramification-theoretic valuation estimate;
in equal characteristic $p$, the image of the integer $p$ in $F$ is
zero and hence $X=0$.

For a wild quadratic extension the corresponding exact correction is
\[
 X=\frac{2\Tr_{K/F}(u)}{1+N_{K/F}(u)}.
\]
Again the denominator is nonzero, and in characteristic two the
correction vanishes because $2=0$ in $F$.  The only remaining quadratic
boundary is excluded in equal characteristic by the fact that a
ramified Artin--Schreier quadratic extension has odd lower break.  Thus
the equal-characteristic theorem is not obtained by formally extending
a mixed-characteristic result: it follows from the exact local
correction formulas proved here.

\subsection*{Organization}
Section~\ref{sec:local-notation} establishes the two realizations of the
local constant, proves nonvanishing, introduces the error term, and
reduces the theorem to the displayed case tree.  The next four sections
develop the cyclic-prime ramification, stationary-phase, finite-field,
and stationary-parameter tools.  We then treat the unramified, tame,
endpoint, stable, odd non-stable, and quadratic cases.  The last section
contains only the exhaustive dispatch.
 
\section{Local constants and the elementary reduction}
\label{sec:local-notation}

\subsection{Local fields, characters, and exact conductors}

Let $E$ be a nonarchimedean local field.  Write
\[
 \cO_E=\{x\in E:v_E(x)\ge0\},
 \qquad
 \p_E=\{x\in E:v_E(x)\ge1\},
\]
and let $k_E=\cO_E/\p_E$ have cardinality $q_E$.  We normalize
$v_E(E^\times)=\mathbf Z$.  Put
\[
 U_E^0=\cO_E^\times,
 \qquad
 U_E^r=1+\p_E^r\quad(r\ge1).
\]
All quasi-characters are continuous homomorphisms
$E^\times\to\C^\times$, and all additive characters are continuous
homomorphisms from the additive group of $E$ to $\C^\times$.

\begin{remark}[Characteristic convention]
\label{rem:characteristic-convention}
Write $p=\operatorname{char}k_E$.  Thus either
$\operatorname{char}E=0$ with residue characteristic $p$, or
$\operatorname{char}E=p$.  An integer occurring in an identity inside
$E$ denotes its image under $\mathbf Z\to E$; in particular $p=0$ in
equal characteristic $p$.
\end{remark}

\begin{definition}[Multiplicative conductor]
The conductor exponent $m_E(\theta)$ is the least $m\ge0$ for which
$\theta$ is trivial on $U_E^m$.  Thus $m_E(\theta)=0$ precisely when
$\theta$ is unramified.
\end{definition}

\begin{definition}[Additive conductor]
For a nontrivial additive character $\psi$, the integer $n_E(\psi)$
is characterized by the condition that $\p_E^{-n_E(\psi)}$ is the
largest fractional ideal on which $\psi$ is trivial.
\end{definition}

These are Langlands's conventions; see \cite[p.~5]{Langlands}.

\subsection{The Haar-integral and finite realizations of the local constant}

Define the total phase
\[
 \ph(z)=
 \begin{cases}
 z/|z|,&z\ne0,\\
 1,&z=0.
 \end{cases}
 \qquad
 \Aphase(z):=\ph(z).
\]
Let $m=m_E(\theta)$ and $n=n_E(\psi)$, and choose
$\gamma\in E^\times$ such that
\begin{equation}\label{eq:admissible-gamma}
 \gamma\cO_E=\p_E^{m+n}.
\end{equation}
For a positive Haar measure $du$ on $U_E^0$, Langlands defines
\begin{equation}\label{eq:def-delta}
 \Delta_E(\theta,\psi)
 =\theta(\gamma)\,
  \ph\left(\int_{U_E^0}
      \psi(u/\gamma)\theta(u)^{-1}\,du\right).
\end{equation}

Put $Q_m=U_E^0/U_E^m$.  For $\bar u\in Q_m$, represented by
$u\in U_E^0$, define
\begin{equation}\label{eq:finite-local-gauss-summand}
 g_{\theta,\psi,\gamma}(\bar u)
 =\psi(u/\gamma)\theta(u)^{-1}
\end{equation}
and
\begin{equation}\label{eq:finite-local-gauss-sum}
 G_E(\theta,\psi;\gamma)
 =\sum_{\bar u\in Q_m}
   g_{\theta,\psi,\gamma}(\bar u).
\end{equation}

\begin{proposition}[Finite realization and nonvanishing]
\label{prop:finite-local-realization}
The summand in \eqref{eq:finite-local-gauss-summand} is well defined.
There is a positive real number $c_m$, depending only on the chosen Haar
measure and $m$, such that
\begin{equation}\label{eq:integral-finite-comparison}
 \int_{U_E^0}\psi(u/\gamma)\theta(u)^{-1}\,du
 =c_m G_E(\theta,\psi;\gamma).
\end{equation}
Moreover
\begin{equation}\label{eq:finite-local-gauss-nonzero}
 G_E(\theta,\psi;\gamma)\ne0.
\end{equation}
Consequently
\begin{equation}\label{eq:delta-finite-realization}
 \Delta_E(\theta,\psi)
 =\theta(\gamma)\ph\bigl(G_E(\theta,\psi;\gamma)\bigr),
\end{equation}
and the zero branch in the definition of $\ph$ is never used for a
local constant.
\end{proposition}

\begin{proof}
If $u/v\in U_E^m$, then $\theta(u)=\theta(v)$ and
$u-v\in\p_E^m$.  By \eqref{eq:admissible-gamma},
$(u-v)/\gamma\in\p_E^{-n}$, so
$\psi(u/\gamma)=\psi(v/\gamma)$.  Thus the summand descends to
$Q_m$.  The integrand is constant on every coset of $U_E^m$ in
$U_E^0$, and all such cosets have the same positive measure $c_m$;
this proves \eqref{eq:integral-finite-comparison}.

It remains to prove nonvanishing.  If $m=0$, the quotient $Q_0$ is a
singleton and $G_E(\theta,\psi;\gamma)=1$.
Assume $m>0$ and put $A=\cO_E/\p_E^m$.  The pairing
\[
 A\times A\longrightarrow\C^\times,
 \qquad
 (\bar x,\bar y)\longmapsto\psi(xy/\gamma)
\]
is perfect.  Indeed, its annihilator in either variable is exactly
$\p_E^m$: this follows from \eqref{eq:admissible-gamma} and the
maximality of $\p_E^{-n}$ in the definition of $n_E(\psi)$.

Extend the function $u\mapsto\theta(u)^{-1}$ on $A^\times$ by zero
to a function $f:A\to\C$, and write
\[
 \widehat f(x)=
 \sum_{u\in A^\times}\theta(u)^{-1}\psi(xu/\gamma).
\]
Then $\widehat f(1)=G_E(\theta,\psi;\gamma)$.  If $x\in A^\times$,
the substitution $u\mapsto x^{-1}u$ gives
\begin{equation}\label{eq:finite-fourier-unit-coefficient}
 \widehat f(x)=\theta(x)\widehat f(1).
\end{equation}
If $x$ is not a unit, exactness of the conductor provides
$a\in U_E^{m-1}$ with $\theta(a)\ne1$.  Since
$x(a-1)u/\gamma\in\p_E^{-n}$ for every $u\in U_E^0$, multiplication
by $a$ permutes $A^\times$ and gives
\[
 \widehat f(x)=\theta(a)^{-1}\widehat f(x),
\]
whence $\widehat f(x)=0$.  If $\widehat f(1)$ were zero,
\eqref{eq:finite-fourier-unit-coefficient} and the preceding vanishing
would make every Fourier coefficient of $f$ zero.  Perfect finite
Fourier duality would then give $f=0$, contradicting $f(1)=1$.
Thus \eqref{eq:finite-local-gauss-nonzero} holds.  Finally the positive
factor $c_m$ does not change phase, proving
\eqref{eq:delta-finite-realization}.
\end{proof}

\begin{corollary}[Nonvanishing of the local Gauss integral]
\label{cor:local-gauss-integral-nonzero}
For every quasi-character $\theta$ and every nontrivial additive
character $\psi$, the integral in \eqref{eq:def-delta} is nonzero.
\end{corollary}

\begin{lemma}[Independence of the admissible denominator]
\label{lem:gamma-independent}
The right side of \eqref{eq:def-delta} is independent of the choice of
$\gamma$ satisfying \eqref{eq:admissible-gamma}.
\end{lemma}

\begin{proof}
Let $\gamma'=a\gamma$ with $a\in U_E^0$.  Multiplication by $a$
permutes $Q_m$, and \eqref{eq:finite-local-gauss-sum} gives
\[
 G_E(\theta,\psi;\gamma')
 =\theta(a)^{-1}G_E(\theta,\psi;\gamma).
\]
The restriction of $\theta$ to $U_E^0$ factors through the finite group
$Q_m$, so $|\theta(a)|=1$.  By nonvanishing,
\[
 \ph\bigl(G_E(\theta,\psi;\gamma')\bigr)
 =\theta(a)^{-1}
  \ph\bigl(G_E(\theta,\psi;\gamma)\bigr).
\]
Since $\theta(\gamma')=\theta(a)\theta(\gamma)$, the two factors
cancel.
\end{proof}

The same comparison shows immediately that the value is independent of
the positive normalization of Haar measure.

For later use, let $k$ be a finite field, let
$\varphi:k\to\C^\times$ be nontrivial, and let
$\eta:k^\times\to\C^\times$ be multiplicative.  Define
\begin{equation}\label{eq:tau-definition}
 \tau_k(\eta,\varphi)
 =-\sum_{x\in k^\times}\eta(x)^{-1}\varphi(x).
\end{equation}
Thus $\tau_k(1,\varphi)=1$.

\begin{lemma}[Normalization of the residual additive character]
\label{lem:residual-normalization}
Let $\psi_F$ have conductor exponent $n$, and fix a uniformizer
$\varpi_F$.  The formula
\[
 \bar\psi(\bar x)
 =\psi_F(\widetilde x\varpi_F^{-n-1})
 \qquad(\bar x\in k_F)
\]
defines a nontrivial additive character of $k_F$, independently of the
lift $\widetilde x$ modulo $\p_F$.  If $\psi_{\rm can}$ is a
prescribed nontrivial additive character of $k_F$, there is a unit
$a\in\cO_F^\times$ such that the residual character attached to
$\psi_F^a(x)=\psi_F(ax)$ is $\psi_{\rm can}$.
\end{lemma}

\begin{proof}
Well-definedness follows because two lifts differ by an element of
$\p_F$, and therefore their arguments differ by an element of
$\p_F^{-n}$, on which $\psi_F$ is trivial.  Nontriviality follows
from the maximality of $\p_F^{-n}$.  Every additive character of the
finite field $k_F$ is uniquely of the form
\[
 x\longmapsto
 \exp\left(\frac{2\pi i}{p}
   \Tr_{k_F/\F_p}(cx)\right),
\]
because the trace pairing $k_F\times k_F\to\F_p$ is nondegenerate.
If $\bar\psi$ and $\psi_{\rm can}$ correspond to $c$ and $c_0$,
choose a unit whose residue is $c_0/c$.
\end{proof}

If $m_E(\theta)=0$ and $\varpi_E$ is a uniformizer, then
$\gamma=\varpi_E^{n_E(\psi)}$ is admissible and
\begin{equation}\label{eq:unramified-character-local-constant}
 \Delta_E(\theta,\psi)
 =\theta(\varpi_E)^{n_E(\psi)}.
\end{equation}
If $m_E(\theta)=1$, define
\[
 \bar\theta(\bar u)=\theta(u),
 \qquad
 \bar\psi(\bar x)
 =\psi(x/\varpi_E^{n_E(\psi)+1}).
\]
Then
\begin{equation}\label{eq:conductor-one-local-gauss}
 \Delta_E(\theta,\psi)
 =-\theta(\varpi_E)^{n_E(\psi)+1}
   \Aphase\bigl(\tau_{k_E}(\bar\theta,\bar\psi)\bigr).
\end{equation}
These are the exact local-to-residue-field formulas used in the
unramified, tame, and endpoint cases.

\subsection{Norm and trace, elementary identities, and the error term}

Let $K/F$ be a finite extension.  Write
\[
 N=N_{K/F}:K^\times\to F^\times,
 \qquad
 T=\Tr_{K/F}:K\to F.
\]
For a nontrivial additive character $\psi_F$ and a quasi-character
$\chi_F$, put
\[
 \psi_K=\psi_F\circ T,
 \qquad
 \chi_K=\chi_F\circ N.
\]
If $K/F$ is abelian, set
\begin{equation}\label{eq:def-S}
 S(K/F)=\{\mu:F^\times\to\C^\times:
   \mu\text{ continuous and }\mu\circ N=1\}.
\end{equation}
By the universal property of the quotient, this is the group of
continuous characters of the topological quotient
$F^\times/N(K^\times)$.  For cyclic $K/F$ of prime degree, the
quotient and $S(K/F)$ are proved directly to have order $\ell$ in
Corollary~\ref{cor:prime-norm-characters} below.  That calculation uses
the norm filtration and not local class field theory.

\begin{lemma}[Change of additive character]\label{lem:additive-scale}
For $a\in E^\times$, put $\psi^a(x)=\psi(ax)$.  Then
\begin{equation}\label{eq:additive-scale}
 \Delta_E(\theta,\psi^a)
 =\theta(a)\Delta_E(\theta,\psi).
\end{equation}
\end{lemma}

\begin{proof}
If $v_E(a)=r$, then $n_E(\psi^a)=n_E(\psi)+r$.  Hence $a\gamma$ is
admissible for $(\theta,\psi^a)$ whenever $\gamma$ is admissible for
$(\theta,\psi)$, and
$\psi^a(u/(a\gamma))=\psi(u/\gamma)$.
\end{proof}

\begin{corollary}[Scaling invariance of the First Main identity]
\label{cor:additive-independence}
For every $a\in F^\times$, the First Main Lemma for $\psi_F$ is
equivalent to the First Main Lemma for $\psi_F^a$.
\end{corollary}

\begin{proof}
By Corollary~\ref{cor:prime-norm-characters}, the group $S(K/F)$ has
$\ell$ elements.  The two sides acquire the respective factors
\[
 \chi_K(a)\prod_{\mu\in S(K/F)}\mu(a)
 \quad\text{and}\quad
 \chi_F(a)^\ell\prod_{\mu\in S(K/F)}\mu(a).
\]
Since $N_{K/F}(a)=a^\ell$ for $a\in F^\times$, these factors agree.
\end{proof}

\begin{lemma}[The trivial character]\label{lem:delta-one}
For every nontrivial additive character $\psi$,
\[
 \Delta_E(1,\psi)=1.
\]
\end{lemma}

\begin{proof}
Here $m_E(1)=0$.  Formula
\eqref{eq:unramified-character-local-constant} gives the result.
\end{proof}

\begin{lemma}[Inverse pairing]\label{lem:inverse-pair}
Let $\mu$ be a finite-order character of $E^\times$.  Then
\begin{equation}\label{eq:inverse-pair}
 \Delta_E(\mu,\psi)\Delta_E(\mu^{-1},\psi)=\mu(-1).
\end{equation}
\end{lemma}

\begin{proof}
Choose a common admissible $\gamma$ and write
\[
 I_\mu=\int_{U_E^0}\psi(u/\gamma)\mu(u)^{-1}\,du.
\]
Both $I_\mu$ and $I_{\mu^{-1}}$ are nonzero by
Corollary~\ref{cor:local-gauss-integral-nonzero}.  Since $\mu$ and
$\psi$ are unitary,
\[
 \overline{I_\mu}
 =\int_{U_E^0}\psi(-u/\gamma)\mu(u)\,du
 =\mu(-1)I_{\mu^{-1}}.
\]
The product of the two phases is therefore $\mu(-1)^{-1}=\mu(-1)$,
and the factors $\mu(\gamma)$ and $\mu^{-1}(\gamma)$ cancel.
\end{proof}

\begin{corollary}\label{cor:S-product-one}
If $K/F$ is cyclic of odd prime degree, then
\begin{equation}\label{eq:S-product-one}
 \prod_{\mu\in S(K/F)}\Delta_F(\mu,\psi_F)=1.
\end{equation}
\end{corollary}

\begin{proof}
The finiteness of $S(K/F)$ follows from
Corollary~\ref{cor:prime-norm-characters}.  For $x\in F^\times$,
viewed as an element of $K^\times$, one has
$N_{K/F}(x)=x^\ell$.  Hence, for every $\mu\in S(K/F)$,
\[
 \mu(x)^\ell=\mu(N_{K/F}(x))=1.
\]
Thus every nontrivial $\mu$ has order $\ell$, which is odd, and in
particular $\mu(-1)=1$.  Pair $\mu$ with $\mu^{-1}$ and use
Lemma~\ref{lem:delta-one} for the trivial character.
\end{proof}

\begin{definition}[Error term]
Define
\begin{equation}\label{eq:error-term}
 \Err_{K/F}(\chi_F)=
 \frac{\Delta_K(\chi_K,\psi_K)
       \prod_{\mu\in S(K/F)}\Delta_F(\mu,\psi_F)}
      {\prod_{\mu\in S(K/F)}\Delta_F(\mu\chi_F,\psi_F)}.
\end{equation}
The denominator is nonzero by
Corollary~\ref{cor:local-gauss-integral-nonzero}, and the First Main
Lemma is the assertion $\Err_{K/F}(\chi_F)=1$.
\end{definition}

\begin{lemma}[Twist invariance]\label{lem:twist-invariance}
For every $\nu\in S(K/F)$,
\[
 \Err_{K/F}(\nu\chi_F)=\Err_{K/F}(\chi_F).
\]
\end{lemma}

\begin{proof}
Since $\nu\circ N=1$, the character on $K^\times$ is unchanged.
Multiplication by $\nu$ permutes the full group $S(K/F)$, including its
trivial character, so the denominator is only reindexed.
\end{proof}

\subsection{Reduction to minimal norm-character orbits and the case tree}

\begin{proposition}[Minimal-orbit reduction]\label{prop:minimal-orbit-reduction}
To prove Theorem~\ref{thm:first-main}, it is enough, for each
$S(K/F)$-orbit of quasi-characters, to prove
$\Err_{K/F}(\chi_F)=1$ for one member of minimal multiplicative
conductor.  For such a member, the proof reduces to the following
mutually exhaustive branches:
\begin{enumerate}[label=\textup{(\roman*)}]
\item $K/F$ is unramified;
\item $K/F$ is ramified tame and $\ell=2$;
\item $K/F$ is ramified tame and $\ell$ is odd;
\item $K/F$ is wild and $\ell=2$;
\item $K/F$ is wild, $\ell$ is odd, and $m_F(\chi_F)\le1$;
\item $K/F$ is wild, $\ell$ is odd, $m_F(\chi_F)>1$, and
      $2(t+1)\le m_F(\chi_F)$;
\item $K/F$ is wild, $\ell$ is odd, $m_F(\chi_F)>1$, and
      $m_F(\chi_F)<2(t+1)$.
\end{enumerate}
Here $t$ is the unique lower break in the ramified wild case.
\end{proposition}

\begin{proof}
The group $S(K/F)$ is finite by
Corollary~\ref{cor:prime-norm-characters}, so every orbit contains a
member of minimal conductor.  Twist invariance transports the result
from that member to the entire orbit.  A ramified extension of prime degree is
totally ramified.  It is either tame, or wild with
$\ell=\operatorname{char}k_F$.  Since $\ell$ is prime, the degree is
quadratic or odd.  In the wild odd-degree case the conductor is either
at most one or greater than one; in the latter range the two final
inequalities are complementary.  Thus the list is exhaustive.
\end{proof}

\section{Ramification notation and the cyclic-prime norm theorem}

Let $K/F$ be a finite Galois extension of nonarchimedean local fields,
with Galois group $G$.  For an integer $i\geq -1$, the lower
ramification group is
\[
 G_i=\{\sigma\in G:\sigma(x)\equiv x\pmod{\p_K^{i+1}}
                      \text{ for every }x\in\cO_K\}.
\]
For a real $u\geq-1$, put $G_u=G_{\lceil u\rceil}$.  Define the
Herbrand function and its inverse by
\begin{equation}\label{eq:Herbrand-def}
 \HerbPhi_{K/F}(u)=\int_0^u\frac{ds}{[G_0:G_s]},
 \qquad
 \HerbPsi_{K/F}=(\HerbPhi_{K/F})^{-1}.
\end{equation}
The superscript ${\rm H}$ is used throughout to prevent confusion with
additive characters.  This notation avoids the conflict noted by Langlands on pp.~28 and 62.

When $K/F$ is cyclic of prime degree $\ell$ and ramified, it is totally
ramified and there is a unique integer $t\geq0$ such that
\begin{equation}\label{eq:one-break-new}
 G=G_0=\cdots=G_t,\qquad G_{t+1}=1.
\end{equation}
Consequently
\begin{equation}\label{eq:Herbrand-prime-new}
 \HerbPsi_{K/F}(u)=
 \begin{cases}
 u,&-1\leq u\leq t,\\
 t+\ell(u-t),&u\geq t.
 \end{cases}
\end{equation}
Indeed the integrand in \eqref{eq:Herbrand-def} is $1$ up to $t$ and
$1/\ell$ after $t$.  This is the cyclic-prime form of \cite[Lemma 6.3, p.~29]{Langlands}.

The only non-elementary ramification statement needed in this section
is the following graded norm theorem.  We state it separately in order
to avoid a circular argument based on correcting norms successively at
levels below the break.

\begin{theorem}[Graded norm theorem for a cyclic prime extension]
\label{thm:graded-norm}
Let $K/F$ be ramified cyclic of prime degree $\ell$, with lower break
$t$.  For $i\ge0$ write
\[
 \operatorname{gr}^iU_E=U_E^i/U_E^{i+1},
\]
where $\operatorname{gr}^0U_E=k_E^\times$ and, for $i\ge1$,
$\operatorname{gr}^iU_E$ is identified with the additive group of
$k_E$.  Then the norm induces maps
\[
 \overline N_i:\operatorname{gr}^iU_K\longrightarrow
                   \operatorname{gr}^iU_F\qquad(0\le i\le t)
\]
with the following properties.
\begin{enumerate}[label=(\roman*)]
\item For $0\le i<t$, $\overline N_i$ is an isomorphism.
\item At $i=t$, its kernel and cokernel both have order $\ell$.  Its
      kernel is the image of $G_t/G_{t+1}$ under the ramification map
\[
 \sigma\longmapsto
 \begin{cases}
  \overline{\sigma(\pi_K)/\pi_K},&t=0,\\[1mm]
  \overline{(\sigma(\pi_K)-\pi_K)/\pi_K^{t+1}},&t>0.
 \end{cases}
\]
\item For every integer $r>t$, put
      $a=\HerbPsi_{K/F}(r)$.  Norm induces an isomorphism
\[
 U_K^a/U_K^{a+1}\xrightarrow{\ \sim\ }
 U_F^r/U_F^{r+1}.
\]
\item The norm image is closed.  Consequently the graded
      isomorphisms in (iii) give surjectivity by successive approximation:
\[
 N(U_K^{\HerbPsi_{K/F}(r)})=U_F^r,
 \qquad
 N(U_K^{\HerbPsi_{K/F}(r)+1})=U_F^{r+1}.
\]
\end{enumerate}
This is the cyclic-prime specialization of
\cite[Chapter~V, \S3, Propositions~4--5 and Corollaries~1--3]{Serre}.
\end{theorem}

\begin{proof}
The graded maps and their kernels are those of
\cite[Chapter~V, \S3, Propositions~4--5]{Serre}; their surjectivity and
the resulting norm equalities are
\cite[Chapter~V, \S3, Corollaries~1--3]{Serre}.  For completeness, the
passage from the graded statement to (iv) is as follows.  Below the unique break the norm maps on successive unit
quotients are isomorphisms.  At the break, the ramification
homomorphism displayed in the statement identifies the kernel; source
and target have the same order, so the kernel and cokernel both have
order $\ell$.  Above the break, Herbrand's change of numbering gives
the isomorphism in (iii).  The subgroup $N(U_K^a)$ is compact and hence
closed in $U_F^0$.  Successive lifting through the graded isomorphisms
therefore converges to a norm and gives the two equalities in (iv).
\end{proof}


\begin{lemma}[Frobenius differences and trace-zero hyperplanes]
\label{lem:frobenius-difference-trace}
Let $k/\F_p$ be a finite field extension and let
$\lambda\in k^\times$.  Then the $\F_p$-linear maps
\[
 \wp(z)=z^p-z,
 \qquad
 \wp_\lambda(z)=z^p-\lambda^{p-1}z
\]
satisfy
\[
 \ker(\wp)=\F_p,
 \qquad
 \operatorname{im}(\wp)=\ker\!\left(\Tr_{k/\F_p}\right),
\]
and
\begin{equation}\label{eq:scaled-frobenius-trace-image}
 \ker(\wp_\lambda)=\lambda\F_p,
 \qquad
 \operatorname{im}(\wp_\lambda)
 =\lambda^p\ker\!\left(\Tr_{k/\F_p}\right).
\end{equation}
Under the perfect trace pairing $(a,x)\mapsto\Tr_{k/\F_p}(ax)$, the
annihilator of the last image is $\lambda^{-p}\F_p$.
\end{lemma}

\begin{proof}
The kernel of $\wp$ is the prime field.  Hence its image has
$\F_p$-dimension $[k:\F_p]-1$.  Frobenius invariance of the trace gives
$\Tr(z^p-z)=0$, so the image is contained in the trace-zero hyperplane;
the two spaces have the same dimension and therefore coincide.  Since
\[
 \wp_\lambda(\lambda y)=\lambda^p(y^p-y),
\]
the scaled kernel and image follow.  The orthogonal complement of
$\ker(\Tr_{k/\F_p})$ is $\F_p$: it contains $\F_p$, and both sides have
dimension one.  Scaling the second variable by $\lambda^p$ gives the
annihilator $\lambda^{-p}\F_p$.
\end{proof}

\begin{proposition}[Critical norm polynomial at a positive wild break]
\label{prop:critical-norm-polynomial}
Let $K/F$ be totally ramified cyclic of prime degree
$p=\operatorname{char}k_F$, with positive lower break $t$, and identify
the residue fields with $k$.
Choose a uniformizer $\varpi_K$, put
$\varpi_F=N_{K/F}(\varpi_K)$, choose a generator $\sigma$ of
$\Gal(K/F)$, and set
\[
 \lambda=
 \overline{\frac{\sigma(\varpi_K)-\varpi_K}
                   {\varpi_K^{t+1}}}\in k^\times.
\]
In the source coordinate $1+\varpi_K^t\widetilde z$ and the target
coordinate $1+\varpi_F^t\widetilde w$, the critical graded norm is the
additive polynomial
\begin{equation}\label{eq:critical-norm-polynomial-extracted}
 P_\lambda(Z)=Z^p-\lambda^{p-1}Z.
\end{equation}
Consequently
\[
 \ker P_\lambda=\lambda\F_p,
 \qquad
 \operatorname{im}P_\lambda
 =\lambda^p\ker\!\left(\Tr_{k/\F_p}\right),
\]
and its annihilator under the trace pairing is
$\lambda^{-p}\F_p$.  The linear coefficient is the normalized trace:
\begin{equation}\label{eq:critical-norm-trace-coefficient}
 \overline{\frac{\Tr_{K/F}(\varpi_K^t)}{\varpi_F^t}}
 =-\lambda^{p-1}.
\end{equation}
\end{proposition}

\begin{proof}
The degree-$p$ endpoint of the exact norm expansion is
\[
 N_{K/F}(\widetilde z\varpi_K^t)
 =\widetilde z^{\,p}\varpi_F^t,
\]
so the leading coefficient of the induced critical polynomial is one.
Its linear term is the trace term in the norm expansion, hence its
coefficient is
$\overline{\Tr_{K/F}(\varpi_K^t)/\varpi_F^t}$.  The graded norm theorem
shows that the induced map is additive and that its kernel is the
ramification line $\lambda\F_p$.  Hence the monic additive polynomial is
\[
 \prod_{j\in\F_p}(Z-j\lambda)
 =Z^p-\lambda^{p-1}Z.
\]
Comparison of the linear coefficients gives
\eqref{eq:critical-norm-trace-coefficient}; the remaining assertions are
Lemma~\ref{lem:frobenius-difference-trace}.
\end{proof}

\begin{lemma}[Strong wild symmetric-coefficient estimate]
\label{lem:wild-symmetric-bound}
Assume that $K/F$ is wildly ramified cyclic of prime degree
$\ell=p=\operatorname{char}k_F$, with lower break $t$.  Put
\[
 T=t+1,\qquad D=v_K(\mathfrak D_{K/F})=(p-1)T.
\]
For $x\in K$ let $s_j(x)$ be the $j$th elementary symmetric function
of its $p$ conjugates.  If $v_K(x)\ge a$, then, for
$1\le j\le p-1$,
\begin{equation}\label{eq:symmetric-bound}
 v_F(s_j(x))\ge
 \left\lfloor\frac{ja+(p-1)T}{p}\right\rfloor.
\end{equation}
For the last coefficient,
\begin{equation}\label{eq:symmetric-bound-last}
 v_F(s_p(x))=v_F(N_{K/F}x)=v_K(x)\ge a.
\end{equation}
\end{lemma}

\begin{proof}
We first recall the trace-ideal equality in the present normalization.
For every integer $q$,
\begin{equation}\label{eq:trace-ideal-preliminary}
 \Tr_{K/F}(\p_K^q)
 =\p_F^{\left\lfloor(q+D)/p\right\rfloor}.
\end{equation}
Indeed, trace duality identifies the inverse different with
$\p_K^{-D}$.  The inclusion
$\Tr(\p_K^q)\subseteq\p_F^r$ is equivalent to
$\Tr(\p_K^{q-pr})\subseteq\cO_F$, hence to
$q-pr\ge-D$.  The largest such $r$ is
$\lfloor(q+D)/p\rfloor$, proving
\eqref{eq:trace-ideal-preliminary}.

Put $P_i(x)=\Tr_{K/F}(x^i)$.  Since
$x^i\in\p_K^{ia}$, equation
\eqref{eq:trace-ideal-preliminary} gives
\begin{equation}\label{eq:power-sum-bound}
 v_F(P_i(x))\ge
 \left\lfloor\frac{ia+D}{p}\right\rfloor
 \qquad(i\ge1).
\end{equation}
We now use Newton's identities
\begin{equation}\label{eq:Newton-symmetric}
 j s_j(x)=\sum_{i=1}^j(-1)^{i-1}s_{j-i}(x)P_i(x),
 \qquad s_0(x)=1.
\end{equation}
Because $j<p$, the integer $j$ is a unit of $\cO_F$.  We prove
\eqref{eq:symmetric-bound} by induction on $j$.  For $j=1$ it is
\eqref{eq:power-sum-bound}.  Suppose it is known below $j$.
The term $i=j$ in \eqref{eq:Newton-symmetric} has the required
valuation by \eqref{eq:power-sum-bound}.  For $i<j$, the induction
hypothesis and \eqref{eq:power-sum-bound} give
\[
 \left\lfloor\frac{(j-i)a+D}{p}\right\rfloor
 +\left\lfloor\frac{ia+D}{p}\right\rfloor.
\]
For real numbers $A,B$ one has
$\lfloor A\rfloor+\lfloor B\rfloor\ge\lfloor A+B\rfloor-1$.
Here the sum before taking floors is
\[
 \frac{ja+D}{p}+\frac Dp.
\]
The extension is wild, so $T\ge2$ and
$D/p=(p-1)T/p\ge1$.  Hence the last displayed integral lower bound is
at least $\lfloor(ja+D)/p\rfloor$.  Every term in
\eqref{eq:Newton-symmetric} therefore has the required valuation.
Finally $s_p(x)=N_{K/F}(x)$, and total ramification gives
$v_F(Nx)=v_K(x)$.
\end{proof}

\begin{corollary}[Norm truncation at an explicit depth]
\label{cor:wild-norm-truncation}
Let $q,r\ge0$ and assume
\begin{equation}\label{eq:wild-norm-truncation-condition}
 2q+(p-1)T\ge pr.
\end{equation}
Then, for every $y\in\p_K^q$,
\begin{equation}\label{eq:wild-norm-truncation}
 N_{K/F}(1+y)
 \equiv1+\Tr_{K/F}(y)+N_{K/F}(y)
 \pmod{\p_F^r}.
\end{equation}
\end{corollary}

\begin{proof}
For $2\le j\le p-1$, Lemma~\ref{lem:wild-symmetric-bound} gives
\[
 v_F(s_j(y))\ge
 \left\lfloor\frac{jq+(p-1)T}{p}\right\rfloor
 \ge
 \left\lfloor\frac{2q+(p-1)T}{p}\right\rfloor\ge r.
\]
The exact norm expansion then gives
\eqref{eq:wild-norm-truncation}.
\end{proof}

\begin{proposition}[Stable-to-critical normalized trace]
\label{prop:stable-critical-trace}
Let $K/F$ be totally ramified cyclic of odd prime degree $p$, with
positive lower break $t$.  Choose $\varpi_K$, put
$\varpi_F=N_{K/F}(\varpi_K)$, and define $\lambda$ as in
Proposition~\ref{prop:critical-norm-polynomial}.  Suppose $d\ge t+1$
and put
\[
 d'=pd-\frac{p-1}{2}t.
\]
Then
\begin{equation}\label{eq:normalized-uniformizer-congruence}
 \frac{\varpi_F}{\varpi_K^p}
 \equiv1\pmod{\p_K^{t+1}},
\end{equation}
\begin{equation}\label{eq:stable-critical-power-congruence}
 \varpi_K^{2d'}
 \equiv \varpi_F^{2d-t}\varpi_K^t
 \pmod{\p_K^{2d'+1}},
\end{equation}
and
\begin{equation}\label{eq:stable-critical-trace-congruence}
 \frac{\Tr_{K/F}(\varpi_K^{2d'})}{\varpi_F^{2d}}
 \equiv
 \frac{\Tr_{K/F}(\varpi_K^t)}{\varpi_F^t}
 \pmod{\p_F}.
\end{equation}
The two normalized traces are units, and their common residue is the
critical norm coefficient
\begin{equation}\label{eq:stable-critical-trace-residue}
 \overline{\frac{\Tr_{K/F}(\varpi_K^{2d'})}{\varpi_F^{2d}}}
 =\overline{\frac{\Tr_{K/F}(\varpi_K^t)}{\varpi_F^t}}
 =-\lambda^{p-1}.
\end{equation}
\end{proposition}

\begin{proof}
For $j\in\F_p$, put
\[
 a_j=\frac{\sigma^j(\varpi_K)}{\varpi_K}-1.
\]
The ramification homomorphism gives $a_j\in\p_K^t$ and
$\overline{a_j/\varpi_K^t}=j\lambda$.  Products of at least two of the
$a_j$ lie in $\p_K^{t+1}$, and $\sum_{j\in\F_p}j=0$.  Hence
\[
 \frac{\varpi_F}{\varpi_K^p}
 =\prod_{j\in\F_p}(1+a_j)
 \equiv1\pmod{\p_K^{t+1}},
\]
which proves \eqref{eq:normalized-uniformizer-congruence}.  Since
$2d'=p(2d-t)+t$, raising the inverse ratio to the power $2d-t$ and
multiplying by $\varpi_K^t$ gives
\eqref{eq:stable-critical-power-congruence}.

The trace-ideal formula \eqref{eq:trace-ideal-preliminary} sends the difference in
\eqref{eq:stable-critical-power-congruence} into $\p_F^{2d+1}$.
After division by $\varpi_F^{2d}$ this is exactly
\eqref{eq:stable-critical-trace-congruence}.  In the critical source
coordinate the linear coefficient of the graded norm is
$\overline{\Tr(\varpi_K^t)/\varpi_F^t}$.  Proposition
\ref{prop:critical-norm-polynomial} identifies it with
$-\lambda^{p-1}$, which is nonzero.  This proves
\eqref{eq:stable-critical-trace-residue} and the unit assertions.
\end{proof}


\begin{theorem}[Cyclic-prime norm filtration]\label{thm:norm-filtration}
Let $K/F$ be ramified cyclic of prime degree $\ell$, with break $t$.
Put $d_{K/F}=v_K(\mathfrak D_{K/F})$.  Then:
\begin{enumerate}[label=(\alph*)]
\item $d_{K/F}=(\ell-1)(t+1)$.
\item For every integer $a$,
\begin{equation}\label{eq:trace-ideal}
 \Tr_{K/F}(\p_K^a)=
 \p_F^{\left\lfloor(a+d_{K/F})/\ell\right\rfloor}.
\end{equation}
\item Norm induces isomorphisms
\begin{equation}\label{eq:norm-below-break}
 K^\times/U_K^t\xrightarrow{\sim}F^\times/U_F^t,
 \qquad
 U_K^r/U_K^t\xrightarrow{\sim}U_F^r/U_F^t
 \quad(0\le r\le t).
\end{equation}
The map $U_K^t/U_K^{t+1}\to U_F^t/U_F^{t+1}$ has cokernel of order
$\ell$.
\item For every integer $r>t$,
\begin{equation}\label{eq:norm-above-break}
 N_{K/F}(U_K^{\HerbPsi_{K/F}(r)})=U_F^r,
 \qquad
 N_{K/F}(U_K^{\HerbPsi_{K/F}(r)+1})=U_F^{r+1}.
\end{equation}
\end{enumerate}
\end{theorem}

\begin{proof}
Choose a uniformizer $\pi_K$.  For every $\sigma\ne1$ one has
$v_K(\sigma(\pi_K)-\pi_K)=t+1$.  The different formula
\[
 v_K(\mathfrak D_{K/F})
 =\sum_{\sigma\ne1}v_K(\sigma(\pi_K)-\pi_K),
\]
which is \cite[Chapter~IV, \S1, Proposition~4]{Serre}, therefore gives
(a).

For (b), the trace-dual of $\p_K^a$ under the pairing
$(x,y)\mapsto\Tr_{K/F}(xy)$ is $\p_K^{-a-d_{K/F}}$ by
\cite[Chapter~III, \S3, Proposition~7]{Serre}.  Write
$\Tr(\p_K^a)=\p_F^c$.  An element $y\in F$ belongs to the dual of this
trace ideal if and only if
\[
 y\p_K^a\subseteq\p_K^{-d_{K/F}},
 \quad\text{equivalently}\quad
 \ell v_F(y)+a\ge-d_{K/F}.
\]
Thus the dual ideal is
$\p_F^{\lceil(-a-d_{K/F})/\ell\rceil}$, whereas the dual of
$\p_F^c$ is $\p_F^{-c}$.  Hence
\[
 -c=\left\lceil\frac{-a-d_{K/F}}{\ell}\right\rceil,
 \qquad
 c=\left\lfloor\frac{a+d_{K/F}}{\ell}\right\rfloor,
\]
which proves \eqref{eq:trace-ideal}.

For (c), apply parts (i) and (ii) of
Theorem~\ref{thm:graded-norm}.  Splicing the graded
isomorphisms for $0\le i<t$ gives the second map in
\eqref{eq:norm-below-break}.  The valuation map is preserved because
the residue degree is one: for $x\in K^\times$,
$v_F(Nx)=v_K(x)$.

If $t=0$ (the tame case), then
$K^\times/U_K^t=K^\times/U_K^0\simeq\mathbf Z$, and norm is the identity
on this valuation quotient.  If $t>0$ (the wild case), the residue norm
on $U_K^0/U_K^1=k_F^\times$ is $x\mapsto x^\ell=x^p$, hence is the
Frobenius automorphism; the remaining positive graded pieces below $t$
are isomorphisms by Theorem~\ref{thm:graded-norm}.  These two separate arguments give
the first map in \eqref{eq:norm-below-break}.  The order-$\ell$ cokernel
at the critical level is exactly Theorem~\ref{thm:graded-norm}(ii).

Part (d) is part (iv) of Theorem~\ref{thm:graded-norm}.
The successive-approximation step is worth spelling out.  Given
$u_0\in U_F^r$, choose $x_0$ whose norm has the same class modulo
$U_F^{r+1}$.  If
$u_0/N(x_0)\in U_F^{r+1}$, choose $x_1$ at the next Herbrand level,
and continue.  The partial products of the $x_i$ are Cauchy, their
norms converge to $u_0$, and continuity of norm gives the desired
preimage.  Closedness of the norm image ensures that no limit leaves
the image.
\end{proof}

\begin{corollary}[Different and additive conductors]\label{cor:additive-conductor}
Let $\psi_K=\psi_F\circ\Tr_{K/F}$.  If $K/F$ is totally ramified of
degree $\ell$, then
\begin{equation}\label{eq:additive-conductor-trace}
 n_K(\psi_K)=\ell n_F(\psi_F)+d_{K/F}.
\end{equation}
\end{corollary}

\begin{proof}
The largest fractional ideal $I$ on which $\psi_F\circ\Tr$ is trivial
is characterized by $\Tr(I)\subseteq\p_F^{-n_F(\psi_F)}$.  By the trace
duality used in the proof of \eqref{eq:trace-ideal}, this ideal is
$\p_K^{-\ell n_F(\psi_F)-d_{K/F}}$.
\end{proof}

\begin{corollary}[Ramified norm quotient and norm characters]
\label{cor:norm-characters}
Let $K/F$ be ramified cyclic of prime degree $\ell$.  Then
$F^\times/N(K^\times)$ is cyclic of order $\ell$, the group
$S(K/F)$ has $\ell$ elements, and every nontrivial
$\mu\in S(K/F)$ has conductor exponent $t+1$.
\end{corollary}

\begin{proof}
The first isomorphism in
Theorem~\ref{thm:norm-filtration}(c) shows that
\[
 F^\times=N(K^\times)U_F^t.
\]
Its injectivity also shows that
\[
 N(K^\times)\cap U_F^t=N(U_K^t).
\]
Consequently
\begin{equation}\label{eq:ramified-norm-quotient-critical}
 F^\times/N(K^\times)\simeq U_F^t/N(U_K^t).
\end{equation}
Taking $r=t+1$ in \eqref{eq:norm-above-break} gives
\[
 U_F^{t+1}=N(U_K^{\HerbPsi_{K/F}(t+1)})\subseteq N(U_K^t),
\]
because $\HerbPsi_{K/F}(t+1)>t$.  Hence
\[
 U_F^t/N(U_K^t)
 \simeq
 \operatorname{coker}\left(
 U_K^t/U_K^{t+1}\longrightarrow U_F^t/U_F^{t+1}
 \right).
\]
By Theorem~\ref{thm:norm-filtration}(c), this cokernel has order
$\ell$.  Thus $F^\times/N(K^\times)$ has order $\ell$, and is
cyclic because $\ell$ is prime.

Moreover, $N(K^\times)$ contains the open subgroup $U_F^{t+1}$, so
the quotient is finite and discrete.  The elements of $S(K/F)$ are
therefore exactly the complex characters of this cyclic quotient.  A
character is determined by the image of a generator, which may be any
$\ell$th root of unity.  Hence $S(K/F)$ has $\ell$ elements.

Finally, \eqref{eq:ramified-norm-quotient-critical} shows that
$U_F^t$ maps onto the norm quotient, whereas $U_F^{t+1}$ maps to the
identity.  A nontrivial norm character is therefore nontrivial on
$U_F^t$ and trivial on $U_F^{t+1}$.  Its conductor exponent is exactly
$t+1$.
\end{proof}

\begin{proposition}[Conductor of norm pullback]\label{prop:pullback-conductor}
Let $m=m_F(\chi_F)$ and $m'=m_K(\chi_F\circ N)$.
\begin{enumerate}[label=(\roman*)]
\item If $m>t+1$, then
\begin{equation}\label{eq:pullback-high}
 m'-1=\HerbPsi_{K/F}(m-1).
\end{equation}
\item If $m\leq t$, then $m'=m$.
\item If $m=t+1$, then $m'\leq m$; moreover $m'<m$ if and only if
there is a $\mu\in S(K/F)$ for which
$m_F(\mu\chi_F)<t+1$.
\end{enumerate}
\end{proposition}

\begin{proof}
For (i), put $a=\HerbPsi_{K/F}(m-1)$.  By
\eqref{eq:norm-above-break}, $N(U_K^{a+1})=U_F^m$ and
$N(U_K^a)=U_F^{m-1}$.  Since $m$ is the exact conductor of $\chi_F$,
the pullback is trivial on $U_K^{a+1}$ and nontrivial on $U_K^a$.
Thus $m'=a+1$, which is \eqref{eq:pullback-high}.

For (ii), first separate $m=0$: an unramified character is trivial on
$U_F^0$, hence its norm pullback is unramified and $m'=0$.  Now assume
$1\le m\le t$.  The character $\chi_F$ is nontrivial on
$U_F^{m-1}/U_F^m$.  The graded norm map is an isomorphism at every
level $<t$, so its pullback is nontrivial on
$U_K^{m-1}/U_K^m$ and trivial on $U_K^m$.  Hence $m'=m$.

For (iii), the pullback is always trivial on $U_K^{t+1}$, so
$m'\le t+1=m$.  It has smaller conductor precisely when the restriction
of $\chi_F$ to the image
\[
 I=\operatorname{im}\bigl(U_K^t/U_K^{t+1}\longrightarrow
                          U_F^t/U_F^{t+1}\bigr)
\]
is trivial.  By Theorem~\ref{thm:norm-filtration}(c), the quotient
$(U_F^t/U_F^{t+1})/I$ has order $\ell$.  The proof of
Corollary~\ref{cor:norm-characters} identifies this quotient with
$F^\times/N(K^\times)$, so its character group is exactly the
restriction of $S(K/F)$ to the last unit layer.  Therefore
$\chi_F|_I=1$ if and only if there is a unique restriction
$\mu|_{U_F^t/U_F^{t+1}}$ with
$(\mu\chi_F)|_{U_F^t/U_F^{t+1}}=1$, equivalently
$m_F(\mu\chi_F)<t+1$.
\end{proof}

This is \cite[Lemmas 8.8 and 8.12, pp.~73 and 78--79]{Langlands}.

\begin{lemma}[Unramified norm, trace, and conductor compatibility]
\label{lem:unramified-compatibility}
Let $K/F$ be an unramified extension of degree $f$, and let $\pi$ be a
uniformizer of $F$.  Then $\pi$ is also a uniformizer of $K$,
$[k_K:k_F]=f$, and the following statements hold.
\begin{enumerate}[label=(\alph*)]
\item Reduction commutes with trace and norm:
\begin{equation}\label{eq:unramified-residue-trace-norm}
 \overline{\Tr_{K/F}(x)}=\Tr_{k_K/k_F}(\bar x)
 \quad(x\in\cO_K),
 \qquad
 \overline{N_{K/F}(u)}=N_{k_K/k_F}(\bar u)
 \quad(u\in\cO_K^\times).
\end{equation}
\item For every integer $a$,
\begin{equation}\label{eq:unramified-trace-ideal}
 \Tr_{K/F}(\p_K^a)=\p_F^a.
\end{equation}
\item For every $r\geq0$,
\begin{equation}\label{eq:unramified-norm-units}
 N_{K/F}(U_K^r)=U_F^r,
 \qquad N_{K/F}(\pi)=\pi^f.
\end{equation}
\item If $\psi_K=\psi_F\circ\Tr_{K/F}$ and
$\chi_K=\chi_F\circ N_{K/F}$, then
\begin{equation}\label{eq:unramified-conductors}
 n_K(\psi_K)=n_F(\psi_F),
 \qquad m_K(\chi_K)=m_F(\chi_F).
\end{equation}
\item The quotient $F^\times/N_{K/F}(K^\times)$ is cyclic of order $f$,
generated by the class of $\pi$.
\end{enumerate}
\end{lemma}

\begin{proof}
Since $K/F$ is unramified, its ramification index is $1$.  Thus
$\p_K=\pi\cO_K$, and the residue extension has degree $f$.
The conjugates of an integral element reduce to the conjugates of its
residue class.  Taking the first and last coefficients of the
characteristic polynomial therefore gives
\eqref{eq:unramified-residue-trace-norm}.

We use two elementary finite-field facts.  The map
$\Tr_{k_K/k_F}$ is onto: if $q=|k_F|$, its polynomial
\[
 X+X^q+\cdots+X^{q^{f-1}}
\]
is nonzero and has degree $q^{f-1}<q^f=|k_K|$, so it cannot vanish on
all of $k_K$; a nonzero $k_F$-linear map to the one-dimensional space
$k_F$ is surjective.  The map $N_{k_K/k_F}:k_K^\times\to k_F^\times$
is onto because $k_K^\times$ is cyclic and
\[
 N_{k_K/k_F}(x)=x^{(q^f-1)/(q-1)}
\]
has image of order $q-1$.

Choose $x\in\cO_K$ whose residue has trace $1$.  Then
$\Tr_{K/F}(x)$ is a unit of $\cO_F$; after multiplying $x$ by its
inverse, we obtain an integral element of trace $1$.  Hence
$\Tr_{K/F}(\cO_K)=\cO_F$.  Since
$\p_K^a=\pi^a\cO_K$ for every $a\in\mathbf Z$, this proves
\eqref{eq:unramified-trace-ideal}.

For $r\geq1$ and $x\in\cO_K$, multiplication of the conjugates gives
\begin{equation}\label{eq:unramified-graded-norm}
 N_{K/F}(1+\pi^r x)
 \equiv 1+\pi^r\Tr_{K/F}(x)\pmod{\p_F^{r+1}}.
\end{equation}
Thus the map induced by norm on
$U_K^r/U_K^{r+1}\to U_F^r/U_F^{r+1}$ is the finite-field trace and is
surjective.  Starting with $v\in U_F^r$, choose a lift of its class,
divide $v$ by the norm of that lift, and repeat at levels
$r+1,r+2,\ldots$.  Completeness makes the product of the successive
lifts converge to an element of $U_K^r$ with norm $v$.  This proves
$N(U_K^r)=U_F^r$ for $r\geq1$.  Surjectivity of the finite-field norm,
followed by the case $r=1$, gives $N(U_K^0)=U_F^0$.  Finally, because
$\pi\in F$, its norm is $\pi^f$.  This proves
\eqref{eq:unramified-norm-units}.

Put $n=n_F(\psi_F)$.  By \eqref{eq:unramified-trace-ideal},
$\psi_F\circ\Tr$ is trivial on $\p_K^{-n}$ and is not trivial on
$\p_K^{-n-1}$, so $n_K(\psi_K)=n$.  By
\eqref{eq:unramified-norm-units}, for every $r\geq0$ the character
$\chi_F\circ N$ is trivial on $U_K^r$ if and only if $\chi_F$ is
trivial on $U_F^r$.  Hence the multiplicative conductors are equal.
This proves \eqref{eq:unramified-conductors}.

Every $y\in K^\times$ has the form $y=\pi^a u$ with
$u\in U_K^0$, so
\[
 N_{K/F}(K^\times)=\pi^{f\mathbf Z}U_F^0
\]
by \eqref{eq:unramified-norm-units}.  The last assertion follows.
\end{proof}

\begin{corollary}[Norm characters in cyclic prime degree]
\label{cor:prime-norm-characters}
Let $K/F$ be cyclic of prime degree $\ell$.  Then
$F^\times/N_{K/F}(K^\times)$ is cyclic of order $\ell$, and
$S(K/F)$ has $\ell$ elements.  In particular, $S(K/F)$ is finite.
\end{corollary}

\begin{proof}
If $K/F$ is ramified, this is
Corollary~\ref{cor:norm-characters}.  Suppose that it is unramified.
Lemma~\ref{lem:unramified-compatibility}(e) gives a cyclic norm
quotient of order $\ell$, and
\eqref{eq:unramified-norm-units} shows that the norm group contains the
open subgroup $U_F^0$.  Thus the quotient is finite and discrete.  Its
continuous complex characters are all of its abstract characters, and
a character is determined by the image of a generator.  There are
exactly $\ell$ possible images, the $\ell$th roots of unity.  Hence
$S(K/F)$ has $\ell$ elements.
\end{proof}

By Lemma~\ref{lem:twist-invariance}, we shall always replace $\chi_F$
by a member of minimal conductor in its $S(K/F)$-orbit.  For such a
representative, Proposition~\ref{prop:pullback-conductor}(iii) implies
\begin{equation}\label{eq:minimal-critical}
 m\leq t+1\quad\Longrightarrow\quad m'=m.
\end{equation}

\section{Lamprecht's formula, with proof}

Let $E$ be nonarchimedean, let $\theta$ be a quasi-character of
conductor $m=2d+\varepsilon>1$, $\varepsilon\in\{0,1\}$, and let
$\gamma\cO_E=\p_E^{m+n_E(\psi)}$.

\begin{lemma}[Finite stationary-phase duality]\label{lem:finite-duality}
Let $n=n_E(\psi)$, let $M,q,r\in\mathbf Z$ with $r\le q$, and let
$\Gamma\in E^\times$ satisfy
\[
 \Gamma\cO_E=\p_E^{M+n}.
\]
Then
\begin{equation}\label{eq:general-stationary-pairing}
 \p_E^{M-q}/\p_E^{M-r}\ \times\
 \p_E^r/\p_E^q
 \longrightarrow\C^\times,
 \qquad(\bar c,\bar x)\longmapsto\psi(cx/\Gamma),
\end{equation}
is a perfect pairing.
\end{lemma}

\begin{proof}
The pairing is well defined: changing $c$ by an element of
$\p_E^{M-r}$ or $x$ by an element of $\p_E^q$ changes $cx/\Gamma$ by
an element of $\p_E^{-n}$, on which $\psi$ is trivial.  The annihilator
of $\p_E^r/\p_E^q$ in the first factor is zero, because
\[
 c\p_E^r/\Gamma\subseteq\p_E^{-n}
 \quad\Longleftrightarrow\quad
 c\in\p_E^{M-r}.
\]
Similarly, the annihilator of the first factor in the second is zero,
because
\[
 \p_E^{M-q}x/\Gamma\subseteq\p_E^{-n}
 \quad\Longleftrightarrow\quad
 x\in\p_E^q.
\]
Thus the pairing is nondegenerate in both variables.  Both finite
additive groups have cardinality $q_E^{q-r}$, hence the induced maps to
the respective character groups are isomorphisms.
\end{proof}

\begin{lemma}[Stationary numerator class at a chosen depth]
\label{lem:stationary-class-existence}
Let $q=m_E(\theta)>1$, let $n=n_E(\psi)$, choose an integer $M$ and
an element $\Gamma\in E^\times$ such that
\[
 \Gamma\cO_E=\p_E^{M+n},
\]
and let $r$ satisfy
\[
 \ceil{q/2}\le r\le q-1.
\]
There is a unique class
\begin{equation}\label{eq:general-stationary-class}
 \operatorname{St}_{\Gamma,r}(\theta)
 \in \p_E^{M-q}/\p_E^{M-r}
\end{equation}
such that every representative $c$ of this class satisfies
\begin{equation}\label{eq:general-stationary-linearization}
 \theta(1+x)=\psi(cx/\Gamma)
 \qquad(x\in\p_E^r).
\end{equation}
Every representative $c$ has exact valuation
\begin{equation}\label{eq:general-stationary-valuation}
 v_E(c)=M-q.
\end{equation}
\end{lemma}

\begin{proof}
For $x,y\in\p_E^r$, the quotient
\[
 \frac{(1+x)(1+y)}{1+x+y}
 =1+\frac{xy}{1+x+y}
\]
belongs to $U_E^q$, because $2r\ge q$.  Hence
$x\mapsto\theta(1+x)$ is an additive character of
$\p_E^r/\p_E^q$.  Lemma~\ref{lem:finite-duality}, applied with the
present $M,q,r,\Gamma$, identifies the full character group of
$\p_E^r/\p_E^q$ with
$\p_E^{M-q}/\p_E^{M-r}$.  This gives the unique class
\eqref{eq:general-stationary-class} satisfying
\eqref{eq:general-stationary-linearization}.

Every representative belongs to $\p_E^{M-q}$.  If it belonged to
$\p_E^{M-q+1}$, then for every $x\in\p_E^{q-1}$ one would have
$cx/\Gamma\in\p_E^{-n}$.  Since $r\le q-1$, such $x$ lies in the
linearization domain, and
\eqref{eq:general-stationary-linearization} would make $\theta$ trivial
on $U_E^{q-1}$, contrary to the exact conductor $q$.  Thus
$v_E(c)=M-q$.  Finally $r\le q-1$ implies
$\p_E^{M-r}\subseteq\p_E^{M-q+1}$, so adding any element of the
ambiguity ideal cannot change this valuation.  Hence every
representative has valuation $M-q$.
\end{proof}

\begin{lemma}[Stationary-class calculus]
\label{lem:stationary-class-calculus}
Use the notation of Lemma~\ref{lem:stationary-class-existence}.
\begin{enumerate}[label=(\alph*)]
\item If $r\le s\le q-1$, restriction from $\p_E^r$ to $\p_E^s$
      sends $\operatorname{St}_{\Gamma,r}(\theta)$ to the class of the
      same representative modulo $\p_E^{M-s}$.
\item If $\Gamma'=a\Gamma$, then $ac$ represents the stationary class
      relative to $\Gamma'$ whenever $c$ represents the class relative
      to $\Gamma$.
\item Let $\theta_1,\theta_2$ be characters and suppose that, on a
      common depth $r$, coefficients $c_1,c_2$ satisfy
      \[
       \theta_i(1+x)=\psi(c_ix/\Gamma)\qquad(x\in\p_E^r).
      \]
      Then $c_1+c_2$ represents the restriction of
      $\theta_1\theta_2$ on that depth, and $jc_1$ represents the
      restriction of the integer power $\theta_1^j$.  These assertions
      are equalities of classes modulo $\p_E^{M-r}$.
\item Suppose $m_E(\theta_1)=m_E(\theta_2)=q$, take
      $\ceil{q/2}\le r\le q-1$, and choose representatives $c_1,c_2$
      relative to the same $\Gamma$.  Then
      \begin{equation}\label{eq:stationary-leading-cancellation}
       m_E(\theta_1\theta_2)<q
       \quad\Longleftrightarrow\quad
       c_1+c_2\in\p_E^{M-q+1}.
      \end{equation}
      If the conductor drops to $q'<q$ and $q'>1$, put
      $r'=\ceil{q'/2}$ and choose a representative $c'$ of the
      stationary class at the new depth $r'$ by
      Lemma~\ref{lem:stationary-class-existence}.  Then
      \begin{equation}\label{eq:dropped-class-old-layer}
       c'\equiv c_1+c_2\pmod{\p_E^{M-r}}.
      \end{equation}
      Thus the old sum records only the restriction of the newly
      defined class on the old layer; it is not automatically a
      representative at the new depth.  If $q'\le1$, there is no
      stationary layer: the conductor-zero or conductor-one formula
      must be used, and the old sum is not retained as a parameter.
\end{enumerate}
\end{lemma}

\begin{proof}
Parts (a) and (b) follow immediately from
\eqref{eq:general-stationary-linearization} and uniqueness.  Part (c)
follows by multiplying the two character identities; uniqueness at the
common depth identifies the resulting class modulo $\p_E^{M-r}$.

For (d), the product has conductor at most $q$.  It has conductor less
than $q$ precisely when it is trivial on $U_E^{q-1}$.  By part (c), its
restriction there is
\[
 x\longmapsto\psi((c_1+c_2)x/\Gamma).
\]
The annihilator of $\p_E^{q-1}$ is $\p_E^{M-q+1}$, which proves
\eqref{eq:stationary-leading-cancellation}.  If the new conductor is
$q'>1$, a representative $c'$ of its stationary class at depth $r'$
induces the product character on every deeper unit layer.  Hence $c'$
and $c_1+c_2$ induce the same additive character on $\p_E^r$, even
when $r\ge q'$ and that character is already trivial.  By the
annihilator computation in
\eqref{eq:general-stationary-pairing}, their difference belongs to
$\p_E^{M-r}$, which proves
\eqref{eq:dropped-class-old-layer}.  Thus a sum on the old stationary
layer is used at a lower conductor only through this restriction of the
newly defined class; it is not treated as a canonical representative
at the new depth.  When the new conductor is $0$ or $1$, the last
assertion of the statement is simply the definition of the applicable
endpoint formula.
\end{proof}

For an admissible element $\gamma$ in Lamprecht's formula one has
$M=q$.  Taking $q=2d+\varepsilon$ and $r=d+\varepsilon$ therefore gives
\[
 \operatorname{St}_{\gamma,d+\varepsilon}(\theta)
 \in\cO_E/\p_E^d,
\]
and \eqref{eq:general-stationary-valuation} says that every
representative of this class is a unit.

\begin{lemma}[Magnitude of a nondegenerate finite quadratic sum]
\label{lem:finite-quadratic-magnitude}
Let $V$ be a finite abelian group and let $q:V\to\C^\times$ satisfy
$|q(x)|=1$, $q(0)=1$, and
\[
 B(x,y)=\frac{q(x+y)}{q(x)q(y)}.
\]
Assume that $B$ is a nondegenerate bicharacter.  Then
\begin{equation}\label{eq:finite-quadratic-magnitude}
 \left|\sum_{x\in V}q(x)\right|^2=|V|.
\end{equation}
In particular, the sum is nonzero.
\end{lemma}

\begin{proof}
Put $S=\sum_xq(x)$.  Since $\overline{q(y)}=q(y)^{-1}$, the change of
variables $x=y+h$ gives
\[
 \begin{aligned}
 |S|^2
 &=\sum_{x,y\in V}q(x)\overline{q(y)}
   =\sum_{h,y\in V}q(y+h)q(y)^{-1}\\
 &=\sum_{h\in V}q(h)\sum_{y\in V}B(h,y).
 \end{aligned}
\]
For $h=0$ the inner sum is $|V|$.  For $h\ne0$, nondegeneracy says that
$B(h,\cdot)$ is a nontrivial character, whose sum is zero.  Hence only
$h=0$ survives and $|S|^2=|V|$.
\end{proof}

\begin{theorem}[Lamprecht]\label{thm:Lamprecht}
Put
\[
 \beta=\beta(\theta):=
 \operatorname{St}_{\gamma,d+\varepsilon}(\theta)
 \in\cO_E/\p_E^d.
\]
Equivalently, $\beta$ is the unique class satisfying
\begin{equation}\label{eq:Lamprecht-linearization}
 \theta(1+x)=\psi(\beta x/\gamma)
 \qquad(x\in\p_E^{d+\varepsilon}).
\end{equation}
Every representative of this class is a unit.  Fix any representative,
again denoted $\beta\in\cO_E^\times$.  If $\varepsilon=0$, then
\begin{equation}\label{eq:Lamprecht-even}
 \Delta_E(\theta,\psi)=
 \theta(\gamma)\psi(\beta/\gamma)\theta(\beta)^{-1}.
\end{equation}
If $\varepsilon=1$, choose $\delta\cO_E=\p_E^d$ and put
\begin{equation}\label{eq:Lamprecht-Hasse-function}
 \phi_\theta(\bar x)=
 \psi(\delta\beta x/\gamma)\theta(1+\delta x)^{-1},
 \qquad \bar x\in k_E.
\end{equation}
Then this function is well defined, nowhere zero, and
\begin{equation}\label{eq:Lamprecht-odd}
 \Delta_E(\theta,\psi)=
 \theta(\gamma)\psi(\beta/\gamma)\theta(\beta)^{-1}
 \Aphase\!\left(\sum_{x\in k_E}\phi_\theta(x)\right),
 \qquad \Aphase(z):=\ph(z).
\end{equation}
Moreover
\begin{equation}\label{eq:Hasse-functional-equation}
 \phi_\theta(x+y)=
 \phi_\theta(x)\phi_\theta(y)\psi_0(xy),
\end{equation}
where
\begin{equation}\label{eq:Lamprecht-residual-character}
 \psi_0(\bar z)=\psi(\beta\delta^2\widetilde z/\gamma)
\end{equation}
is a well-defined nontrivial additive character of $k_E$.  Consequently
\begin{equation}\label{eq:Lamprecht-Hasse-magnitude}
 \left|\sum_{x\in k_E}\phi_\theta(x)\right|=q_E^{1/2}.
\end{equation}
\end{theorem}

\begin{proof}
Apply Lemma~\ref{lem:stationary-class-existence} with
\[
 q=M=m=2d+\varepsilon,
 \qquad \Gamma=\gamma,
 \qquad r=d+\varepsilon.
\]
Since $M-q=0$ and $M-r=d$, it gives precisely the class
$\beta\in\cO_E/\p_E^d$ characterized by
\eqref{eq:Lamprecht-linearization};
\eqref{eq:general-stationary-valuation} gives $v_E(\beta)=0$ for every
representative.  This proves the first assertions of the theorem.

Put
\[
 I=\int_{U_E^0}\psi(u/\gamma)\theta(u)^{-1}\,du,
 \qquad A=\p_E^{d+\varepsilon}/\p_E^m.
\]
For each $x\in A$, substitute $u\mapsto u(1+x)$ in the integral.  Haar
measure on $U_E^0$ is translation invariant, and
\eqref{eq:Lamprecht-linearization} gives
\[
 I=\int_{U_E^0}\psi(u/\gamma)\theta(u)^{-1}
       \psi((u-\beta)x/\gamma)\,du.
\]
Average this identity over $x\in A$.  Character orthogonality and
Lemma~\ref{lem:finite-duality}, specialized to
$M=q=m$, $r=d+\varepsilon$, and $\Gamma=\gamma$, give
\begin{equation}\label{eq:Lamprecht-selector}
 \frac1{|A|}\sum_{x\in A}\psi((u-\beta)x/\gamma)
 =\begin{cases}
 1,&u-\beta\in\p_E^d,\\
 0,&u-\beta\notin\p_E^d.
 \end{cases}
\end{equation}
Hence only the compact open subset
$\beta(1+\p_E^d)$ survives.

Assume first that $m=2d$.  Write $u=\beta(1+z)$ with
$z\in\p_E^d$.  Since $z\in\p_E^d$, the linearization identity with the variable $z$ gives directly
$\theta(1+z)=\psi(\beta z/\gamma)$.  Therefore on the surviving set
\[
 \psi(u/\gamma)\theta(u)^{-1}
 =\psi(\beta/\gamma)\theta(\beta)^{-1}.
\]
The integral is this nonzero constant times the positive measure of the
surviving set.  Taking its phase proves \eqref{eq:Lamprecht-even}.

Now assume that $m=2d+1$.  Since $m>1$, one has $d\ge1$.  Partition
$\beta(1+\p_E^d)$ into its $q_E$ cosets modulo
$\beta(1+\p_E^{d+1})$.  They are represented by
$u_x=\beta(1+\delta\widetilde x)$, $\bar x\in k_E$.  If
$w\in\p_E^{d+1}$, then the quotient of the integrand at
$u_x(1+w)$ by its value at $u_x$ is
\[
 \psi(u_xw/\gamma)\theta(1+w)^{-1}.
\]
Linearization gives $\theta(1+w)=\psi(\beta w/\gamma)$, while
$(u_x-\beta)w\in\p_E^{2d+1}=\p_E^m$; hence the displayed quotient is
one.  Thus the integrand is constant on each of the $q_E$ cosets and
its value at $u_x$ is
\[
 \psi(\beta/\gamma)\theta(\beta)^{-1}\phi_\theta(\bar x).
\]
This proves \eqref{eq:Lamprecht-odd}, once the sum is known to be
nonzero.

Changing a lift of $\bar x$ changes $\delta x$ by an element of
$\p_E^{d+1}$, and the same cancellation proves that
$\phi_\theta$ is well defined.  It is nowhere zero because both factors
in its definition lie in $\C^\times$.

For lifts $x,y\in\cO_E$, put
\[
 w=\frac{\delta^2xy}{1+\delta(x+y)}.
\]
Then
\[
 (1+\delta x)(1+\delta y)=(1+\delta(x+y))(1+w).
\]
The additive linear factors in \eqref{eq:Lamprecht-Hasse-function}
cancel, so
\[
 \frac{\phi_\theta(x+y)}{\phi_\theta(x)\phi_\theta(y)}
 =\theta(1+w)=\psi(\beta w/\gamma).
\]
Since $w-\delta^2xy\in\p_E^{3d}$ and
$3d\ge2d+1=m$, the last expression equals
$\psi(\beta\delta^2xy/\gamma)$.  This proves
\eqref{eq:Hasse-functional-equation} and
\eqref{eq:Lamprecht-residual-character}.  The character $\psi_0$ is
well defined because its coefficient has valuation $-n_E(\psi)-1$;
it is nontrivial for the same reason.  The bicharacter
$(x,y)\mapsto\psi_0(xy)$ is nondegenerate: for $x\ne0$, multiplication
by $x$ is a bijection of $k_E$, so $y\mapsto\psi_0(xy)$ is nontrivial.
Lemma~\ref{lem:finite-quadratic-magnitude} now gives
\eqref{eq:Lamprecht-Hasse-magnitude}; in particular the finite sum and
the original integral are nonzero.
\end{proof}

This is \cite[Lemma 8.1, pp.~60--62, and Lemma 9.3, pp.~86--87]{Langlands}.

\begin{lemma}[Independence of the stationary representative]
\label{lem:complete-stationary-invariant}
The complete stationary expression in Lamprecht's formula is independent
of the representative of the class
$\beta\in\cO_E/\p_E^d$.  More precisely, if
$\beta'=\beta+h$ with $h\in\p_E^d$, then the right sides of
\eqref{eq:Lamprecht-even} and \eqref{eq:Lamprecht-odd}, respectively,
are unchanged when $\beta$ is replaced by $\beta'$ and the Hasse
function is formed with $\beta'$.
\end{lemma}

\begin{proof}
Suppose first that $m=2d$.  Since $\beta$ is a unit,
$z=h/\beta$ lies in $\p_E^d$, and the linearization formula gives
$\theta(1+z)=\psi(\beta z/\gamma)=\psi(h/\gamma)$.  Hence
\[
 \frac{\psi(\beta'/\gamma)\theta(\beta')^{-1}}
      {\psi(\beta/\gamma)\theta(\beta)^{-1}}
 =\psi(h/\gamma)\theta(1+h/\beta)^{-1}=1.
\]

Now let $m=2d+1$.  Choose $\delta\cO_E=\p_E^d$ and write
\[
 a=\frac{h}{\beta\delta}\in\cO_E.
\]
Let $\phi_\beta$ and $\phi_{\beta'}$ denote the Hasse functions formed
with $\beta$ and $\beta'$.  Directly from
\eqref{eq:Lamprecht-Hasse-function},
\begin{equation}\label{eq:Hasse-representative-change}
 \phi_{\beta'}(x)
 =\phi_\beta(x)\,\psi(\delta h\widetilde x/\gamma).
\end{equation}
The functional equation \eqref{eq:Hasse-functional-equation} gives
\[
 \phi_\beta(x+a)
 =\phi_\beta(x)\phi_\beta(a)
   \psi(\beta\delta^2a\widetilde x/\gamma).
\]
Since $\beta\delta a=h$, the last additive factor is the one in
\eqref{eq:Hasse-representative-change}.  Therefore
\[
 \phi_{\beta'}(x)=\phi_\beta(a)^{-1}\phi_\beta(x+a),
 \qquad
 \sum_{x\in k_E}\phi_{\beta'}(x)
 =\phi_\beta(a)^{-1}\sum_{x\in k_E}\phi_\beta(x).
\]
On the other hand,
\[
 \frac{\psi(\beta'/\gamma)\theta(\beta')^{-1}}
      {\psi(\beta/\gamma)\theta(\beta)^{-1}}
 =\psi(h/\gamma)\theta(1+\delta a)^{-1}
 =\phi_\beta(a).
\]
Because the Hasse sums are nonzero, taking phases shows that these two
multipliers cancel.  Thus the complete odd-conductor stationary factor
is also unchanged.
\end{proof}

For a chosen representative $\beta$ of the stationary class, set
\[
 \Delta_{2,E}(\theta)=
 \psi(\beta/\gamma)\theta(\beta)^{-1},
 \qquad
 \Delta_{3,E}(\theta)=
 \begin{cases}
 1,&m_E(\theta)\text{ even},\\
 \Aphase(\sum\phi_\theta),&m_E(\theta)\text{ odd}.
 \end{cases}
\]
Then
\begin{equation}\label{eq:Delta-factorization-new}
 \Delta_E(\theta,\psi)=\theta(\gamma)\Delta_{2,E}(\theta)
                                      \Delta_{3,E}(\theta).
\end{equation}
When the conductor is odd, $\Delta_{2,E}(\theta)$ and
$\Delta_{3,E}(\theta)$ separately depend on the chosen representative;
their product is independent by
Lemma~\ref{lem:complete-stationary-invariant}.  Every later comparison
therefore either keeps these two factors together or applies the
corresponding translation to the critical function.

\begin{lemma}[Exact stable twist]\label{lem:stable-twist-new}
If $m_E(\theta)=2d+\varepsilon>1$ and
$m_E(\nu)\leq d$, then, for any representative $\beta$ of the
stationary class of $\theta$,
\begin{equation}\label{eq:stable-twist-new}
 \Delta_E(\nu\theta,\psi)=
 \nu(\gamma/\beta)\Delta_E(\theta,\psi).
\end{equation}
The right-hand side is independent of the chosen representative.
\end{lemma}

\begin{proof}
The character $\nu$ is trivial on $1+\p_E^d$, so the stationary
classes of $\theta$ and $\nu\theta$ coincide; choose the same
representative $\beta$ for both.  Their Hasse functions are then
identical.  Taking the ratio in
\eqref{eq:Delta-factorization-new} gives
$\nu(\gamma)\nu(\beta)^{-1}$.  If
$\beta'\equiv\beta\pmod{\p_E^d}$, then
$\beta'/\beta\in U_E^d$ and $\nu(\beta'/\beta)=1$, proving the final
independence assertion.
\end{proof}

The formulation above is the form of \cite[Lemma 9.4, pp.~87--88]{Langlands} used below.

\begin{lemma}[Conductors and leading classes in a minimal norm-character orbit]
\label{lem:minimal-orbit-stationary}
Let $K/F$ be ramified cyclic of prime degree $\ell$, let $t$ be its
break, put $T=t+1$, and let $\tau$ be a nontrivial element of
$S(K/F)$.  Choose $\chi_F$ of minimal conductor in its
$S(K/F)$-orbit and put $m=m_F(\chi_F)$.  Then, for every nontrivial
power $\tau^j$,
\begin{equation}\label{eq:minimal-orbit-conductors}
 m_F(\tau^j\chi_F)=\max\{m,T\}.
\end{equation}
If $T>1$ and $m=T$, choose a common denominator $\Gamma$, write
\[
 \Gamma\cO_F=\p_F^{M+n_F(\psi_F)},
\]
choose a common depth
\[
 \ceil{T/2}\le r\le T-1,
\]
and let $a$ and $b$ represent the stationary restrictions of $\tau$
and $\chi_F$.  Then the leading
classes of $b+ja$ do not vanish:
\begin{equation}\label{eq:minimal-orbit-no-leading-cancellation}
 b+ja\notin\p_F^{M-T+1}.
\end{equation}
Thus $b+ja$ represents a character of exact conductor $T$; no
lower-conductor Lamprecht parameter is obtained by reusing a cancelled
sum on the old layer.
\end{lemma}

\begin{proof}
A nontrivial norm character has conductor $T$ by
Corollary~\ref{cor:norm-characters}.  If $m<T$, then $\chi_F$ is trivial
on $U_F^{T-1}$, whereas $\tau^j$ is not, so the product has conductor
$T$.  If $m>T$, then $\tau^j$ is trivial on $U_F^{m-1}$, whereas
$\chi_F$ is not, so the product has conductor $m$.  If $m=T$, every
member of the orbit has conductor at least $T$ by the minimal choice of
$\chi_F$, and at most $T$ because both factors are trivial on $U_F^T$.
This proves \eqref{eq:minimal-orbit-conductors}.

Assume $T>1$ and $m=T$.  Lemma~\ref{lem:stationary-class-calculus}(c)
identifies $b+ja$ with the stationary restriction of
$\tau^j\chi_F$ on the common layer.  If its leading class vanished,
equation \eqref{eq:stationary-leading-cancellation} would give
conductor less than $T$, contradicting
\eqref{eq:minimal-orbit-conductors}.  This proves
\eqref{eq:minimal-orbit-no-leading-cancellation} and the final
assertion.
\end{proof}

\section{Finite-field identities}

Let $k$ be a finite field.  We use the normalized Gauss sum
$\tau_k(\chi,\psi)$ defined in \eqref{eq:tau-definition}; in
particular $\tau_k(1,\psi)=1$.  The local conductor-zero and
conductor-one formulas were already isolated in
\eqref{eq:unramified-character-local-constant} and
\eqref{eq:conductor-one-local-gauss}.  We now establish the finite-field
identities used to compare those formulas and the higher stationary phases.

\begin{lemma}[Finite-field Gauss-sum magnitude]
\label{lem:finite-gauss-magnitude}
For
\[
 G_k(\chi,\psi)=\sum_{x\in k^\times}\chi(x)\psi(x)
\]
one has $G_k(1,\psi)=-1$.  If $\chi$ is nontrivial, then
\begin{equation}\label{eq:finite-gauss-magnitude}
 |G_k(\chi,\psi)|^2=|k|.
\end{equation}
Consequently every $G_k(\chi,\psi)$ and every
$\tau_k(\chi,\psi)$ is nonzero.
\end{lemma}

\begin{proof}
The trivial-character formula is
$\sum_{x\in k^\times}\psi(x)=-1$, because the sum over all of $k$ is
zero.  Assume $\chi$ nontrivial and write $q=|k|$.  Since the values of
$\chi$ and $\psi$ have absolute value one,
\[
 \begin{aligned}
 |G_k(\chi,\psi)|^2
 &=\sum_{x,y\in k^\times}\chi(x/y)\psi(x-y)\\
 &=\sum_{t\in k^\times}\chi(t)
       \sum_{y\in k^\times}\psi((t-1)y).
 \end{aligned}
\]
The inner sum is $q-1$ when $t=1$ and is $-1$ when $t\ne1$.
Therefore
\[
 |G_k(\chi,\psi)|^2
 =(q-1)-\sum_{t\ne1}\chi(t).
\]
Character orthogonality gives $\sum_{t\in k^\times}\chi(t)=0$, so
$\sum_{t\ne1}\chi(t)=-1$ and the last expression is $q$.
\end{proof}


\begin{theorem}[Hasse--Davenport lift]\label{thm:HD-lift}
For a finite extension $k'/k$ of degree $f$,
\begin{equation}\label{eq:HD-lift}
 \tau_{k'}(\chi\circ N_{k'/k},\psi\circ\Tr_{k'/k})
 =\tau_k(\chi,\psi)^f.
\end{equation}
\end{theorem}

This is the classical lifting relation of Davenport and Hasse
\cite{DavenportHasse}.  We include the Euler-product proof in the
normalization used here.

\begin{proof}
It is convenient first to use the ordinary Gauss sum
\[
 G_k(\alpha,\psi)=\sum_{x\in k^\times}\alpha(x)\psi(x).
\]
For a monic polynomial
\[
 P(X)=X^r+a_1X^{r-1}+\cdots+a_r\in k[X],\qquad a_r\ne0,
\]
put
\[
 w(P)=\chi((-1)^r a_r)\psi(-a_1).
\]
If $P,Q$ are monic with nonzero constant terms, then the constant terms
multiply and the negative of the coefficient of degree
$\deg(PQ)-1$ is the sum of the roots; hence $w(PQ)=w(P)w(Q)$.  Therefore
\begin{align*}
 L(T)
 &=\prod_{\substack{P\ {\rm monic\ irreducible}\\ P\ne X}}
       (1-w(P)T^{\deg P})^{-1}\\
 &=\sum_{\substack{P\ {\rm monic}\\ P(0)\ne0}}w(P)T^{\deg P}.
\end{align*}
For degree $r\ge2$, the coefficient is zero, because $a_1$ ranges freely
through $k$ and $\sum_{a_1\in k}\psi(-a_1)=0$.  The coefficient in
degree $1$ is
\[
 \sum_{a_1\ne0}\chi(-a_1)\psi(-a_1)=G_k(\chi,\psi).
\]
Thus
\begin{equation}\label{eq:HD-L}
 L(T)=1+G_k(\chi,\psi)T.
\end{equation}
On the other hand, taking logarithms of the Euler product and grouping
closed points after extension of scalars gives
\begin{equation}\label{eq:HD-log}
 \log L(T)=\sum_{r\ge1}\frac{T^r}{r}
 \sum_{x\in k_r^\times}
     \chi(N_{k_r/k}x)\psi(\Tr_{k_r/k}x),
\end{equation}
where $k_r/k$ is the degree-$r$ extension.  Indeed, an element of $k_r$
of degree $d\mid r$ contributes through its monic minimal polynomial
with multiplicity $d$, exactly as in the usual Euler-product proof for
the zeta function of the affine line.  Comparing \eqref{eq:HD-L} and
\eqref{eq:HD-log} yields
\begin{equation}\label{eq:HD-G}
 G_{k_r}(\chi\circ N,\psi\circ\Tr)
 =(-1)^{r-1}G_k(\chi,\psi)^r.
\end{equation}
Now apply \eqref{eq:HD-G} to $\chi^{-1}$.  Since
$\tau_k(\chi,\psi)=-G_k(\chi^{-1},\psi)$, the signs cancel:
\[
 \tau_{k_r}(\chi\circ N,\psi\circ\Tr)
 =-(-1)^{r-1}G_k(\chi^{-1},\psi)^r
 =\tau_k(\chi,\psi)^r.
\]
Taking $r=f$ proves the theorem.
\end{proof}

\begin{theorem}[Hasse--Davenport multiplication formula]
\label{thm:HD-product}
Let $\ell\mid |k^\times|$, let $\eta$ have order $\ell$, and let
$\chi$ be any multiplicative character.  Then
\begin{equation}\label{eq:HD-product}
 \chi(\ell^\ell)\tau_k(\chi^\ell,\psi)
 \prod_{j=1}^{\ell-1}\tau_k(\eta^j,\psi)
 =\tau_k(\chi,\psi)
 \prod_{j=1}^{\ell-1}\tau_k(\chi\eta^j,\psi).
\end{equation}
Equivalently, for
$G(\alpha)=\sum_{x\in k^\times}\alpha(x)\psi(x)$ with $G(1)=-1$,
\begin{equation}\label{eq:ordinary-HD-product}
 \prod_{j=0}^{\ell-1}G(\alpha\eta^j)
 =\alpha(\ell^{-\ell})G(\alpha^\ell)
   \prod_{j=1}^{\ell-1}G(\eta^j).
\end{equation}
\end{theorem}

\begin{proof}
The ordinary identity \eqref{eq:ordinary-HD-product} is the classical
Davenport--Hasse multiplication formula \cite{DavenportHasse}.  Apply
it with
$\alpha=\chi^{-1}$.  Inversion permutes the nontrivial powers of
$\eta$, and $G(\alpha)=-\tau_k(\alpha^{-1},\psi)$.  Both sides contain
exactly $\ell$ ordinary Gauss sums when the factor
$G(\chi^{-\ell})$ is counted, so the signs cancel.  Finally
$\chi^{-1}(\ell^{-\ell})=\chi(\ell^\ell)$, which gives
\eqref{eq:HD-product} in Langlands's normalization.
\end{proof}

\begin{lemma}[Quadratic phase]\label{lem:quadratic-phase-new}
Assume $\operatorname{char}k\neq2$.  Let $\nu$ be the quadratic
character and put
\[
 Q_\psi(a,b)=\Aphase\!\left(\sum_{x\in k}
     \psi\left(\frac a2x^2+bx\right)\right),\qquad a\neq0.
\]
Then
\begin{equation}\label{eq:quadratic-phase-new}
 Q_\psi(a,b)=\nu(a)\psi(-b^2/2a)Q_\psi(1,0),
 \qquad Q_\psi(1,0)^2=\nu(-1).
\end{equation}
The defining sum has absolute value $|k|^{1/2}$.
\end{lemma}

\begin{proof}
Let
\[
 S(a,b)=\sum_{x\in k}\psi\left(\frac a2x^2+bx\right).
\]
Completing the square gives
\[
 S(a,b)=\psi(-b^2/2a)S(a,0).
\]
For $c\ne0$, every $y\in k^\times$ has $1+\nu(y)$ square roots, while
zero has one square root.  Hence
\begin{align}
 \sum_{x\in k}\psi(cx^2)
 &=\sum_{y\in k}(1+\nu(y))\psi(cy)\notag\\
 &=\sum_{y\in k^\times}\nu(y)\psi(cy)
  =\nu(c)G_k(\nu,\psi).\label{eq:quad-gauss}
\end{align}
The additive-character part disappeared because
$\sum_{y\in k}\psi(cy)=0$.  Taking $c=a/2$ and comparing with $a=1$
proves the first phase identity.  Lemma~\ref{lem:finite-gauss-magnitude}
shows at once that $|S(a,b)|=|k|^{1/2}$.

It remains to compute the square of the phase.  Put
$G=G_k(\nu,\psi)$.  Changing variables $x=ty$ in the product of the
two Gauss sums gives
\[
 G^2=\sum_{t\in k^\times}\nu(t)
       \sum_{y\in k^\times}\psi((t+1)y).
\]
Here the inner sum is $|k|-1$ for $t=-1$ and is $-1$ for
$t\ne-1$.  Therefore
\[
 \begin{aligned}
 G^2
 &=\nu(-1)(|k|-1)-\sum_{t\ne-1}\nu(t)\\
 &=\nu(-1)(|k|-1)+\nu(-1)
  =\nu(-1)|k|,
 \end{aligned}
\]
because $\sum_{t\in k^\times}\nu(t)=0$.  Since
$S(1,0)=\nu(1/2)G$ and $\nu(1/2)^2=1$, dividing by absolute values
gives $Q_\psi(1,0)^2=\nu(-1)$.
\end{proof}

\begin{theorem}[Hasse-function lift]\label{thm:Hasse-lift}
Let $k'/k$ be finite and let a nowhere-zero function $\phi:k\to\C^\times$
satisfy
\[
 \phi(x+y)=\phi(x)\phi(y)\psi(xy).
\]
Put
\[
 \omega_{k'/k}(x)=\sum_{\{\sigma,\tau\}}x^\sigma x^\tau,
 \qquad
 \phi_{k'/k}(x)=
 \phi(\Tr x)\psi(-\omega_{k'/k}(x)).
\]
Here the sum defining $\omega_{k'/k}$ is over the unordered pairs of
distinct $k$-embeddings $k'\hookrightarrow\overline{k}$.  Then
$\phi_{k'/k}$ satisfies the analogous functional equation over
$k'$ and
\begin{equation}\label{eq:Hasse-lift-sum}
 -\sum_{x\in k'}\phi_{k'/k}(x)
 =\left(-\sum_{x\in k}\phi(x)\right)^{[k':k]}.
\end{equation}
\end{theorem}

\begin{proof}
The identity
\[
 \omega(x+y)=\omega(x)+\omega(y)+\Tr(x)\Tr(y)-\Tr(xy)
\]
together with the functional equation for $\phi$ proves the corresponding
functional equation for $\phi_{k'/k}$.

For the sum identity, let
\[
 P(X)=X^r+a_1X^{r-1}+a_2X^{r-2}+\cdots
\]
be monic and put
\[
 w(P)=\phi(-a_1)\psi(-a_2),
\]
where $a_2=0$ in degree $1$.  If $Q$ has first coefficients $b_1,b_2$,
then the first two coefficients of $PQ$ are
$a_1+b_1$ and $a_2+b_2+a_1b_1$.  Hence
\begin{align*}
 w(PQ)
 &=\phi(-a_1-b_1)\psi(-a_2-b_2-a_1b_1)\\
 &=\phi(-a_1)\phi(-b_1)\psi(a_1b_1)
   \psi(-a_2-b_2-a_1b_1)=w(P)w(Q).
\end{align*}
Thus
\[
 L(T)=\prod_{P\ {\rm monic\ irreducible}}
       (1-w(P)T^{\deg P})^{-1}
     =\sum_{P\ {\rm monic}}w(P)T^{\deg P}.
\]
For degree $r\ge2$, the coefficient is zero because $a_2$ ranges freely
and $\psi$ is nontrivial.  In degree $1$ it is $\sum_{x\in k}\phi(x)$.
Therefore
\begin{equation}\label{eq:Hasse-L}
 L(T)=1+\left(\sum_{x\in k}\phi(x)\right)T.
\end{equation}
We next justify the closed-point expansion, including compatibility
with repeated closed points.  Let $P$ be monic irreducible of degree
$d$, choose a root $\xi\in k_d$, and put
\[
 s=\Tr_{k_d/k}(\xi),\qquad e=\omega_{k_d/k}(\xi).
\]
The first two coefficients of $P$ are $-s$ and $e$, so
\[
 w(P)=\phi(s)\psi(-e).
\]
If $r=hd$ and $\xi$ is viewed in $k_r$, each of its $d$ conjugates
occurs $h$ times.  Consequently
\begin{equation}\label{eq:Hasse-repeated-trace-omega}
 \Tr_{k_r/k}(\xi)=hs,
 \qquad
 \omega_{k_r/k}(\xi)=he+\binom h2s^2.
\end{equation}
Iteration of the functional equation for $\phi$ gives
\begin{equation}\label{eq:Hasse-iterated-functional-equation}
 \phi(hs)=\phi(s)^h\psi\!\left(\binom h2s^2\right).
\end{equation}
Combining the last two displays yields
\begin{equation}\label{eq:Hasse-closed-point-compatibility}
 \phi(\Tr_{k_r/k}\xi)\psi(-\omega_{k_r/k}(\xi))
 =\bigl(\phi(s)\psi(-e)\bigr)^h=w(P)^h.
\end{equation}
Thus the $d$ roots belonging to the closed point $P$ contribute
$d\,w(P)^h$ to the sum over $k_r$.  Expanding the logarithm of the
Euler product and regrouping by $r=h\deg P$ therefore gives
\begin{equation}\label{eq:Hasse-log}
 \log L(T)
 =\sum_{r\ge1}\frac{T^r}{r}
 \sum_{x\in k_r}\phi(\Tr_{k_r/k}x)
              \psi(-\omega_{k_r/k}(x)).
\end{equation}
Comparing the logarithm of \eqref{eq:Hasse-L} with
\eqref{eq:Hasse-log} gives
\[
 \sum_{x\in k_r}\phi_{k_r/k}(x)
 =(-1)^{r-1}\left(\sum_{x\in k}\phi(x)\right)^r.
\]
Multiplying by $-1$ is exactly \eqref{eq:Hasse-lift-sum}.
\end{proof}

This is \cite[Lemma 9.1, pp.~81--84]{Langlands}.




\begin{proposition}[Ramified Hasse-phase comparison]
\label{thm:ramified-Hasse-comparison}
Let $K/F$ be totally ramified cyclic of odd prime degree $\ell$, with
single break $t$.  Suppose
$m_F(\chi_F)=2d+1>1$, $2d+1\ge t+1$, and
$m_K(\chi_K)=2d'+1$, where
\[
 d'=\ell d-\frac{\ell-1}{2}t.
\]
Assume compatible representatives of the stationary classes in
Proposition~\ref{prop:high-parameter-table} have been chosen.  In either
of the following situations,
\[
 t=0\qquad\text{or}\qquad d\ge t+1,
\]
the residual factors satisfy
\begin{equation}\label{eq:ramified-Hasse-comparison}
 \Delta_{3,K}(\chi_K)=\Delta_{3,F}(\chi_F)^\ell.
\end{equation}
\end{proposition}

\begin{proof}
Put $n=n_F(\psi_F)$ and choose uniformizers with
$\varpi_F=N_{K/F}(\varpi_K)$.  The different exponent is
$(\ell-1)(t+1)$, and therefore
\[
 n_K(\psi_K)=\ell n+(\ell-1)(t+1),
 \qquad
 (2d'+1)+n_K(\psi_K)=\ell(2d+1+n).
\]
Consequently
\[
 \Gamma:=\varpi_F^{2d+1+n}\in F^\times
\]
is admissible both downstairs and upstairs.  By
Proposition~\ref{prop:high-parameter-table}(b), the downstairs and
upstairs stationary classes admit the same representative
$\beta\in\cO_F^\times$; choose this representative in both Lamprecht
factorizations.  If the conductor is odd, the corresponding critical
functions are understood with the simultaneous translation required by
Lemma~\ref{lem:complete-stationary-invariant}.

Let $k$ be the common residue field and, for $x\in k$, choose a
lift $\widetilde x\in\cO_F\subset\cO_K$.  Write $E_2(y)$ for the second
elementary symmetric function of the conjugates of $y$.  The lower
critical function is
\begin{equation}\label{eq:ramified-Hasse-lower-function}
 \phi(x)=
 \psi_F\!\left(\frac{\beta\varpi_F^d\widetilde x}{\Gamma}\right)
 \chi_F(1+\varpi_F^d\widetilde x)^{-1},
\end{equation}
and its polar character is
\[
 \psi_0(z)=
 \psi_F\!\left(\frac{\beta\varpi_F^{2d}\widetilde z}{\Gamma}\right).
\]
Put $Y=\varpi_K^{d'}\widetilde x$.  We claim that the upper critical
function is
\begin{equation}\label{eq:ramified-Hasse-upper-function}
 \phi_K(x)=
 \psi_F\!\left(\frac{\beta\Tr_{K/F}(Y)}{\Gamma}\right)
 \chi_F\!\left(1+\Tr_{K/F}(Y)+E_2(Y)\right)^{-1}.
\end{equation}
Indeed, it is enough to prove
\begin{equation}\label{eq:ramified-Hasse-norm-truncation-local}
 N_{K/F}(1+Y)
 \equiv1+\Tr_{K/F}(Y)+E_2(Y)
       \pmod{\p_F^{2d+1}}.
\end{equation}
If $t=0$, then $d'=\ell d$.  Every term of degree at least three in
the norm has valuation at least $3d\ge2d+1$, and the norm term has
valuation $\ell d\ge2d+1$.  If $t>0$, then $\ell=p$ is the residue
characteristic and $d\ge t+1$.  For $3\le j\le p-1$,
Lemma~\ref{lem:wild-symmetric-bound} gives
\[
 v_F(s_j(Y))\ge
 \left\lfloor\frac{j d'+(p-1)(t+1)}p\right\rfloor
 \ge2d+1,
\]
while
\[
 v_F(NY)=d'=pd-\frac{p-1}{2}t\ge2d+1.
\]
This proves \eqref{eq:ramified-Hasse-norm-truncation-local}.  The same
estimates, with the valuation of the lift increased by one, show that
\eqref{eq:ramified-Hasse-upper-function} is independent of the lift of
$x$.  They also give
\[
 \Tr(Y)\in\p_F^d,
 \qquad E_2(Y)\in\p_F^{2d},
\]
and one extra power of $\p_F$ when $\widetilde x\in\p_K$.

Set
\begin{equation}\label{eq:ramified-Hasse-epsilon}
 \epsilon=
 \frac{\Tr_{K/F}(\varpi_K^{2d'})}{\varpi_F^{2d}}\in\cO_F.
\end{equation}
The preceding valuation estimates show that this quotient is integral.
We now show that it is a unit and determine its residue.  If $t=0$,
then
\[
 c=\overline{\varpi_K^{d'}/\varpi_F^d}\in k^\times
\]
is fixed by $\Gal(K/F)$, because $d'=\ell d$.  Hence
\begin{equation}\label{eq:ramified-epsilon-tame}
 \bar\epsilon=\ell c^2.
\end{equation}
If $t>0$, Proposition~\ref{prop:stable-critical-trace} gives
\begin{equation}\label{eq:ramified-epsilon-wild}
 \bar\epsilon=-\lambda^{p-1}.
\end{equation}
In particular, $\epsilon$ is a unit in both cases.

For $x\in k$ write
\[
 Y_x=\varpi_K^{d'}\widetilde x,\qquad
 A_x=\Tr(Y_x),\qquad B_x=E_2(Y_x),\qquad
 U_x=1+A_x+B_x.
\]
For each fixed pair $x,y$, lift $x+y$ by
$\widetilde x+\widetilde y$; the lift-independence proved above permits
this choice, and then $Y_{x+y}=Y_x+Y_y$.  The identity
\[
 B_{x+y}=B_x+B_y+A_xA_y-\Tr(Y_xY_y)
\]
gives
\[
 U_xU_y
 \equiv U_{x+y}+\Tr(Y_xY_y)
 \pmod{\p_F^{2d+1}}.
\]
Indeed, the omitted products $A_xB_y$, $A_yB_x$, and $B_xB_y$ have
valuation at least $3d\ge2d+1$.  Since
$U_{x+y}\equiv1\pmod{\p_F^d}$ and
$\Tr(Y_xY_y)\in\p_F^{2d}$, division by $U_{x+y}$ gives
\begin{equation}\label{eq:ramified-Hasse-unit-quotient}
 \frac{U_xU_y}{U_{x+y}}
 \equiv1+\Tr(Y_xY_y)\pmod{\p_F^{2d+1}}.
\end{equation}
The additive factors in
\eqref{eq:ramified-Hasse-upper-function} cancel in the polar quotient,
whereas its multiplicative factors give
$\chi_F(U_xU_y/U_{x+y})$.  Lamprecht linearization on
$\p_F^{d+1}$ and \eqref{eq:ramified-Hasse-unit-quotient} therefore
imply
\begin{align*}
 \frac{\phi_K(x+y)}{\phi_K(x)\phi_K(y)}
 &=\psi_F\!\left(
    \frac{\beta\Tr(Y_xY_y)}{\Gamma}\right)\\
 &=\psi_0(\bar\epsilon xy),
\end{align*}
because
$\Tr(Y_xY_y)=\epsilon\varpi_F^{2d}\widetilde x\widetilde y$.
Thus
\begin{equation}\label{eq:ramified-Hasse-polar}
 \phi_K(x+y)=\phi_K(x)\phi_K(y)\psi_0(\bar\epsilon xy),
\end{equation}
and $\phi_K$ is a nondegenerate quadratic phase on $k$.

We compare its normalized sum with that of $\phi$.  First assume that
$k$ has odd characteristic.  The polar character of $\phi$ is
$(x,y)\mapsto\psi_0(xy)$, so for a unique $\alpha_0\in k$ one has
\[
 \phi(x)=\psi_0(x^2/2+\alpha_0x).
\]
Put
\[
 T_x=\overline{\Tr(Y)/\varpi_F^d},
 \qquad
 Q_x=\overline{\Tr(Y^2)/\varpi_F^{2d}}.
\]
Using $\Tr(Y)/\varpi_F^d$ itself as a lift of $T_x$ in
\eqref{eq:ramified-Hasse-lower-function}, one has
\[
 \phi(T_x)=
 \psi_F\!\left(\frac{\beta\Tr(Y)}{\Gamma}\right)
 \chi_F(1+\Tr(Y))^{-1}.
\]
Because $E_2(Y)\in\p_F^{2d}$,
\[
 \frac{1+\Tr(Y)}{1+\Tr(Y)+E_2(Y)}
 \equiv1-E_2(Y)\pmod{\p_F^{2d+1}}.
\]
A second application of Lamprecht linearization therefore gives the
exact residual identity
\begin{equation}\label{eq:ramified-Hasse-lower-upper-relation}
 \phi_K(x)=
 \phi(T_x)\,
 \psi_0\!\left(
  -\overline{\frac{E_2(Y)}{\varpi_F^{2d}}}\right).
\end{equation}
Now
$2E_2(Y)=(\Tr Y)^2-\Tr(Y^2)$, so
\[
 -\overline{\frac{E_2(Y)}{\varpi_F^{2d}}}
 =\frac{Q_x-T_x^2}{2}.
\]
Substituting
$\phi(T_x)=\psi_0(T_x^2/2+\alpha_0T_x)$ in
\eqref{eq:ramified-Hasse-lower-upper-relation} yields
\begin{equation}\label{eq:ramified-Hasse-quadratic-form}
 \phi_K(x)=\psi_0(Q_x/2+\alpha_0T_x).
\end{equation}
If $t=0$, then $T_x=\ell cx$ and $Q_x=\ell c^2x^2$; after replacing
$x$ by $c^{-1}x$,
\begin{equation}\label{eq:ramified-Hasse-tame-form}
 \phi_K(x)=\psi_0(\ell x^2/2+\ell\alpha_0x).
\end{equation}
If $t>0$, the trace-ideal formula gives $T_x=0$, while
$Q_x=\bar\epsilon x^2$, and hence
\begin{equation}\label{eq:ramified-Hasse-wild-form}
 \phi_K(x)=\psi_0(\bar\epsilon x^2/2).
\end{equation}
Let $\nu$ be the quadratic character of $k^\times$ and put
$Q=Q_{\psi_0}(1,0)$.  In the tame case, the phase of the lower sum is
$Q\psi_0(-\alpha_0^2/2)$ and that of the upper sum is
\[
 \nu(\ell)Q\psi_0(-\ell\alpha_0^2/2).
\]
Since $Q^2=\nu(-1)$, the required power identity is equivalent to
\begin{equation}\label{eq:ramified-Hasse-tame-sign}
 \nu(\ell)=\nu(-1)^{(\ell-1)/2}.
\end{equation}
Let $|k|=p^f$.  If $f$ is even, both sides are $1$.  If $f$ is odd,
the order of $p$ modulo $\ell$ is odd because $\ell\mid p^f-1$;
therefore $(p/\ell)=1$.  Quadratic reciprocity now gives exactly
\eqref{eq:ramified-Hasse-tame-sign}.

In the wild case, $\ell=p$ and
$\bar\epsilon=-\lambda^{p-1}$, so
$\nu(\bar\epsilon)=\nu(-1)$.  Formula
\eqref{eq:ramified-Hasse-wild-form} and $Q^2=\nu(-1)$ reduce the desired
identity to
\[
 \nu(-1)^{(p-3)/2}=1.
\]
If $p\equiv1\pmod4$, then $\nu(-1)=1$; if
$p\equiv3\pmod4$, the exponent $(p-3)/2$ is even.  This proves the
comparison in odd residue characteristic.

It remains to treat residue characteristic two.  Then $t=0$, because
$\ell$ is odd.  With the same $c$ as above, the norm calculation gives
\begin{equation}\label{eq:ramified-Hasse-char-two-power}
 \phi_K(c^{-1}x)
 =\phi(\ell x)\psi_0\!\left(-\binom{\ell}{2}x^2\right)
 =\phi(x)^\ell.
\end{equation}
The second equality follows by iterating
$\phi(x+y)=\phi(x)\phi(y)\psi_0(xy)$.  Put
$S=\sum_x\phi(x)$ and $H=\Aphase(S)$.  Choose $c_0\in k^\times$ so
that $\psi_0(x^2)=\psi_0(c_0x)$.  Then
\[
 \overline{\phi(x)}=\phi(x)\psi_0(c_0x),
 \qquad
 \overline S=\phi(c_0)^{-1}S,
 \qquad
 H^2=\phi(c_0).
\]
As $\phi(x)^4=1$, equation
\eqref{eq:ramified-Hasse-char-two-power} shows that the desired power
identity is equivalent to
\[
 \phi(c_0)^{(\ell-1)/2}=1\quad(\ell\equiv1\bmod4),
 \qquad
 \phi(c_0)^{(\ell+1)/2}=1\quad(\ell\equiv3\bmod4).
\]
Write $\psi_0(x)=(-1)^{\Tr_{k/\F_2}(a x)}$.  One may take
$c_0=a^{2^{f-1}-1}$, where $f=[k:\F_2]$; hence
\[
 \phi(c_0)^2=\psi_0(c_0^2)=(-1)^f.
\]
The two displayed powers are automatic for
$\ell\equiv1,7\pmod8$.  If $\ell\equiv3,5\pmod8$, then
$(2/\ell)=-1$; since $\ell\mid2^f-1$, the order of $2$ modulo
$\ell$ and therefore $f$ are even.  Thus $\phi(c_0)^2=1$, and the
remaining two residue classes also follow.

In every case the phase of the upper sum is the $\ell$th power of the
phase of the lower sum.  This is precisely
\eqref{eq:ramified-Hasse-comparison}.
\end{proof}

\begin{lemma}[Affine pencil of quadratic phases]\label{lem:affine-pencil}
Let $k$ have odd characteristic $p$, normalize $\psi$ so that
$\psi(x^p)=\psi(x)$, and take $\rho\neq0$ with
$1+j\rho\neq0$ for every $j\in\F_p$.  Choose $a_0,b_0\in k$ with
\[
 a_0^p=1-\rho^{1-p},\qquad
 b_0^p=\sigma-\rho^{-1}\tau.
\]
Then
\begin{equation}\label{eq:affine-pencil}
 Q_\psi(a_0,b_0)
 \prod_{j\in\F_p^\times}Q_\psi(j\rho,j\tau)
 =\prod_{j\in\F_p}Q_\psi(1+j\rho,\sigma+j\tau).
\end{equation}
\end{lemma}

\begin{proof}
Apply Lemma~\ref{lem:quadratic-phase-new}.  The quadratic-character
parts agree because
\[
 \prod_{j\in\F_p}(1+j\rho)=1-\rho^{p-1},
 \qquad
 \prod_{j\neq0}j=-1,
\]
and $\nu(a_0)=\nu(1-\rho^{1-p})$.  The additive parts agree by
\begin{equation}\label{eq:rational-sum}
 \sum_{j\in\F_p}\frac{(\sigma+j\tau)^2}{1+j\rho}
 =\frac{\rho^{p-3}(\sigma\rho-\tau)^2}{\rho^{p-1}-1}.
\end{equation}
To prove \eqref{eq:rational-sum}, use
\[
 \sum_j\frac1{1+j\rho}=\frac{\rho^{p-1}}{\rho^{p-1}-1},\quad
 \sum_j\frac{j}{1+j\rho}=-\frac{\rho^{p-2}}{\rho^{p-1}-1},\quad
 \sum_j\frac{j^2}{1+j\rho}=\frac{\rho^{p-3}}{\rho^{p-1}-1},
\]
which follow from the logarithmic derivative of
$\prod_j(X+j)=X^p-X$.  Finally the right side of
\eqref{eq:rational-sum} is $(b_0^2/a_0)^p$, and the normalization
$\psi(z^p)=\psi(z)$ finishes the proof.
\end{proof}


\begin{proposition}[Homogeneous affine-pencil quotient]
\label{prop:affine-pencil-homogeneous}
Under the hypotheses and notation of
Lemma~\ref{lem:affine-pencil}, let $c\in k^\times$ and choose
$c_0\in k^\times$ with $c_0^p=c$.  Then
\begin{align}
&\frac{
 Q_\psi(ca_0^p,cb_0^p)
 \displaystyle\prod_{j\in\F_p^\times}Q_\psi(cj\rho,cj\tau)}
 {\displaystyle\prod_{j\in\F_p}
 Q_\psi(c(1+j\rho),c(\sigma+j\tau))}
\notag\\[1mm]
&\hspace{18mm}=
\frac{
 Q_\psi(a_0,b_0)
 \displaystyle\prod_{j\in\F_p^\times}Q_\psi(j\rho,j\tau)}
 {\displaystyle\prod_{j\in\F_p}
 Q_\psi(1+j\rho,\sigma+j\tau)}.
\label{eq:affine-pencil-homogeneous}
\end{align}
Thus a common nonzero scale on all actual rows does not change the
normalized affine-pencil quotient.
\end{proposition}

\begin{proof}
For $a,c\ne0$, completion of the square gives the scaling identity
\begin{equation}\label{eq:quadratic-phase-common-scale}
 Q_\psi(ca,cb)
 =\nu(c)\,
  \psi\!\left(-\frac{(c-1)b^2}{2a}\right)Q_\psi(a,b),
\end{equation}
where $\nu$ is the quadratic character.  Frobenius invariance also gives
$Q_\psi(a^p,b^p)=Q_\psi(a,b)$.  Apply
\eqref{eq:quadratic-phase-common-scale} to every row.  The product of
the quadratic-character factors cancels because numerator and
denominator contain the same number of rows and $p-1$ is even.  The
unit-indexed additive corrections have total exponent zero because
$\sum_{j\in\F_p^\times}j=0$.  For the denominator, the rational-sum
identity \eqref{eq:rational-sum} shows that its total correction is the
Frobenius of the upper correction.  The relation $c_0^p=c$ and
$\psi(z^p)=\psi(z)$ identify these two corrections.  The remaining
common factor therefore cancels from the quotient, proving
\eqref{eq:affine-pencil-homogeneous}.
\end{proof}

\begin{lemma}[Artin--Schreier phase]\label{lem:AS-phase-new}
Let $k$ have characteristic $2$ and let
$\psi_0(x)=(-1)^{\Tr_{k/\F_2}(x)}$.  Then
\[
 \psi_0(c+c^2)=1\qquad(c\in k).
\]
\end{lemma}

\begin{proof}
The absolute trace is Frobenius-invariant, so
$\Tr(c+c^2)=0$.
\end{proof}

\section{Stationary parameters under norm}\label{sec:stationary-parameters}

We now establish the stationary-parameter comparison under norm.  The
formulation is specialized to cyclic prime degree.

Throughout this section a symbol such as $\beta(\theta)$ denotes the
stationary \emph{class} furnished by
Lemma~\ref{lem:stationary-class-existence}, unless a representative is
explicitly chosen.  Thus a formula
\[
 \beta(\theta)\equiv c\pmod{\p_E^r}
\]
means that $c$ represents that class at the indicated precision; it
does not assert that the character canonically determines the element
$c\in E$.  When we say that one may take $\beta(\theta)=c$, we are
making a deliberate choice of representative.  Products and powers of
characters are handled by Lemma~\ref{lem:stationary-class-calculus}; if
a conductor drops but remains greater than one, the stationary class is
redefined at the new conductor and only its restriction to the old
layer is represented by the old sum.  If it drops to conductor zero or
one, the corresponding endpoint formula is used and there is no
stationary class.  Finally, changing a representative is always accompanied
by the translation of the odd critical function prescribed by
Lemma~\ref{lem:complete-stationary-invariant}.  Hence every complete
Lamprecht factor used below is independent of these choices.

Let
$m=m_F(\chi_F)=2d+\varepsilon>1$ and
$m'=m_K(\chi_F\circ N)=2d'+\varepsilon'$.  We use one of the following depth
conventions.
\begin{enumerate}[label=(\roman*)]
\item If $K/F$ is unramified, then $m'=m$.
\item If $K/F$ is ramified with break $t$ and $m>t+1$, then
\[
 m'=\HerbPsi_{K/F}(m-1)+1.
\]
\item If $K/F$ is ramified, $2\le m\le t+1$, and $\chi_F$ has minimal
      conductor in its $S(K/F)$-orbit, then $m'=m$;
      this is \eqref{eq:minimal-critical}.
\end{enumerate}
For $x\in\p_K^{d'+\varepsilon'}$ put
\begin{equation}\label{eq:P-map}
 P_{K/F}(x)=N(1+x)-1\pmod{\p_F^m}.
\end{equation}

For every finite separable extension, multiplication of the conjugates
of $1+x$ gives the exact identity
\begin{equation}\label{eq:norm-expansion}
 N_{K/F}(1+x)
 =1+\sum_{j=1}^{[K:F]}s_j(x),
 \qquad s_1(x)=\Tr_{K/F}(x),\quad
 s_{[K:F]}(x)=N_{K/F}(x),
\end{equation}
where $s_j(x)$ is the $j$th elementary symmetric function of the
$F$-conjugates of $x$.  This identity is polynomial and remains valid
regardless of the residue characteristic because $K/F$ itself is
separable.

\begin{lemma}[Norm filtration at stationary depth]
\label{lem:stationary-depth-norm}
In each of the three depth conventions above,
\begin{align}
 N(U_K^{d'+\varepsilon'})&\subseteq U_F^{d+\varepsilon},
                                                        \label{eq:stationary-depth-first-inclusion}\\
 N(U_K^{m'})&\subseteq U_F^m,             \label{eq:stationary-depth-second-inclusion}\\
 N(U_K^{d'})&\subseteq U_F^d.             \label{eq:stationary-parameter-ambiguity}
\end{align}
In the unramified case these inclusions are equalities.  Consequently,
if $b,b'\in\cO_K^\times$ and
$b'\equiv b\pmod{\p_K^{d'}}$, then
\begin{equation}\label{eq:norm-of-stationary-class-well-defined}
 N(b')\equiv N(b)\pmod{\p_F^d}.
\end{equation}
Thus the norm of a stationary numerator class modulo $\p_K^{d'}$ is a
well-defined class modulo $\p_F^d$.
\end{lemma}

\begin{proof}
Put
\[
 q=\left\lceil\frac m2\right\rceil=d+\varepsilon,
 \qquad
 q'=\left\lceil\frac {m'}2\right\rceil=d'+\varepsilon'.
\]
In the unramified case Lemma~\ref{lem:unramified-compatibility}(c)
gives equality at every unit depth.

Assume that $K/F$ is ramified.  We first prove
\eqref{eq:stationary-depth-first-inclusion} in the high range.  Set
$r=m-1\ge t$ and $s=\lfloor r/2\rfloor=q-1$.  Then
\[
 m'=\HerbPsi_{K/F}(r)+1,
 \qquad
 q'-1=\left\lfloor\frac{\HerbPsi_{K/F}(r)}2\right\rfloor.
\]
We claim that
\begin{equation}\label{eq:half-Herbrand-inequality}
 \left\lfloor\frac{\HerbPsi_{K/F}(r)}2\right\rfloor
 \ge \HerbPsi_{K/F}(s).
\end{equation}
If $s\le t$, then the right side is $s$, whereas
$\HerbPsi_{K/F}(r)\ge r\ge2s$.  If $s>t$, write
$r=2s+e$ with $e\in\{0,1\}$.  The explicit formula
\eqref{eq:Herbrand-prime-new} gives
\[
 \HerbPsi_{K/F}(r)-2\HerbPsi_{K/F}(s)
   =(\ell-1)t+\ell e\ge0,
\]
which again proves \eqref{eq:half-Herbrand-inequality}.  Hence
\begin{equation}\label{eq:stationary-half-depth-bound}
 q'\ge\HerbPsi_{K/F}(q-1)+1.
\end{equation}

If $q-1<t$, norm preserves the filtration below the break, so
$N(U_K^q)\subseteq U_F^q$.  If $q-1=t$, the critical norm map gives
$N(U_K^{t+1})\subseteq U_F^{t+1}$.  If $q-1>t$, part~(iv) of
Theorem~\ref{thm:graded-norm} gives
$N(U_K^{\HerbPsi_{K/F}(q-1)+1})=U_F^q$.  Together with
\eqref{eq:stationary-half-depth-bound} and the decreasing unit
filtration, these three cases prove
\eqref{eq:stationary-depth-first-inclusion}.

For \eqref{eq:stationary-depth-second-inclusion}, if $m>t+1$,
part~(iv) of the graded norm theorem, applied with $r=m-1$, gives
$N(U_K^{m'})=U_F^m$.  If $m=t+1$, then $m'=t+1$ and the critical norm
map gives $N(U_K^{m'})\subseteq U_F^m$.

We next prove \eqref{eq:stationary-parameter-ambiguity} in the high
range.  Since $\HerbPsi_{K/F}(u)\ge u$, one has $m'\ge m$ and hence
$d'\ge d$.  If $d\le t$, norm preserves the filtration at depth $d$,
so
\[
 N(U_K^{d'})\subseteq N(U_K^d)\subseteq U_F^d.
\]
If $d=t+1$, the same conclusion follows from the critical inclusion
$N(U_K^{t+1})\subseteq U_F^{t+1}$.  Finally suppose $d\ge t+2$ and put
\[
 R=\HerbPsi_{K/F}(d-1)+1.
\]
Using \eqref{eq:Herbrand-prime-new} and $m=2d+\varepsilon$, one obtains
\begin{align*}
 m'-2R
 &=\HerbPsi_{K/F}(2d+\varepsilon-1)+1
     -2\bigl(\HerbPsi_{K/F}(d-1)+1\bigr)\\
 &=(\ell-1)(t+1)+\ell\varepsilon\ge0.
\end{align*}
Thus $d'=\lfloor m'/2\rfloor\ge R$.  Part~(iv) of
Theorem~\ref{thm:graded-norm}, applied to $d-1>t$, now gives
\[
 N(U_K^{d'})\subseteq N(U_K^R)=U_F^d.
\]
This proves the third inclusion in the high range.

In the low range $m'=m\le t+1$, so $q'=q$ and $d'=d\le t$.
Below the break, norm maps $U_K^i$ into $U_F^i$ for every $i\le t$;
at the last possible step $m=t+1$ one uses the critical inclusion
$N(U_K^{t+1})\subseteq U_F^{t+1}$.  Applying these statements at
$i=q$, $i=m$, and $i=d$ proves all three inclusions.

For the final assertion, write $b'=b(1+z)$ with
$z\in\p_K^{d'}$.  Then $1+z\in U_K^{d'}$, so
\eqref{eq:stationary-parameter-ambiguity} gives
$N(b')/N(b)\in U_F^d$, which is equivalent to
\eqref{eq:norm-of-stationary-class-well-defined}.
\end{proof}

\begin{lemma}[Truncated norm homomorphism and adjoint]
\label{lem:truncated-norm-adjoint}
The formula \eqref{eq:P-map} induces an additive homomorphism
\[
 P_{K/F}:\p_K^{d'+\varepsilon'}/\p_K^{m'}
 \longrightarrow
 \p_F^{d+\varepsilon}/\p_F^m.
\]
For each admissible pair $(\gamma_F,\gamma_K)$, the pairings
\[
 (\bar x,\bar y)\longmapsto
 \psi_F(xy/\gamma_F),\qquad
 (\bar x',\bar y')\longmapsto
 \psi_K(x'y'/\gamma_K)
\]
between
$\cO_F/\p_F^d$ and
$\p_F^{d+\varepsilon}/\p_F^m$, and between
$\cO_K/\p_K^{d'}$ and
$\p_K^{d'+\varepsilon'}/\p_K^{m'}$, are perfect.  Consequently there
is a unique additive adjoint
\begin{equation}\label{eq:Pstar}
 P^*_{K/F;\gamma_F,\gamma_K}:
 \cO_F/\p_F^d\longrightarrow\cO_K/\p_K^{d'}
\end{equation}
satisfying
\begin{equation}\label{eq:Pstar-def}
 \psi_F\left(\frac{xP_{K/F}(y)}{\gamma_F}\right)
 =\psi_K\left(
   \frac{P^*_{K/F;\gamma_F,\gamma_K}(x)y}{\gamma_K}\right).
\end{equation}
When the denominator pair is fixed, we abbreviate this map to
$P^*_{K/F}$.  Thus the adjoint notation always includes the chosen
pair, even when it is suppressed.  If $P_{K/F}=\Tr_{K/F}$ at the
required precision and $\gamma_K/\gamma_F\in\cO_K^\times$, then
\begin{equation}\label{eq:Pstar-trace-denominator-ratio}
 P^*_{K/F;\gamma_F,\gamma_K}(x)
 \equiv\frac{\gamma_K}{\gamma_F}x\pmod{\p_K^{d'}}.
\end{equation}
In particular, when
$\gamma_F=\gamma_K=\gamma\in F^\times$, the adjoint is the natural
inclusion $\cO_F/\p_F^d\to\cO_K/\p_K^{d'}$.
\end{lemma}

\begin{proof}
The first inclusion in Lemma~\ref{lem:stationary-depth-norm} places
$P_{K/F}(x)$ in the stated target.  If $x'$ differs from $x$ by an
element of $\p_K^{m'}$, then
\[
 \frac{1+x'}{1+x}\in U_K^{m'},
\]
so the second inclusion shows that their norms are congruent modulo
$U_F^m$; hence $P_{K/F}$ is well defined.

Let $x,y\in\p_K^{d'+\varepsilon'}$.  Since
$2(d'+\varepsilon')\ge m'$, one has $xy\in\p_K^{m'}$.  Multiplicativity
of norm gives
\[
 1+P(x+y+xy)=(1+P(x))(1+P(y)).
\]
The preceding well-definedness permits deletion of $xy$.  Moreover
$P(x),P(y)\in\p_F^{d+\varepsilon}$ and
$2(d+\varepsilon)\ge m$, so $P(x)P(y)\in\p_F^m$.  Reducing the displayed
identity modulo $\p_F^m$ gives $P(x+y)=P(x)+P(y)$.

Perfectness of the two pairings is Lemma~\ref{lem:finite-duality}, once
over $F$ and once over $K$.  For fixed $x\in\cO_F/\p_F^d$, the left
side of \eqref{eq:Pstar-def} is a character of the source of $P$.
Perfectness over $K$ supplies a unique representing class in
$\cO_K/\p_K^{d'}$.  Character multiplication and uniqueness show that
this representing class depends additively on $x$.  This is the desired
adjoint.

Finally suppose that $P_{K/F}=\Tr_{K/F}$ at the required precision and
that $\gamma_K/\gamma_F\in\cO_K^\times$.  Trace transitivity gives
\[
 \psi_F\!\left(\frac{x\Tr_{K/F}(y)}{\gamma_F}\right)
 =\psi_K\!\left(\frac{xy}{\gamma_F}\right)
 =\psi_K\!\left(
   \frac{(\gamma_K/\gamma_F)xy}{\gamma_K}\right).
\]
Perfectness over $K$ proves
\eqref{eq:Pstar-trace-denominator-ratio}; the common-denominator case is
its specialization.
\end{proof}

\begin{corollary}[Stationary class under norm pullback]
\label{cor:stationary-class-under-norm}
Let
\[
 \beta_F=\operatorname{St}_{\gamma_F,d+\varepsilon}(\chi_F)
       \in\cO_F/\p_F^d,
 \qquad
 \beta_K=\operatorname{St}_{\gamma_K,d'+\varepsilon'}(\chi_F\circ N)
       \in\cO_K/\p_K^{d'}.
\]
Then, with
$P^*_{K/F}=P^*_{K/F;\gamma_F,\gamma_K}$,
\begin{equation}\label{eq:stationary-class-adjoint}
 \beta_K=P^*_{K/F}(\beta_F).
\end{equation}
In particular, every explicit formula for $P^*_{K/F}(\beta_F)$ below
is an equality of stationary classes.  Any element representing its
right-hand side may be chosen as the upstairs stationary numerator;
its complete Lamprecht factor is independent of that choice by
Lemma~\ref{lem:complete-stationary-invariant}.
\end{corollary}

\begin{proof}
For $y\in\p_K^{d'+\varepsilon'}$, the definition of $P_{K/F}$ and
Lemma~\ref{lem:stationary-depth-norm} give
\begin{align*}
 (\chi_F\circ N)(1+y)
 &=\chi_F(1+P_{K/F}(y))\\
 &=\psi_F\!\left(\frac{\beta_FP_{K/F}(y)}{\gamma_F}\right)\\
 &=\psi_K\!\left(\frac{P^*_{K/F}(\beta_F)y}{\gamma_K}\right).
\end{align*}
The uniqueness assertion in
Lemma~\ref{lem:stationary-class-existence}, applied over $K$, identifies
this class with $\beta_K$.
\end{proof}

In the unramified case, $d'=d$ and
$\varepsilon'=\varepsilon$.  If
$y\in\p_K^{d+\varepsilon}$, the $j$th elementary symmetric term in
$N(1+y)$ has valuation at least $j(d+\varepsilon)$; hence every term
with $j\geq2$ lies in $\p_F^m$, and
\begin{equation}\label{eq:unramified-P-trace}
 P_{K/F}(y)\equiv\Tr_{K/F}(y)\pmod{\p_F^m}.
\end{equation}

\begin{lemma}[Integer powers and Teichmuller representatives]
\label{lem:teichmueller-stationary}
Assume that $K/F$ is wildly ramified of degree $p$ and put
$T=t+1$ and $s_0=\lceil T/2\rceil$.  For $j\in\F_p$, let
$j_{\mathbf Z}$ be its integer representative in
$\{0,\ldots,p-1\}$ and let $[j]\in\cO_F$ be its Teichmuller lift.  If
\[
 \tau(1+x)=\psi_F(Ax)\qquad(x\in\p_F^{s_0}),
\]
then
\begin{equation}\label{eq:teichmueller-stationary}
 \psi_F(j_{\mathbf Z}Ax)=\psi_F([j]Ax)
 \qquad(x\in\p_F^{s_0}).
\end{equation}
Thus every stationary numerator for the integer power $\tau^j$ may be
represented by the Teichmuller scalar $[j]$ at all depths occurring
below.  The Teichmuller lifts satisfy
\begin{equation}\label{eq:teichmueller-polynomials}
 \prod_{j\in\F_p}(X-[j])=X^p-X,
 \qquad
 \prod_{j\in\F_p}(X+[j])=X^p+(-1)^pX,
\end{equation}
and
\begin{equation}\label{eq:teichmueller-sum}
 \sum_{j\in\F_p}[j]=
 \begin{cases}0,&p\text{ odd},\\1,&p=2.\end{cases}
\end{equation}
\end{lemma}

\begin{proof}
In equal characteristic, the prime-field element $j_{\mathbf Z}$ is
already its Teichmuller lift.  Suppose that $F$ has characteristic zero
and put $e=v_F(p)$.  Since $p=\Tr_{K/F}(1)$, the trace-ideal formula
\eqref{eq:trace-ideal} with $a=0$ gives
\[
 e\ge \left\lfloor\frac{(p-1)T}{p}\right\rfloor
   \ge \left\lfloor\frac T2\right\rfloor.
\]
Moreover $[j]-j_{\mathbf Z}\in p\cO_F$.  Hence for
$x\in\p_F^{s_0}$,
\[
 ([j]-j_{\mathbf Z})x\in
 \p_F^{e+s_0}\subseteq\p_F^T.
\]
The additive character $x\mapsto\psi_F(Ax)$ represents the last
nontrivial unit layer of the conductor-$T$ character $\tau$, and is
therefore trivial on $\p_F^T$.  This proves
\eqref{eq:teichmueller-stationary}.

The elements $[j]$ are precisely the $p$ roots of $X^p-X$ in the
Teichmuller copy of the prime field.  This proves the first polynomial
identity; replacing $X$ by $-X$ gives the second.  The coefficient of
$X^{p-1}$ gives the sum formula.
\end{proof}

\begin{remark}[Notation for powers of the norm character]
From this point onward, an exponent $j$ in $\tau^j$ denotes the integer
representative of a class in $\F_p$, while the same symbol inside a
field-valued formula denotes its Teichmuller lift.  The preceding lemma
proves that these two uses give the same stationary character at every
precision used below.
\end{remark}

\begin{lemma}[Norm representatives at subcritical precision]
\label{lem:norm-stationary-representatives}
Let $K/F$ be ramified cyclic of prime degree with break $t$.  Let
$h\in\mathbf Z$, let $0\le s\le t$, and let $c\in F^\times$ have
$v_F(c)=h$.  There is $c_1\in K^\times$ with $v_K(c_1)=h$ such that
\begin{equation}\label{eq:norm-representative-precision}
 c-N(c_1)\in\p_F^{h+s}.
\end{equation}
Consequently every stationary coefficient class of the form
$c+\p_F^{h+s}$, with $s\le t$, has a representative which is an exact
norm.  Any finite collection of such classes admits norm
representatives chosen independently and hence simultaneously.
\end{lemma}

\begin{proof}
The norm isomorphism
\[
 K^\times/U_K^t\xrightarrow{\sim}F^\times/U_F^t
\]
of Theorem~\ref{thm:norm-filtration}(c) supplies $c_1$ with
$c/N(c_1)\in U_F^t$.  Total ramification gives
$v_K(c_1)=v_F(N(c_1))=h$.  Since $c$ and $N(c_1)$ have the same
valuation and their quotient lies in $U_F^t$, their difference belongs
to $\p_F^{h+t}\subseteq\p_F^{h+s}$.  Applying this argument separately
to each member of a finite collection proves the last assertion.
\end{proof}

\subsection{High-conductor parameters}

We first give the wild non-stable calculation in odd prime degree.
Let $K/F$ be totally and wildly ramified of odd prime degree $p$, put
\[
 T=t+1,\qquad D=(p-1)T,
\]
and suppose
\[
 T\le m=m_F(\chi_F)=2d+\varepsilon\le2T-1,
 \qquad m=T+v\quad(0\le v\le t).
\]
As throughout the ramified analysis, assume that $\chi_F$ has minimal
conductor in its $S(K/F)$-orbit.  Then
\begin{equation}\label{eq:mprime-high}
 m'=m_K(\chi_K)=T+pv=pm-D=2d'+\varepsilon.
\end{equation}
For $v>0$ this follows from
Proposition~\ref{prop:pullback-conductor}(i), while for $v=0$ it is
\eqref{eq:minimal-critical}.  In particular, $m$ and $m'$ have the same
parity.  Put
\[
 r_0=\left\lfloor\frac T2\right\rfloor,
 \qquad s_0=\left\lceil\frac T2\right\rceil.
\]
Relative to the denominator $\gamma_F$, the stationary class of a
generator $\tau$ of $S(K/F)$ on $\p_F^{s_0}$ is
\begin{equation}\label{eq:tau-high-stationary-class}
 \operatorname{St}_{\gamma_F,s_0}(\tau)
 \in\p_F^v/\p_F^{v+r_0},
\end{equation}
because $m=T+v$ and $r_0+s_0=T$.  Since $r_0\le t$,
Lemma~\ref{lem:norm-stationary-representatives} permits the choice of a
representative
\begin{equation}\label{eq:alpha-high}
 \alpha=N(\alpha_1),\qquad v_K(\alpha_1)=v,
\end{equation}
for which
\begin{equation}\label{eq:tau-high-linearization}
 \tau(1+x)=\psi_F(\alpha x/\gamma_F)
 \qquad(x\in\p_F^{s_0}).
\end{equation}
Likewise
\[
 \operatorname{St}_{\gamma_F,d+\varepsilon}(\chi_F)
 \in\cO_F/\p_F^d,
\]
and $d\le t$.  The same lemma therefore permits an exact norm
representative
\begin{equation}\label{eq:beta-high-norm}
 \beta=N(\beta_1),\qquad \beta_1\in U_K,
\end{equation}
satisfying
\begin{equation}\label{eq:chi-high-linearization}
 \chi_F(1+x)=\psi_F(\beta x/\gamma_F)
 \qquad(x\in\p_F^{d+\varepsilon}).
\end{equation}
The two norm representatives are chosen independently and hence
simultaneously.  If $m=T$, then
Lemma~\ref{lem:minimal-orbit-stationary} shows that none of the leading
classes $\beta+j\alpha$ for $j\in\F_p$ cancels; thus every twist in the
table below still has exact conductor $T$.

\begin{lemma}[Norm-character conversion]\label{lem:norm-character-conversion}
For $y\in\p_K^{s_0}$,
\begin{equation}\label{eq:conversion}
 \psi_F\!\left(\frac{\alpha N(y)}{\gamma_F}\right)
 =\psi_F\!\left(-\frac{\alpha\Tr(y)}{\gamma_F}\right).
\end{equation}
\end{lemma}

\begin{proof}
Corollary~\ref{cor:wild-norm-truncation}, with $q=s_0$ and $r=T$,
gives
\[
 N(1+y)\equiv1+\Tr(y)+N(y)\pmod{\p_F^T},
\]
because $2s_0\ge T$.  The trace-ideal formula and total ramification
give
\[
 v_F(\Tr(y))\ge
 \left\lfloor\frac{s_0+D}{p}\right\rfloor\ge s_0,
 \qquad v_F(N(y))=v_K(y)\ge s_0.
\]
Thus the sum lies in the layer on which
\eqref{eq:tau-high-linearization} holds.  Since $\tau$ is trivial on
norms and has conductor $T$, applying $\tau$ to $N(1+y)$ proves
\eqref{eq:conversion}.
\end{proof}

\begin{proposition}[High-conductor stationary-class table]
\label{prop:high-parameter-table}
In each case choose an admissible denominator $\gamma_F$ downstairs and
put $\gamma_K=\gamma_F$, with the same element of $F^\times$ viewed in
$K^\times$.  The conductor identities in the proof show that this is
admissible upstairs.  Throughout the statement,
$P^*_{K/F}$ denotes
$P^*_{K/F;\gamma_F,\gamma_K}$ for this common-denominator pair.
The following statements hold.
\begin{enumerate}[label=(\alph*)]
\item If $K/F$ is unramified, then
      $P^*_{K/F}(\beta)=\beta$ in $\cO_K/\p_K^{d'}$.
\item If $K/F$ is totally ramified of odd prime degree and
      $d\ge t+1$, then
      $P^*_{K/F}(\beta)=\beta$ in $\cO_K/\p_K^{d'}$.
\item In the odd wild non-stable range above, the upstairs stationary
      class is
\begin{equation}\label{eq:Pstar-intermediate}
 \beta_K:=\beta(\chi_K)
 \equiv \beta-\frac{\beta_1\alpha}{\alpha_1}
 \pmod{\p_K^{d'}}.
\end{equation}
For $j\in\F_p$, represented by its Teichmuller lift, the stationary
class of the twist is
\begin{equation}\label{eq:twist-beta-intermediate}
 \beta_j:=\beta(\tau^j\chi_F)
 \equiv N(\beta_1+j\alpha_1)
 \equiv\beta+j\alpha\pmod{\p_F^d}.
\end{equation}
Every right-hand side in these congruences may be selected as the
corresponding stationary representative.  Moreover the norm of the
upstairs class is well defined modulo $\p_F^d$ and satisfies
\begin{equation}\label{eq:high-product-congruence}
 N(\beta_K)\equiv
 \prod_{j\in\F_p}(\beta+j\alpha)
 =\beta^p-\beta\alpha^{p-1}\pmod{\p_F^d}.
\end{equation}
\item In the tamely ramified quadratic case,
      $P^*_{K/F}(\beta)=\beta$ in $\cO_K/\p_K^{d'}$.
\end{enumerate}
All assertions in (a), (b), and (d) are equalities of stationary
classes as in Corollary~\ref{cor:stationary-class-under-norm}; no
canonical field representative is asserted.
\end{proposition}

\begin{proof}
We first verify the denominator compatibility used in the statement.
In the unramified case,
$m'=m$ and $n_K(\psi_K)=n_F(\psi_F)$, so an admissible
$\gamma_F$ remains admissible in $K$.  In every ramified case in this
high-conductor table, of degree $\ell$, one has
\[
 m'=\ell m-D,
 \qquad n_K(\psi_K)=\ell n_F(\psi_F)+D,
\]
and therefore
\[
 m'+n_K(\psi_K)=\ell\bigl(m+n_F(\psi_F)\bigr)
 =v_K(\gamma_F).
\]
Thus $\gamma_K=\gamma_F$ is admissible upstairs in all four parts.

In the unramified case, equation~\eqref{eq:unramified-P-trace}, trace
transitivity, and the common denominator identify the adjoint with the
inclusion $F\hookrightarrow K$.
Corollary~\ref{cor:stationary-class-under-norm} then proves (a).

For (b), put $q=d'+\varepsilon'=\lceil m'/2\rceil$ and separate
the tame and wild cases.  If $K/F$ is tamely ramified of odd degree
$\ell$, then
\[
 m'-1=\ell(m-1).
\]
For $x\in\p_K^q$ and $2\le j\le\ell-1$, every summand defining
$s_j(x)$ has $K$-valuation at least $jq$.  Since $s_j(x)\in F$ and
$v_K|_F=\ell v_F$, it follows that
\[
 v_F(s_j(x))\ge\left\lceil\frac{jq}{\ell}\right\rceil
 \ge\left\lceil\frac{2q}{\ell}\right\rceil\ge m.
\]
Also $v_F(N(x))=v_K(x)\ge q\ge m$.  Thus all nonlinear norm terms
vanish modulo $\p_F^m$.

If $K/F$ is wildly ramified, its degree is the residue characteristic
$p$, $D=(p-1)T$, and
\[
 m'=p(m-T)+T.
\]
The assumption $d\ge T$ gives $q\ge m$.  Moreover
$2q+D\ge m'+D=pm$, so the strong symmetric estimate gives
$v_F(s_j(x))\ge m$ for $2\le j\le p-1$.  Again the norm term has
valuation at least $m$.  Hence in either case $P=\Tr$ at the required
precision and $P^*(\beta)=\beta$.  The class interpretation again
follows from Corollary~\ref{cor:stationary-class-under-norm}.

For (c), put $q=\lceil m'/2\rceil$.  Since
$2q+D\ge pm$, Corollary~\ref{cor:wild-norm-truncation} gives
\begin{equation}\label{eq:high-upper-trunc}
 N(1+x)-1\equiv\Tr(x)+N(x)\pmod{\p_F^m}
 \qquad(x\in\p_K^q).
\end{equation}
The elementary identity
$q-v\ge d+\varepsilon\ge s_0$ shows that
$y=\beta_1x/\alpha_1$ lies in the domain of
Lemma~\ref{lem:norm-character-conversion}.  Because
$\beta=N(\beta_1)$ exactly, the conversion identity gives
\[
 \psi_F\!\left(\frac{\beta N(x)}{\gamma_F}\right)
 =\psi_F\!\left(-\frac{\alpha\Tr(\beta_1x/\alpha_1)}
                         {\gamma_F}\right).
\]
Combining this with the trace term in
\eqref{eq:high-upper-trunc} proves that the class on the right of
\eqref{eq:Pstar-intermediate} is $P^*_{K/F}(\beta)$.  The element
$\gamma_F$, viewed in $K$, is admissible upstairs because
\[
 m'+n_K(\psi_K)=p\bigl(m+n_F(\psi_F)\bigr).
\]
Corollary~\ref{cor:stationary-class-under-norm} now proves
\eqref{eq:Pstar-intermediate} as an equality of stationary classes.

On the common stationary layer, Lemma~\ref{lem:stationary-class-calculus}(c)
adds the classes of $\chi_F$ and $\tau^j$, giving the class of
$\beta+j\alpha$.  Its conductor is exactly $m$ by
Lemma~\ref{lem:minimal-orbit-stationary}; in particular, the boundary
case $m=T$ involves no unrecorded conductor drop.  Factoring
$N(\beta_1+j\alpha_1)=\beta N(1+j\alpha_1/\beta_1)$, the strong
symmetric estimate gives, for every intermediate term, valuation at
least $d$, since
\[
 iv+(p-1)T\ge pd\qquad(1\le i\le p-1).
\]
The two endpoint terms are $\beta$ and $j\alpha$, because $j^p=j$ for
the chosen Teichmuller lift.  This proves
\eqref{eq:twist-beta-intermediate} as a class congruence.

For the product congruence, choose the normalized representative
\[
 \beta_K^0=\beta-\frac{\beta_1\alpha}{\alpha_1}
\]
of the class in \eqref{eq:Pstar-intermediate}.  By
\eqref{eq:norm-of-stationary-class-well-defined}, the class of
$N(\beta_K)$ modulo $\p_F^d$ is the class of $N(\beta_K^0)$, so this
choice entails no loss of information.  Put
$y=\beta_1\alpha/\alpha_1$.  Then
$v_K(y)=(p-1)v$ and $N(y)=\beta\alpha^{p-1}$.  In the expansion of
$N(\beta-y)$ every intermediate term has valuation at least $d$,
because
\[
 i(p-1)v+(p-1)T\ge pd\qquad(1\le i\le p-1).
\]
The endpoints give $\beta^p-\beta\alpha^{p-1}$, which equals
$\prod_{j\in\F_p}(\beta+j\alpha)$.  This proves
\eqref{eq:high-product-congruence} independently of the representative
of $\beta_K$.

Finally, in the tame quadratic case the stationary source is
$\p_K^m$ and the quadratic norm term is already in $\p_F^m$; hence
$P=\Tr$ and $P^*(\beta)=\beta$.  Corollary~\ref{cor:stationary-class-under-norm}
again supplies the class interpretation.
\end{proof}

\begin{proposition}[Compatible wild-quadratic stationary classes and representatives]
\label{prop:wild-quadratic-high-parameters}
Assume $p=2$, $K/F$ is wildly ramified, and the character is in the
high-conductor range.  Representatives of the relevant stationary
classes may be chosen in the form
\[
 \alpha=N(\alpha_1),\qquad
 \beta=\delta_hN(\beta_1),\qquad \delta_h\in U_F^t,
\]
with
$\tau(1+x)=\psi_F(\alpha\delta_hx/\gamma_F)$ on the required unit
layer.  For these choices, the stationary classes satisfy
\begin{equation}\label{eq:Pstar-wild-two}
 \beta_K\equiv
 \beta-\frac{\beta_1\alpha\delta_h}{\alpha_1}
 \pmod{\p_K^{d'}},
 \qquad
 \beta(\tau\chi_F)\equiv\beta+\alpha\delta_h
 \pmod{\p_F^d},
\end{equation}
and their norm-product class satisfies
\begin{equation}\label{eq:high-product-two}
 N(\beta_K)\equiv\beta(\beta+\alpha\delta_h)
 \pmod{\p_F^d}.
\end{equation}
The congruences are equalities of stationary classes.  The displayed
right-hand sides are selected representatives, not canonically
determined field elements.
\end{proposition}

\begin{proof}
Put
\[
 T=t+1,\qquad m=T+v=2d+\varepsilon,
 \qquad m'=2m-T=2d'+\varepsilon',
\]
and set
\[
 r=\left\lfloor\frac T2\right\rfloor,
 \qquad s=\left\lceil\frac T2\right\rceil.
\]
Then
\begin{equation}\label{eq:wild-two-depth-identities}
 d'+\varepsilon'=m-r,
 \qquad d'+\varepsilon'-v=s,
 \qquad r+s=T.
\end{equation}

Choose any representative $\beta$ of the stationary class of $\chi_F$.
By Theorem~\ref{thm:norm-filtration}(c), choose
$\beta_1\in\cO_K^\times$ such that
$N(\beta_1)/\beta\in U_F^t$, and put
\[
 \delta_h=\frac{\beta}{N(\beta_1)}\in U_F^t.
\]
Thus $\beta=\delta_hN(\beta_1)$ exactly.  In general one need not have
$d\le t$, so no congruence
$N(\beta_1)\equiv\beta\pmod{\p_F^d}$ is asserted.  The factor
$\delta_h$ is essential and will be retained in every subsequent
stationary-class formula.

Let $c_\tau$ be any representative of the stationary coefficient class
of $\tau$ relative to $\gamma_F$ on $\p_F^s$.  That class lies in
\[
 \p_F^v/\p_F^{v+r}.
\]
Apply Lemma~\ref{lem:norm-stationary-representatives} to the element
$c_\tau\delta_h^{-1}$ with $(h,s)=(v,r)$.  It supplies
$\alpha_1\in K^\times$ with $v_K(\alpha_1)=v$ such that, for
$\alpha=N(\alpha_1)$,
\[
 \alpha\delta_h\equiv c_\tau\pmod{\p_F^{v+r}}.
\]
Thus $\alpha\delta_h$ is a permitted representative and
\begin{equation}\label{eq:wild-two-tau-linearization}
 \tau(1+x)=
 \psi_F\!\left(\frac{\alpha\delta_hx}{\gamma_F}\right)
 \qquad(x\in\p_F^s).
\end{equation}
The choices of $\beta_1$ and $\alpha_1$ are independent; this proves
simultaneous existence without asserting uniqueness.

For $y\in\p_K^s$, the quadratic norm identity is exact:
\[
 N(1+y)=1+\Tr(y)+N(y).
\]
The trace-ideal formula puts both $\Tr(y)$ and $N(y)$ in $\p_F^s$.
Since $\tau$ is trivial on norms, applying
\eqref{eq:wild-two-tau-linearization} gives the conversion identity
\begin{equation}\label{eq:wild-two-conversion}
 \psi_F\!\left(\frac{\alpha\delta_hN(y)}{\gamma_F}\right)
 =\psi_F\!\left(-\frac{\alpha\delta_h\Tr(y)}{\gamma_F}\right).
\end{equation}

Let $x\in\p_K^{d'+\varepsilon'}$ and put
$y=\beta_1x/\alpha_1$.  By
\eqref{eq:wild-two-depth-identities}, $y\in\p_K^s$, and the identities
$\beta=\delta_hN(\beta_1)$ and $\alpha=N(\alpha_1)$ give
\[
 \beta N(x)=\alpha\delta_hN(y).
\]
Using \eqref{eq:wild-two-conversion} in the exact expansion
$P_{K/F}(x)=\Tr(x)+N(x)$ yields
\begin{align*}
 \psi_F\!\left(\frac{\beta P_{K/F}(x)}{\gamma_F}\right)
 &=\psi_F\!\left(
   \frac{\beta\Tr(x)-
    \alpha\delta_h\Tr(\beta_1x/\alpha_1)}{\gamma_F}\right)\\
 &=\psi_K\!\left(
   \frac{(\beta-\beta_1\alpha\delta_h/\alpha_1)x}{\gamma_F}
   \right).
\end{align*}
Here $\gamma_F$, viewed in $K$, is admissible because
$m'+n_K(\psi_K)=2(m+n_F(\psi_F))$.  By
Corollary~\ref{cor:stationary-class-under-norm}, this proves the first
class congruence in \eqref{eq:Pstar-wild-two}.

Since $d+\varepsilon\ge s$, the two stationary restrictions of $\tau$
and $\chi_F$ are defined on a common layer.  Lemma
\ref{lem:stationary-class-calculus}(c) therefore gives the second class
congruence in \eqref{eq:Pstar-wild-two}.  At $m=T$, its leading class
is nonzero by Lemma~\ref{lem:minimal-orbit-stationary}, so no hidden
conductor drop occurs.

Finally choose the normalized representative
\[
 \beta_K^0=\beta-\frac{\beta_1\alpha\delta_h}{\alpha_1}
\]
of the upstairs class, and put
$c=\beta_1\alpha\delta_h/\alpha_1$.  By
\eqref{eq:norm-of-stationary-class-well-defined}, it is enough to
compute $N(\beta_K^0)$ modulo $\p_F^d$.  The exact quadratic norm
formula gives
\[
 N(\beta-c)=\beta^2-\beta\Tr(c)+N(c),
 \qquad N(c)=\beta\alpha\delta_h.
\]
Now $v_K(c)=v$, and the trace-ideal formula, with different exponent
$T$, gives
\[
 v_F(\Tr(c))\ge
 \left\lfloor\frac{v+T}{2}\right\rfloor=d.
\]
Consequently
\[
 N(\beta_K)\equiv N(\beta_K^0)
 \equiv\beta^2+\beta\alpha\delta_h
 =\beta(\beta+\alpha\delta_h)\pmod{\p_F^d},
\]
which proves \eqref{eq:high-product-two} independently of the chosen
upstairs representative.
\end{proof}

\subsection{Low-conductor parameters}

Assume
\[
 2\le m=2d+\varepsilon\le T,
 \qquad v=T-m.
\]
Choose $\chi_F$ of minimal conductor in its $S(K/F)$-orbit; then
$m_K(\chi_K)=m$.  Put
\[
 r_0=\left\lfloor\frac T2\right\rfloor,
 \qquad s_0=\left\lceil\frac T2\right\rceil.
\]
Choose
\[
 \delta\cO_F=\p_F^{T+n_F(\psi_F)},
 \qquad \varepsilon_1\cO_K=\p_K^v,
 \qquad \varepsilon=N(\varepsilon_1),
\]
and
\begin{equation}\label{eq:gammas-low}
 \gamma_F=\delta/\varepsilon,
 \qquad \gamma_K=\delta/\varepsilon_1.
\end{equation}
Relative to $\delta$, the stationary class of $\tau$ on
$\p_F^{s_0}$ lies in $\cO_F/\p_F^{r_0}$; relative to $\gamma_F$, the
stationary class of $\chi_F$ lies in $\cO_F/\p_F^d$.  Since
$r_0,d\le t$, Lemma~\ref{lem:norm-stationary-representatives} permits
independent, hence simultaneous, choices
$\alpha_1,\beta_1\in U_K$ and exact norm representatives
\begin{equation}\label{eq:low-norm-representatives}
 \alpha=N(\alpha_1),\qquad \beta=N(\beta_1),
\end{equation}
such that
\begin{equation}\label{eq:low-mu}
 \tau(1+x)=\psi_F(\alpha x/\delta)
 \quad(x\in\p_F^{s_0}),
 \qquad
 \chi_F(1+x)=\psi_F(\beta x/\gamma_F)
 \quad(x\in\p_F^{d+\varepsilon}).
\end{equation}
For $j\ne0$, Lemma~\ref{lem:minimal-orbit-stationary} gives
$m_F(\tau^j\chi_F)=T$.  In particular, when $m=T$, the leading class
$j\alpha+\beta$ does not cancel.

\begin{proposition}[Low-conductor stationary-class table]
\label{prop:low-parameter-table}
The upstairs stationary class is
\begin{equation}\label{eq:low-beta-K}
 \beta_K:=\beta(\chi_K)
 \equiv\frac{\beta\varepsilon}{\varepsilon_1}
       -\frac{\beta_1\alpha}{\alpha_1}
 \pmod{\p_K^d}.
\end{equation}
For $j\in\F_p$, let $\widetilde\beta_j$ denote the stationary
numerator class of $\tau^j\chi_F$ after all factors are written with
the common additive denominator $\delta$.  Thus
$\widetilde\beta_0=\varepsilon\beta$, while for $j\ne0$ it is the
ordinary conductor-$T$ stationary numerator class relative to
$\delta$.  Then
\begin{equation}\label{eq:low-beta-twists}
 \widetilde\beta_j
 \equiv N(j\alpha_1+\varepsilon_1\beta_1)
 \equiv j\alpha+\varepsilon\beta
 \pmod{\p_F^{r_0}}.
\end{equation}
Every displayed right-hand side may be chosen as a representative of
the indicated class.  Moreover the norm of the upstairs class is
well defined modulo $\p_F^d$ and
\begin{equation}\label{eq:low-product-congruence}
 N(\beta_K)\equiv
 \beta\prod_{j\in\F_p^\times}(\varepsilon\beta+j\alpha)
 \pmod{\p_F^d}.
\end{equation}
\end{proposition}

\begin{proof}
Put $q=d+\varepsilon=\lceil m/2\rceil$.  Since
$2q+(p-1)T\ge pm$, Corollary~\ref{cor:wild-norm-truncation} gives
\begin{equation}\label{eq:low-norm-trunc}
 N(1+x)-1\equiv\Tr(x)+N(x)\pmod{\p_F^m}
 \qquad(x\in\p_K^q).
\end{equation}
For $y=\varepsilon_1\beta_1x/\alpha_1$ one has
\[
 v_K(y)\ge q+v=T-\floor{m/2}\ge s_0.
\]
The trace-ideal formula and total ramification give
\[
 v_F(\Tr(y))\ge
 \floor{\frac{v_K(y)+(p-1)T}{p}}
 \ge T-\ceil{\frac{\floor{m/2}}p}
 \ge\ceil{\frac T2}=s_0,
 \qquad v_F(N(y))=v_K(y)\ge s_0.
\]
Moreover
\[
 2v_K(y)+(p-1)T
 \ge2\bigl(T-\floor{m/2}\bigr)+(p-1)T
 \ge pT,
\]
so Corollary~\ref{cor:wild-norm-truncation}, now used with target
precision $T$, gives
\[
 N(1+y)\equiv1+\Tr(y)+N(y)\pmod{\p_F^T}.
\]
The character $\tau$ is trivial on norms and on $U_F^T$, while its
linearization is valid on $\p_F^{s_0}$.  Applying $\tau$ to this
congruence therefore gives
\begin{equation}\label{eq:low-conversion}
 \psi_F\!\left(-\frac{\alpha\Tr(y)}\delta\right)
 =\psi_F\!\left(\frac{\alpha N(y)}\delta\right)
 =\psi_F\!\left(\frac{\varepsilon\beta N(x)}\delta\right).
\end{equation}
For the element on the right of \eqref{eq:low-beta-K},
\begin{align*}
 \psi_K\!\left(\frac{\beta_Kx}{\gamma_K}\right)
 &=\psi_F\!\left(
   \frac{\Tr(\varepsilon_1\beta_Kx)}\delta\right)\\
 &=\psi_F\!\left(
   \frac{\varepsilon\beta\Tr(x)
         -\alpha\Tr(\varepsilon_1\beta_1x/\alpha_1)}\delta\right).
\end{align*}
Using \eqref{eq:low-conversion} and
\eqref{eq:low-norm-trunc}, this is exactly the pullback of the
stationary character of $\chi_F$.  Corollary
\ref{cor:stationary-class-under-norm} proves
\eqref{eq:low-beta-K} as an equality of stationary classes.

For $j=0$, the first congruence in
\eqref{eq:low-beta-twists} is the exact equality
$N(\varepsilon_1\beta_1)=\varepsilon\beta$.  For $j\ne0$, put
$w=\varepsilon_1\beta_1/(j\alpha_1)$, so $v_K(w)=v$.  The strong
symmetric estimate gives
\[
 v_F(s_i(w))\ge
 \left\lfloor\frac{iv+(p-1)T}{p}\right\rfloor\ge r_0
 \qquad(1\le i\le p-1),
\]
and therefore
$N(j\alpha_1+\varepsilon_1\beta_1)
 \equiv j\alpha+\varepsilon\beta\pmod{\p_F^{r_0}}$.
By Lemma~\ref{lem:stationary-class-calculus}(c), this is precisely the
sum of the two stationary classes with common denominator $\delta$.
Its exact conductor is $T$ for $j\ne0$ by
Lemma~\ref{lem:minimal-orbit-stationary}; this explicitly excludes a
boundary conductor drop.

Put
\[
 u=\frac{\varepsilon_1\beta_1}{\alpha_1},
 \qquad n=N(u)=\frac{\varepsilon\beta}{\alpha},
\]
and choose the normalized representative
\[
 \beta_K^0=\frac{\alpha}{\varepsilon_1}(n-u)
\]
of the class in \eqref{eq:low-beta-K}.  By
\eqref{eq:norm-of-stationary-class-well-defined}, $N(\beta_K)$ modulo
$\p_F^d$ may be computed from this representative.  Thus
\[
 N(\beta_K)\equiv N(\beta_K^0)
 =\frac{\alpha^p}{\varepsilon}N(n-u)\pmod{\p_F^d}.
\]
In the expansion of $N(n-u)$ the two endpoints are
$n^p+(-1)^pn$.  For $1\le i\le p-1$, the intermediate term
$n^{p-i}s_i(u)$, after multiplication by
$\alpha^p/\varepsilon$, has valuation at least
\[
 L_i=(p-i-1)v+
 \left\lfloor\frac{iv+(p-1)T}{p}\right\rfloor.
\]
Because $T$ is an integer,
\[
 \left\lfloor\frac{iv+(p-1)T}{p}\right\rfloor
 =T-\left\lceil\frac{T-iv}{p}\right\rceil.
\]
Moreover $i\ge1$ and $v=T-m\ge0$, so
$T-iv\le T-v=m=2d+\varepsilon$.  Since $p\ge2$,
\[
 \left\lceil\frac{T-iv}{p}\right\rceil
 \le\left\lceil\frac m2\right\rceil=d+\varepsilon.
\]
Consequently
\[
 \begin{aligned}
 L_i
 &\ge(p-i-1)v+T-d-\varepsilon\\
 &=(p-i)v+d\\
 &\ge d.
 \end{aligned}
\]
Thus every intermediate term vanishes modulo $\p_F^d$, and only the
endpoints survive.  Since the Teichmuller representatives satisfy
\[
 \prod_{j\in\F_p}(X+j)=X^p+(-1)^pX,
\]
we obtain
\[
 N(\beta_K)\equiv
 \frac{\alpha^p}{\varepsilon}\bigl(n^p+(-1)^pn\bigr)
 =\beta\prod_{j\ne0}(\varepsilon\beta+j\alpha)
 \pmod{\p_F^d}.
\]
This proves \eqref{eq:low-product-congruence} independently of the
chosen representative of the upstairs class.
\end{proof}

\begin{remark}[Characteristic independence of the stationary-class tables]
The stationary-class calculus, the proofs of
Propositions~\ref{prop:high-parameter-table} and
\ref{prop:low-parameter-table}, and the representative-independence
argument use only the cyclic-prime norm filtration, trace duality, the
exact norm polynomial, and Lemma~\ref{lem:wild-symmetric-bound}.
Hence the stationary-class tables, the compatible representatives
selected from them, and the residual calculations based on the complete
stationary factors are valid in both mixed and equal characteristic.
\end{remark}

\section{Unramified and tame cases}

\begin{proposition}[Unramified case]\label{prop:unramified-new}
The First Main Lemma holds when $K/F$ is unramified.
\end{proposition}

\begin{proof}
Put
\[
 k=k_F,\qquad k'=k_K,\qquad
 n=n_F(\psi_F),\qquad m=m_F(\chi_F).
\]
Choose a uniformizer $\pi$ of $F$.  By
Lemma~\ref{lem:unramified-compatibility}, it is also a uniformizer of
$K$, the residue extension $k'/k$ has degree $\ell$, and
\begin{equation}\label{eq:unramified-common-conductors}
 n_K(\psi_K)=n,\qquad m_K(\chi_K)=m.
\end{equation}

\smallskip
\noindent\emph{Step 1: reduction to a power identity.}
Lemma~\ref{lem:unramified-compatibility}(e) identifies
$F^\times/N(K^\times)$ with the cyclic group generated by the class of
$\pi$, of order $\ell$.  Hence every $\mu\in S(K/F)$ is unramified,
and the values $\mu(\pi)$ run through all $\ell$th roots of unity.
Twisting by such a $\mu$ does not change the conductor of $\chi_F$.
Use the same admissible element $\gamma_F=\pi^{m+n}$ for $\chi_F$ and
$\mu\chi_F$.  Since $\mu$ is trivial on $U_F^0$, their unit integrals
are equal, and therefore
\begin{equation}\label{eq:unramified-twist-delta}
 \Delta_F(\mu\chi_F,\psi_F)
 =\mu(\pi)^{m+n}\Delta_F(\chi_F,\psi_F).
\end{equation}
For $\mu$ itself, \eqref{eq:unramified-character-local-constant} gives
\begin{equation}\label{eq:unramified-mu-delta}
 \Delta_F(\mu,\psi_F)=\mu(\pi)^n.
\end{equation}
Both formulas remain valid when $m=0$.  The values $\mu(\pi)$ are
the roots of $X^\ell-1$, whose product is $(-1)^{\ell-1}$; hence
\[
 \prod_{\mu\in S(K/F)}\mu(\pi)=(-1)^{\ell-1}.
\]
The quotient of the product on the right side of
\eqref{eq:first-main} by the product of norm-character factors on its
left side is
\[
 (-1)^{m(\ell-1)}\Delta_F(\chi_F,\psi_F)^\ell.
\]
Thus the First Main Lemma is equivalent to
\begin{equation}\label{eq:unramified-power-new}
 \Delta_K(\chi_K,\psi_K)=
 (-1)^{m(\ell-1)}\Delta_F(\chi_F,\psi_F)^\ell.
\end{equation}

\smallskip
\noindent\emph{Step 2: conductor $m=0$.}
Formula \eqref{eq:unramified-character-local-constant} gives
\[
 \Delta_F(\chi_F,\psi_F)=\chi_F(\pi)^n,
 \qquad
 \Delta_K(\chi_K,\psi_K)=\chi_K(\pi)^n.
\]
By Lemma~\ref{lem:unramified-compatibility}(c),
$N_{K/F}(\pi)=\pi^\ell$, so
$\chi_K(\pi)=\chi_F(\pi)^\ell$.  This proves
\eqref{eq:unramified-power-new} for $m=0$.

\smallskip
\noindent\emph{Step 3: conductor $m=1$.}
Let
\[
 \bar\chi(\bar u)=\chi_F(u),
 \qquad
 \bar\psi(\bar x)=\psi_F(x/\pi^{n+1})
\]
be the residual multiplicative and additive characters on $k$.  By
\eqref{eq:unramified-residue-trace-norm}, the corresponding characters
on $k'$ are
\begin{equation}\label{eq:unramified-residual-lifts}
 \bar\chi_K=\bar\chi\circ N_{k'/k},
 \qquad
 \bar\psi_K=\bar\psi\circ\Tr_{k'/k}.
\end{equation}
The conductor-one formula \eqref{eq:conductor-one-local-gauss} gives
\[
 \Delta_F(\chi_F,\psi_F)
 =-\chi_F(\pi)^{n+1}
   \Aphase\!\left(\tau_k(\bar\chi,\bar\psi)\right)
\]
and
\[
 \Delta_K(\chi_K,\psi_K)
 =-\chi_K(\pi)^{n+1}
   \Aphase\!\left(\tau_{k'}(\bar\chi_K,\bar\psi_K)\right).
\]
Theorem~\ref{thm:HD-lift} and
\eqref{eq:unramified-residual-lifts} imply
\[
 \tau_{k'}(\bar\chi_K,\bar\psi_K)
 =\tau_k(\bar\chi,\bar\psi)^\ell.
\]
Since $\Aphase(z^\ell)=\Aphase(z)^\ell$ for $z\ne0$, and since
$\chi_K(\pi)=\chi_F(\pi)^\ell$, this yields
\[
 \Delta_K(\chi_K,\psi_K)
 =(-1)^{\ell-1}\Delta_F(\chi_F,\psi_F)^\ell,
\]
which is \eqref{eq:unramified-power-new} because $m=1$.

\smallskip
\noindent\emph{Step 4: conductor $m>1$.}
Write $m=2d+\varepsilon$, with $\varepsilon\in\{0,1\}$, and put
$\gamma=\pi^{m+n}$.  Proposition~\ref{prop:high-parameter-table}(a)
says that the stationary classes of $\chi_F$ and $\chi_K$ admit the
same representative $\beta\in\cO_F^\times$.  Choose this representative
for both characters and define
\[
 B_E=\chi_E(\gamma)\psi_E(\beta/\gamma)
      \chi_E(\beta)^{-1}\qquad(E=F,K).
\]
The factorization \eqref{eq:Delta-factorization-new} reads
\begin{equation}\label{eq:unramified-Lamprecht-factorization}
 \Delta_E(\chi_E,\psi_E)=B_E\Delta_{3,E}(\chi_E).
\end{equation}
Since $\beta,\gamma\in F$, one has
\[
 N_{K/F}(\beta)=\beta^\ell,
 \quad N_{K/F}(\gamma)=\gamma^\ell,
 \quad \Tr_{K/F}(\beta/\gamma)=\ell\beta/\gamma.
\]
Moreover $\psi_F(\ell z)=\psi_F(z)^\ell$ for every $z\in F$, because
$\psi_F$ is an additive character.  Therefore
\begin{equation}\label{eq:unramified-elementary-power}
 B_K=B_F^\ell.
\end{equation}
If $\varepsilon=0$, both residual factors are $1$ by definition, and
\eqref{eq:unramified-power-new} follows from
\eqref{eq:unramified-Lamprecht-factorization}--
\eqref{eq:unramified-elementary-power}; its sign is $1$ because $m$ is
even.

It remains to treat $\varepsilon=1$.  Then $m=2d+1$ and $d\geq1$.
Take $\delta=\pi^d$.  Let $\phi_F:k\to\C^\times$ and
$\phi_K:k'\to\C^\times$ be the Hasse functions from
\eqref{eq:Lamprecht-Hasse-function}:
\[
 \phi_F(\bar x)=
 \psi_F(\delta\beta x/\gamma)\chi_F(1+\delta x)^{-1},
\]
\[
 \phi_K(\bar x)=
 \psi_K(\delta\beta x/\gamma)\chi_K(1+\delta x)^{-1}.
\]
The residual additive character in
\eqref{eq:Hasse-functional-equation} is
\begin{equation}\label{eq:unramified-Hasse-additive-character}
 \psi_0(\bar z)=\psi_F(\beta\delta^2z/\gamma).
\end{equation}
Indeed, expanding $(1+\delta x)(1+\delta y)$ modulo $\p_F^m$ in the
proof of Theorem~\ref{thm:Lamprecht} gives this character.

For $x\in\cO_K$, let $E_j(x)$ be the $j$th elementary symmetric
function of its $F$-conjugates.  Thus
$E_1(x)=T(x):=\Tr_{K/F}(x)$ and
\[
 E_2(x)=\sum_{\{\sigma,\tau\}}x^\sigma x^\tau,
\]
the sum being over unordered pairs of distinct $F$-embeddings of $K$.
The norm expansion is
\[
 N(1+\delta x)
 =1+\delta T(x)+\delta^2E_2(x)
   +\sum_{j=3}^{\ell}\delta^jE_j(x),
\]
where the last sum is empty if $\ell=2$.  Each term in that sum lies in
$\p_F^{3d}\subseteq\p_F^{2d+1}$, and the same is true of the cross term
$\delta^3T(x)E_2(x)$.  Hence
\begin{equation}\label{eq:unramified-Hasse-norm-expansion}
 N(1+\delta x)
 \equiv(1+\delta T(x))(1+\delta^2E_2(x))
 \pmod{\p_F^{m}}.
\end{equation}
Because $\delta^2E_2(x)\in\p_F^{2d}\subseteq\p_F^{d+1}$,
Lamprecht linearization \eqref{eq:Lamprecht-linearization} gives
\[
 \chi_F(1+\delta^2E_2(x))^{-1}
 =\psi_F(-\beta\delta^2E_2(x)/\gamma)
 =\psi_0(-\overline{E_2(x)}).
\]
Substituting this and \eqref{eq:unramified-Hasse-norm-expansion} in the
definition of $\phi_K$ yields
\begin{equation}\label{eq:unramified-Hasse-lift-identification}
 \phi_K(\bar x)
 =\phi_F\!\left(\Tr_{k'/k}(\bar x)\right)
  \psi_0\!\left(-\omega_{k'/k}(\bar x)\right).
\end{equation}
Here \eqref{eq:unramified-residue-trace-norm} identifies
$\overline{T(x)}$ with the residue-field trace.  The same reduction of
the characteristic polynomial, now using its second coefficient,
identifies $\overline{E_2(x)}$ with the function $\omega_{k'/k}$ in
Theorem~\ref{thm:Hasse-lift}.  Thus
\eqref{eq:unramified-Hasse-lift-identification} says exactly that
$\phi_K$ is the Hasse-function lift of $\phi_F$.

Theorem~\ref{thm:Hasse-lift} gives
\[
 -\sum_{x\in k'}\phi_K(x)
 =\left(-\sum_{x\in k}\phi_F(x)\right)^\ell,
\]
or equivalently
\[
 \sum_{x\in k'}\phi_K(x)
 =(-1)^{\ell-1}\left(\sum_{x\in k}\phi_F(x)\right)^\ell.
\]
Taking phases gives
\begin{equation}\label{eq:unramified-residual-power}
 \Delta_{3,K}(\chi_K)
 =(-1)^{\ell-1}\Delta_{3,F}(\chi_F)^\ell.
\end{equation}
Combining \eqref{eq:unramified-Lamprecht-factorization},
\eqref{eq:unramified-elementary-power}, and
\eqref{eq:unramified-residual-power} proves
\eqref{eq:unramified-power-new}.  Since $m$ is odd in this last case,
$(-1)^{\ell-1}=(-1)^{m(\ell-1)}$.  This completes the proof.
\end{proof}

\begin{proposition}[Tame conductor-one local-to-residue comparison]
\label{prop:tame-conductor-one-comparison}
Let $K/F$ be tamely ramified cyclic of prime degree $\ell$ and let
$\chi_F$ have conductor one and minimal conductor in its
$S(K/F)$-orbit.  Then $m_K(\chi_F\circ N)=1$.  After choosing compatible
uniformizers and reducing the conductor-one local constants by
\eqref{eq:conductor-one-local-gauss}, the First Main Lemma is exactly
the Hasse--Davenport multiplication identity
\eqref{eq:HD-product} for the residue character of $\chi_F$ and the
order-$\ell$ residue character attached to $K/F$.  In particular the
First Main Lemma holds in this case.
\end{proposition}

\begin{proof}
The break is $t=0$.  Proposition~\ref{prop:pullback-conductor}(iii)
shows that $m_K(\chi_F\circ N)\le1$; if it were $0$, some twist of
$\chi_F$ by a member of $S(K/F)$ would have conductor $0$, contrary to
the minimal choice of $\chi_F$.  Thus the upper conductor is exactly
one.

Choose $\varpi_K$ and put $\varpi_F=N(\varpi_K)$.  Let
$n=n_F(\psi_F)$.  Since the tame different has exponent $\ell-1$,
\[
 n_K(\psi_K)=\ell n+\ell-1.
\]
The element
\[
 \Gamma=\varpi_F^{n+1}
\]
has $K$-valuation $\ell(n+1)=n_K(\psi_K)+1$, so it is an admissible
element for every conductor-one factor over both fields.  Define
\[
 \bar\psi(x)=\psi_F(\widetilde x/\Gamma),
 \qquad
 \bar\chi(x)=\chi_F(\widetilde x)
 \quad(x\in k^\times),
\]
where $k$ is the common residue field.  On residue classes the norm and
trace are
\[
 \overline{N_{K/F}(u)}=\bar u^{\,\ell},
 \qquad
 \overline{\Tr_{K/F}(u)}=\ell\bar u.
\]
Hence the upper residual characters are
\begin{equation}\label{eq:tame-one-upper-residual-characters}
 \bar\chi_K=\bar\chi^{\,\ell},
 \qquad
 \bar\psi_K(x)=\bar\psi(\ell x).
\end{equation}

Let $S(K/F)=\{\mu_j:0\le j<\ell\}$, with $\mu_0=1$, and let
$\eta$ be the residue character of $\mu_1$.  Then $\eta$ has order
$\ell$ and the residue character of $\mu_j$ is $\eta^j$.  Moreover
$\mu_j(\varpi_F)=1$, because $\varpi_F=N(\varpi_K)$.  Every
$\mu_j\chi_F$ has conductor one by minimality.

The coset calculation proving
\eqref{eq:conductor-one-local-gauss}, used with the common admissible
element $\Gamma$, gives
\begin{align*}
 \Delta_K(\chi_K,\psi_K)
 &=-\chi_F(\Gamma)^\ell
   \Aphase\!\left(\tau_k(\bar\chi^{\,\ell},
                              \bar\psi(\ell\,\cdot))\right),\\
 \Delta_F(\mu_j\chi_F,\psi_F)
 &=-\chi_F(\Gamma)
   \Aphase\!\left(\tau_k(\bar\chi\eta^j,\bar\psi)\right),\\
 \Delta_F(\mu_j,\psi_F)
 &=-\Aphase\!\left(\tau_k(\eta^j,\bar\psi)\right)
 \quad(1\le j<\ell),
\end{align*}
while $\Delta_F(1,\psi_F)=1$.  The powers of $\chi_F(\Gamma)$ and
the signs are the same on the two sides of the First Main Lemma and
cancel.

For a multiplicative character $\alpha$ and $c\in k^\times$, the
definition \eqref{eq:tau-definition} gives
\[
 \tau_k(\alpha,\bar\psi(c\,\cdot))
 =\alpha(c)\tau_k(\alpha,\bar\psi).
\]
Thus
\[
 \tau_k(\bar\chi^{\,\ell},\bar\psi(\ell\,\cdot))
 =\bar\chi(\ell^\ell)
   \tau_k(\bar\chi^{\,\ell},\bar\psi).
\]
After the preceding cancellations, the First Main Lemma is therefore
exactly
\[
 \bar\chi(\ell^\ell)
 \tau_k(\bar\chi^{\,\ell},\bar\psi)
 \prod_{j=1}^{\ell-1}\tau_k(\eta^j,\bar\psi)
 =
 \tau_k(\bar\chi,\bar\psi)
 \prod_{j=1}^{\ell-1}\tau_k(\bar\chi\eta^j,\bar\psi),
\]
which is the Hasse--Davenport multiplication formula
\eqref{eq:HD-product}.  This proves the proposition.
\end{proof}

\begin{proposition}[Tamely ramified odd degree]\label{prop:tame-odd-new}
The First Main Lemma holds for a tamely ramified cyclic extension of odd
prime degree.
\end{proposition}

\begin{proof}
Here $t=0$ and every nontrivial norm character has conductor $1$.
Put $n=n_F(\psi_F)$.  For $m=0$, choose a uniformizer
$\varpi_K$ of $K$ and set $\varpi_F=N(\varpi_K)$; this is a uniformizer
of $F$.  Since
$n_K=\ell n+(\ell-1)$, formula
\eqref{eq:unramified-character-local-constant} gives
\[
 \Delta_K(\chi_K)=
 \chi_F(\varpi_F)^{\ell n+\ell-1}.
\]
Every nontrivial $\mu\in S(K/F)$ is trivial on $\varpi_F$ and has
conductor one, so \eqref{eq:conductor-one-local-gauss} gives
\[
 \Delta_F(\mu\chi_F)=
 \chi_F(\varpi_F)^{n+1}\Delta_F(\mu).
\]
Multiplying this for the $\ell-1$ nontrivial $\mu$ and including
$\Delta_F(\chi_F)=\chi_F(\varpi_F)^n$ gives exponent
$n+(\ell-1)(n+1)=\ell n+\ell-1$; the factors
$\prod_{\mu\ne1}\Delta_F(\mu)$ agree on both sides.  This proves the
unramified-character case.

For $m=1$, apply Proposition~\ref{prop:tame-conductor-one-comparison}; it identifies every
local scaling and reduces the equality exactly to
Theorem~\ref{thm:HD-product}.
If $m=2d+\varepsilon>1$, then $1\leq d$ and
Lemma~\ref{lem:stable-twist-new} gives
\[
 \prod_{\mu\in S}\Delta_F(\mu\chi_F,\psi_F)
 =\Delta_F(\chi_F,\psi_F)^\ell,
\]
because the product of all characters in the odd-order group $S$ is
trivial.  Corollary~\ref{cor:S-product-one} removes the other norm
characters.  The parameter comparison in
Proposition~\ref{prop:high-parameter-table}(b) and
Theorem~\ref{thm:ramified-Hasse-comparison} give
$\Delta_K(\chi_K,\psi_K)=\Delta_F(\chi_F,\psi_F)^\ell$.
\end{proof}

\begin{proposition}[Tamely ramified quadratic case]\label{prop:tame-two-new}
The First Main Lemma holds for a tamely ramified quadratic extension.
\end{proposition}

\begin{proof}
Put $n=n_F(\psi_F)$ and write $S(K/F)=\{1,\mu\}$.  The residue
characteristic is odd.  Choose a uniformizer $\varpi_K$ with
$\sigma(\varpi_K)=-\varpi_K$ and put
$\varpi_F=N(\varpi_K)=-\varpi_K^2$.  Then $\mu(\varpi_F)=1$, and the
restriction of $\mu$ to $\cO_F^\times$ is the quadratic character
$\nu$ of the common residue field.

If $m_F(\chi_F)=0$, then $n_K(\psi_K)=2n+1$ and
\[
 \Delta_K(\chi_K)=\chi_F(\varpi_F)^{2n+1},
\]
while
\[
 \Delta_F(\chi_F)\Delta_F(\mu\chi_F)
 =\chi_F(\varpi_F)^n
  \chi_F(\varpi_F)^{n+1}\Delta_F(\mu).
\]
Multiplying the first display by the remaining factor $\Delta_F(\mu)$
on the left-hand side of \eqref{eq:first-main} gives exactly the second
display.  Thus this is the desired formula.  If $m_F(\chi_F)=1$,
Proposition~\ref{prop:tame-conductor-one-comparison}, specialized to
$\ell=2$, supplies the exact uniformizer, different, norm, trace, and
residual-character identifications and reduces the assertion to
Theorem~\ref{thm:HD-product}.

Now let $m=2d+\varepsilon\ge2$ and put
\[
 \gamma=\varpi_F^{m+n}.
\]
This element is admissible for the two downstairs characters.  It is
also admissible upstairs, because in the tame quadratic case
\[
 m_K(\chi_K)=2m-1,\qquad n_K(\psi_K)=2n+1,
\]
and hence
\[
 v_K(\gamma)=2(m+n)=m_K(\chi_K)+n_K(\psi_K).
\]
Since $\mu$ is trivial on $U_F^d$,
Lemma~\ref{lem:stationary-class-calculus}(c) identifies the stationary
classes of $\chi_F$ and $\mu\chi_F$ on the Lamprecht layer;
Proposition~\ref{prop:high-parameter-table}(d) identifies the pullback
class upstairs.  Choose the same representative $\beta$ for all three
classes, translating an odd critical function when required by
Lemma~\ref{lem:complete-stationary-invariant}.  Then
Lemma~\ref{lem:stable-twist-new} gives
\begin{equation}\label{eq:tame-stable-twist}
 \Delta_F(\mu\chi_F)=\mu(\gamma/\beta)\Delta_F(\chi_F).
\end{equation}
Since $\gamma$ is a power of the norm $\varpi_F$, one has
$\mu(\gamma)=1$.  Since $\mu|_{\cO_F^\times}=\nu$ and $\nu$ is
quadratic,
\[
 \mu(\gamma/\beta)=\mu(\beta)^{-1}
 =\nu(\bar\beta)^{-1}=\nu(\bar\beta).
\]
By \eqref{eq:Delta-factorization-new}, the nonresidual factor upstairs
is exactly the square of the corresponding one downstairs:
\begin{align*}
 \chi_K(\gamma)\Delta_{2,K}(\chi_K)
 &=\chi_F(\gamma^2)\,
   \psi_F(2\beta/\gamma)\,\chi_F(\beta^2)^{-1}\\
 &=\bigl(\chi_F(\gamma)\Delta_{2,F}(\chi_F)\bigr)^2.
\end{align*}
Indeed, $N_{K/F}(a)=a^2$ and $\Tr_{K/F}(a)=2a$ for $a\in F$.
Thus \eqref{eq:tame-stable-twist}, together with $\Delta_F(1)=1$,
cancels all nonresidual factors in \eqref{eq:first-main}.  It remains
to prove
\begin{equation}\label{eq:tame-quadratic-residual}
 \Delta_{3,K}(\chi_K)\Delta_F(\mu)
 =\nu(\bar\beta)\Delta_{3,F}(\chi_F)^2.
\end{equation}
Let
\[
 \psi_0(\bar x)=\psi_F(x/\varpi_F^{n+1}),
 \qquad G=\sum_{x\in k^\times}\nu(x)\psi_0(x).
\]
Then $\Delta_F(\mu)=\Aphase(G)$ and
$\Aphase(G)^2=\nu(-1)$.

If $m=2d$, then $m_K(\chi_K)=4d-1$.  Taking
$\delta_K=\varpi_K^{2d-1}$, the trace term in the Hasse function is zero
and
\[
 N(1+\delta_Kx)=1+\varpi_F^{2d-1}x^2.
\]
Consequently
\[
 \Delta_{3,K}=\Aphase\!\left(\sum_{x\in k}\psi_0(-\bar\beta x^2)\right)
 =\nu(-\bar\beta)\Aphase(G).
\]
Multiplying by $\Delta_F(\mu)$ gives
$\nu(-\bar\beta)\nu(-1)=\nu(\bar\beta)$, proving \eqref{eq:tame-quadratic-residual}.

If $m=2d+1$, the Hasse function over $F$ has the form
\[
 \phi_F(x)=\psi_0(\bar\beta x^2/2+cx)
\]
for a uniquely determined $c\in k$.  Here
$m_K(\chi_K)=4d+1$.  With
$\delta_K=\varpi_K^{2d}=(-1)^d\varpi_F^d$, direct substitution gives
\[
 \phi_K(x)=\phi_F((-1)^d x)^2
 =\psi_0(\bar\beta x^2+2(-1)^dcx).
\]
Lemma~\ref{lem:quadratic-phase-new} now gives
\begin{align*}
 \Delta_{3,K}
 &=\nu(2\bar\beta)\psi_0(-c^2/\bar\beta)Q_{\psi_0}(1,0),\\
 \Delta_{3,F}^{\,2}
 &=\psi_0(-c^2/\bar\beta)\nu(-1),\\
 \Delta_F(\mu)&=\nu(2)Q_{\psi_0}(1,0).
\end{align*}
Since $Q_{\psi_0}(1,0)^2=\nu(-1)$, these three formulas imply
\eqref{eq:tame-quadratic-residual}.  Combining with \eqref{eq:tame-stable-twist} completes
the proof.
\end{proof}

\section{Wild odd degree: stable range}

\begin{proposition}\label{prop:wild-stable-new}
Let $K/F$ be wildly ramified of odd prime degree $p$, with break $t$.
If $m_F(\chi_F)=2d+\varepsilon>1$ and $d\geq t+1$, then the First Main
Lemma holds.
\end{proposition}

\begin{proof}
Every nontrivial norm character has conductor $t+1\leq d$.  The stable
twist formula and Corollary~\ref{cor:S-product-one} reduce the theorem to
\[
 \Delta_K(\chi_K,\psi_K)=\Delta_F(\chi_F,\psi_F)^p.
\]
The elementary factors agree by
Proposition~\ref{prop:high-parameter-table}(b); note that
$m'+n_K=p(m+n_F)$ by \eqref{eq:pullback-high} and
\eqref{eq:additive-conductor-trace}.  If the conductor is even there is
no residual factor.  If it is odd,
Theorem~\ref{thm:ramified-Hasse-comparison} gives the required
$p$th-power identity for the residual factors.
\end{proof}

\section{The non-stable wild calculation}

The only remaining cases are wild and non-stable.  We now give the
complete local calculation.  Propositions
\ref{prop:high-parameter-table} and \ref{prop:low-parameter-table}
determine the stationary classes and exhibit compatible representatives
in the high- and low-conductor ranges, with the quadratic high range
governed by
Proposition~\ref{prop:wild-quadratic-high-parameters}; the resulting finite phases and valuation estimates are
then evaluated explicitly.

\subsection{The endpoint conductors in wild prime degree}\label{sec:wild-endpoints}
Before separating odd and quadratic degree, we dispose of the two
conductors for which there is no stationary layer.  The calculation is
valid for every wildly ramified cyclic extension $K/F$ of prime degree
$p$, including $p=2$.  Let $t$ be its break, put
$n=n_F(\psi_F)$, and choose uniformizers with
$\varpi_F=N_{K/F}(\varpi_K)$.  Then
\begin{equation}\label{eq:endpoint-additive-conductor}
 n_K(\psi_K)=pn+(p-1)(t+1)
\end{equation}
by Corollary~\ref{cor:additive-conductor}.

\medskip
\noindent\emph{Conductor $0$.}
If $\chi_F$ is unramified, then
\[
 \Delta_K(\chi_K,\psi_K)
 =\chi_F(\varpi_F)^{n_K(\psi_K)}.
\]
For $j\in\F_p$, let $\tau^j$ run through $S(K/F)$, with
$\tau^0=1$.  The character $\tau^j$ has conductor $t+1$ for
$j\ne0$, whereas $\tau^0$ has conductor $0$.  Since an unramified
character is constant on the unit group, direct substitution in
\eqref{eq:def-delta} gives
\[
 \Delta_F(\chi_F\tau^j,\psi_F)
 =\chi_F(\varpi_F)^{m_F(\tau^j)+n}
  \Delta_F(\tau^j,\psi_F).
\]
Multiplication over $j$ gives the exponent
\[
 n+(p-1)(t+1+n)=pn+(p-1)(t+1)=n_K(\psi_K),
\]
so the First Main Lemma follows.

\medskip
\noindent\emph{Conductor $1$.}
We now give the finite-sum comparison, retaining every local-to-residue
scaling factor.
Choose an admissible element $\gamma_F$ for $(\chi_F,\psi_F)$, so
$v_F(\gamma_F)=n+1$, and use it to write the residual additive
character as
\[
 \psi_0(\bar x)=\psi_F(x/\gamma_F),\qquad x\in\cO_F.
\]
After the normalization of Lemma~\ref{lem:residual-normalization},
$\psi_0(x^p)=\psi_0(x)$.

Let $f(T)$ be the minimal polynomial of $\varpi_K$ and put
\begin{equation}\label{eq:endpoint-gamma-K}
 \gamma_K=\gamma_F\frac{f'(\varpi_K)}{\varpi_K^{p-1}}.
\end{equation}
The different formula gives
$v_K(f'(\varpi_K))=(p-1)(t+1)$, so
$v_K(\gamma_K)=1+n_K(\psi_K)$; hence $\gamma_K$ is admissible for
$(\chi_K,\psi_K)$.  For $0\le i\le p-1$, partial fractions for
$T^i/f(T)$ at infinity give
\begin{equation}\label{eq:endpoint-lagrange-trace}
 \Tr_{K/F}\!\left(\frac{\varpi_K^i}{f'(\varpi_K)}\right)
 =\begin{cases}0,&i<p-1,\\1,&i=p-1.\end{cases}
\end{equation}
If $x\in\cO_K$ has residue $a\in k$, then
$x-a\in\p_K$ and the trace-ideal formula \eqref{eq:trace-ideal}
shows that
\[
 \Tr_{K/F}\!\left(\frac{(x-a)\varpi_K^{p-1}}
                         {f'(\varpi_K)}\right)\in\p_F.
\]
Consequently the residual additive character determined by
$\gamma_K$ is again $\psi_0$.  If $\bar\chi$ denotes the residue-field
character of $\chi_F$, then that of $\chi_K$ is
$x\mapsto\bar\chi(x^p)$.  Frobenius invariance of $\psi_0$ therefore
identifies the two finite sums by the change of variable $x\mapsto x^p$.
Thus
\begin{equation}\label{eq:endpoint-upper-ratio}
 \frac{\Delta_K(\chi_K,\psi_K)}
      {\Delta_F(\chi_F,\psi_F)}
 =\chi_F\!\left(\frac{N(\gamma_K)}{\gamma_F}\right).
\end{equation}

It remains to identify the residue class of the element on the right.
Let $\sigma$ generate $\Gal(K/F)$ and set
\[
 \lambda=
 \overline{\frac{\sigma(\varpi_K)-\varpi_K}
                   {\varpi_K^{t+1}}}\in k^\times.
\]
The ramification map sends $\sigma^j$ to $j\lambda$.  Proposition~\ref{prop:critical-norm-polynomial} gives, in these
source and target coordinates,
\begin{equation}\label{eq:endpoint-critical-norm}
 P(Z)=Z^p-\lambda^{p-1}Z,
 \qquad
 \operatorname{im}P=\lambda^p\ker(\Tr_{k/\F_p}).
\end{equation}
Choose an admissible
$\gamma_\tau$ for $\tau$ and a representative $b_\tau$ of its
stationary numerator class such that
\[
 \tau(1+x)=\psi_F(b_\tau x/\gamma_\tau)
\]
on the Lamprecht layer, and put $c_\tau=\gamma_\tau/b_\tau$.
\begin{lemma}[Endpoint annihilator and discriminant congruence]
\label{lem:endpoint-discriminant-congruence}
With the notation of the conductor-$1$ endpoint calculation, the
associated reciprocal stationary element
$c_\tau=\gamma_\tau/b_\tau$ and the ramification coefficient $\lambda$
satisfy
\begin{equation}\label{eq:endpoint-c-tau}
 \left(\overline{\frac{c_\tau}{\varpi_F^t\gamma_F}}\right)^{p-1}
 =\lambda^{p(p-1)}.
\end{equation}
Moreover, for
$D=f'(\varpi_K)/\varpi_K^{(p-1)(t+1)}$, one has
\[
 \bar D=(-1)^p\lambda^{p-1},\qquad
 \overline{N(D)}=(-1)^p\lambda^{p(p-1)}.
\]
\end{lemma}

\begin{proof}
The critical norm polynomial is
$P(Z)=Z^p-\lambda^{p-1}Z$.  Its image is
$\lambda^p\ker(\Tr_{k/\F_p})$, and the annihilator of this image
under the trace pairing is $\lambda^{-p}\F_p$.  On the critical
unit quotient write $x=\varpi_F^t\widetilde z$.  Since
$\tau(1+x)=\psi_F(x/c_\tau)$ and
$\psi_0(\bar y)=\psi_F(y/\gamma_F)$, the residual coefficient of
$\tau$ is
\[
 a_\tau=
 \overline{\frac{\varpi_F^t\gamma_F}{c_\tau}}.
\]
The character $\tau$ is trivial on the norm image.  Hence
$\psi_0(a_\tau P(z))=1$ for every $z\in k$, so
$a_\tau\in\lambda^{-p}\F_p$.  It is nonzero because $\tau$ has
conductor $t+1$; consequently
$a_\tau=\lambda^{-p}u$ for some $u\in\F_p^\times$.  Taking
inverses and then $(p-1)$st powers gives
\[
 \left(\overline{\frac{c_\tau}
 {\varpi_F^t\gamma_F}}\right)^{p-1}
 =(\lambda^p u^{-1})^{p-1}
 =\lambda^{p(p-1)},
\]
which is \eqref{eq:endpoint-c-tau}.

For the second assertion, write
\[
 f'(\varpi_K)=\prod_{j=1}^{p-1}
 (\varpi_K-\sigma^j\varpi_K).
\]
After division by $\varpi_K^{(p-1)(t+1)}$, the residues of the factors
are the nonzero elements $-j\lambda$ of the ramification line.
Their product is $(-1)^p\lambda^{p-1}$.  Taking the norm to $F$ raises
this residue to the $p$th power and gives the last identity.
\end{proof}

Put
\[
 D=\frac{f'(\varpi_K)}{\varpi_K^{(p-1)(t+1)}}\in\cO_K^\times.
\]
Since the nontrivial ramification classes are $j\lambda$,
\begin{equation}\label{eq:endpoint-D-residue}
 \bar D=(-1)^p\lambda^{p-1},
 \qquad
 \overline{N(D)}=(-1)^p\lambda^{p(p-1)}.
\end{equation}
Using \eqref{eq:endpoint-gamma-K} and $N(\varpi_K)=\varpi_F$, we obtain
\[
 \frac{N(\gamma_K)}{\gamma_F}
   =(\gamma_F\varpi_F^t)^{p-1}N(D).
\]
Equations \eqref{eq:endpoint-c-tau} and
\eqref{eq:endpoint-D-residue} give the congruence
\begin{equation}\label{eq:endpoint-discriminant-congruence}
 \frac{N(\gamma_K)}{\gamma_F}
 \equiv(-1)^p c_\tau^{p-1}\pmod{U_F^1}.
\end{equation}
Because $m_F(\chi_F)=1$, $\chi_F$ depends only on valuation and residue,
so the congruence may be inserted in \eqref{eq:endpoint-upper-ratio}.

Finally, Lemma~\ref{lem:stationary-class-calculus}(c) shows that
$jb_\tau$ represents the stationary numerator class of $\tau^j$ on the
same layer.  Its conductor remains $t+1$ because $\tau^j$ is a
nontrivial norm character.  Choose this representative; changing
$b_\tau$ within its permitted class changes neither the residue
congruence \eqref{eq:endpoint-discriminant-congruence} nor the complete
Lamprecht factor.  Since
$m_F(\chi_F)=1\le\lfloor(t+1)/2\rfloor$, the exact stable-twist formula
gives
\[
 \Delta_F(\tau^j\chi_F,\psi_F)
 =\chi_F\!\left(\frac{c_\tau}{j}\right)
  \Delta_F(\tau^j,\psi_F).
\]
The product of the Teichmuller representatives of
$\F_p^\times$ is $-1$ for odd $p$ and is $1$ for $p=2$; hence
\[
 \prod_{j\in\F_p^\times}\frac{c_\tau}{j}
 =(-1)^p c_\tau^{p-1}.
\]
After cancelling the common factors
$\prod_{j\ne0}\Delta_F(\tau^j,\psi_F)$, the First Main Lemma is
exactly \eqref{eq:endpoint-upper-ratio} together with
\eqref{eq:endpoint-discriminant-congruence}.  Thus the endpoint
conductors $0$ and $1$ follow from the internal trace, finite-field, and
stable-twist arguments and Lemma~\ref{lem:endpoint-discriminant-congruence}.

\begin{lemma}[Characteristic-independent polar coefficient and stationary translation]
\label{lem:critical-polar-coordinate}
Let $E$ be a nonarchimedean local field, let $\psi_E$ be a
nontrivial additive character, put $p=\operatorname{char}k_E$, let
$m_E(\theta)=2d+1>1$, and choose an admissible element
$\gamma_\theta$ with
\[
 \gamma_\theta\cO_E=\p_E^{2d+1+n_E(\psi_E)}.
\]
Choose $c\in\p_E^d\setminus\p_E^{d+1}$ and a representative
$\beta_\theta\in\cO_E^\times$ of the stationary numerator class.  Fix
a nontrivial Frobenius-invariant additive character
$\psi_0:k_E\to\C^\times$, so that
$\psi_0(z^p)=\psi_0(z)$, and put
\[
 \Phi_{\theta,\beta_\theta,c}(\bar x)=
 \psi_E(\beta_\theta c\widetilde x/\gamma_\theta)
 \theta(1+c\widetilde x)^{-1}
 \qquad(\bar x\in k_E).
\]
There is a unique element
$\kappa_\theta(c)\in k_E^\times$ such that
\begin{equation}\label{eq:raw-critical-A}
 \psi_E\!\left(
   \frac{\beta_\theta c^2\widetilde z}{\gamma_\theta}
 \right)
 =\psi_0\bigl(\kappa_\theta(c)z\bigr)
 \qquad(z\in k_E).
\end{equation}
It is the polar coefficient of the complete critical function:
\begin{equation}\label{eq:critical-polar-identity}
 \frac{\Phi_{\theta,\beta_\theta,c}(x+y)}
      {\Phi_{\theta,\beta_\theta,c}(x)
       \Phi_{\theta,\beta_\theta,c}(y)}
 =\psi_0\bigl(\kappa_\theta(c)xy\bigr).
\end{equation}
If $\lambda\in k_E^\times$ and $\widetilde\lambda\in\cO_E^\times$
is any lift, then
\begin{equation}\label{eq:critical-polar-scaling}
 \kappa_\theta(\widetilde\lambda c)
 =\lambda^2\kappa_\theta(c).
\end{equation}
Precomposition with Frobenius $x\mapsto x^p$ replaces the polar
coefficient by $\kappa_\theta(c)^{1/p}$, and precomposition with
inverse Frobenius replaces it by $\kappa_\theta(c)^p$.

If $\beta_\theta' =\beta_\theta+c_0$ is another permitted
representative, define $\Delta B_\theta(c_0,c)\in k_E$ uniquely by
\begin{equation}\label{eq:critical-translation-character}
 \psi_E\!\left(
   \frac{c_0c\widetilde z}{\gamma_\theta}
 \right)
 =\psi_0\bigl(\Delta B_\theta(c_0,c)z\bigr)
 \qquad(z\in k_E).
\end{equation}
Then
\begin{equation}\label{eq:critical-translation-function}
 \Phi_{\theta,\beta_\theta',c}(x)
 =\Phi_{\theta,\beta_\theta,c}(x)
  \psi_0\bigl(\Delta B_\theta(c_0,c)x\bigr).
\end{equation}
Thus changing the representative leaves the polar coefficient fixed
and changes only the affine character; the corresponding elementary
stationary factor must be translated at the same time as in
Lemma~\ref{lem:complete-stationary-invariant}.
\end{lemma}

\begin{proof}
The coefficient
$\beta_\theta c^2/\gamma_\theta$ has valuation
$-n_E(\psi_E)-1$.  Hence the left side of
\eqref{eq:raw-critical-A} is a well-defined nontrivial additive
character of $k_E$.  The trace pairing of the finite field is perfect,
so every additive character has a unique expression
$z\mapsto\psi_0(\kappa z)$; exactness of the valuation makes
$\kappa\ne0$.  This proves existence and uniqueness of
$\kappa_\theta(c)$.

For lifts of $x,y\in k_E$, one has
\[
 \frac{(1+cx)(1+cy)}{1+c(x+y)}
 =1+\frac{c^2xy}{1+c(x+y)}.
\]
The argument on the right lies in $\p_E^{2d}\subseteq\p_E^{d+1}$,
where the stationary linearization is valid, and it differs from
$c^2xy$ by an element of $\p_E^{3d}\subseteq\p_E^{2d+1}$.  Therefore
\[
 \frac{\Phi_{\theta,\beta_\theta,c}(x+y)}
      {\Phi_{\theta,\beta_\theta,c}(x)
       \Phi_{\theta,\beta_\theta,c}(y)}
 =\psi_E\!\left(
   \frac{\beta_\theta c^2\widetilde x\widetilde y}
        {\gamma_\theta}
 \right),
\]
and \eqref{eq:critical-polar-identity} follows from
\eqref{eq:raw-critical-A}.  Replacing $c$ by
$\widetilde\lambda c$ proves
\eqref{eq:critical-polar-scaling}.  Since
$\psi_0(z^p)=\psi_0(z)$, the two Frobenius assertions follow by
reindexing the polar identity.

For another representative, $c_0$ belongs to $\p_E^d$.  Hence
$z\mapsto\psi_E(c_0c\widetilde z/\gamma_\theta)$ is a well-defined
additive character of $k_E$, possibly trivial, and so determines the
unique element in \eqref{eq:critical-translation-character}.
Substitution in the definition of $\Phi$ gives
\eqref{eq:critical-translation-function}.  The last assertion is
exactly the representative-change compatibility of
Lemma~\ref{lem:complete-stationary-invariant}.
\end{proof}


\begin{corollary}[One-layer-deeper representative changes]
\label{cor:deep-representative-change}
In the notation of Lemma~\ref{lem:critical-polar-coordinate}, suppose
that $\beta_\theta'=\beta_\theta+c_0$ with
$c_0\in\p_E^{d+1}$.  Then
\[
 \Delta B_\theta(c_0,c)=0,
 \qquad
 \Phi_{\theta,\beta_\theta',c}
 =\Phi_{\theta,\beta_\theta,c}.
\]
\end{corollary}

\begin{proof}
Since $v_E(c)=d$ and
$v_E(\gamma_\theta)=2d+1+n_E(\psi_E)$,
\[
 v_E(c_0c/\gamma_\theta)\ge -n_E(\psi_E).
\]
The additive character is therefore trivial on this quotient.  The
translation character in
\eqref{eq:critical-translation-character} is trivial, so its coefficient
is zero; equation \eqref{eq:critical-translation-function} gives the
function equality.
\end{proof}

\subsection{Odd residue characteristic: critical phases and residual coefficients}\label{sec:nonstable-odd}

We first isolate the exact coefficient extraction from Lamprecht's
critical sum and then compute the resulting stationary coefficients.

Let $k$ have odd characteristic.  By
Corollary~\ref{cor:additive-independence} and
Lemma~\ref{lem:residual-normalization}, we normalize the residual
additive character once and for all as
\[
 \psi_0(z)=\exp\!\left(\frac{2\pi i}{p}
             \Tr_{k/\F_p}(z)\right).
\]
In particular $\psi_0(z^p)=\psi_0(z)$.  For $A\ne0$ and $B\in k$ put
\begin{equation}\label{eq:QAB}
 Q_{\psi_0}(A,B)=
 |k|^{-1/2}\sum_{x\in k}\psi_0(Ax^2/2+Bx).
\end{equation}
Completing the square gives
\begin{equation}\label{eq:quadratic-phase}
 Q_{\psi_0}(A,B)
 =q_0\nu(A)\psi_0\!\left(-\frac{B^2}{2A}\right),
\end{equation}
where $\nu$ is the quadratic character and
$q_0=|k|^{-1/2}\sum_x\psi_0(x^2/2)$.

\begin{lemma}[Powers of the basic quadratic phase]
\label{lem:q0-powers}
One has
\begin{equation}\label{eq:q0-powers}
 q_0^2=\nu(-1),\qquad
 q_0^{p-1}=\nu(-1),\qquad
 q_0^p=q_0\nu(-1).
\end{equation}
\end{lemma}

\begin{proof}
The first identity is Lemma~\ref{lem:quadratic-phase-new}.  In
particular $q_0^4=1$.  If $p\equiv1\pmod4$, then $-1$ is a square in
$\F_p\subset k$, so $\nu(-1)=1$, and $p-1$ is divisible by $4$;
hence $q_0^{p-1}=1=\nu(-1)$.  If $p\equiv3\pmod4$, then
$p-1\equiv2\pmod4$, and therefore
$q_0^{p-1}=q_0^2=\nu(-1)$.  Multiplication by $q_0$ gives the last
identity.
\end{proof}

For every local character $\theta$ occurring below, we use the parity
convention
\begin{equation}\label{eq:odd-phase-parity-convention}
 \mathscr Q(\theta)=
 \begin{cases}
  Q_{\psi_0}(A_\theta,B_\theta),&m(\theta)\text{ odd},\\
  1,&m(\theta)\text{ even}.
 \end{cases}
\end{equation}
Thus a coefficient pair is required only in the odd-conductor rows; all
even-conductor residual factors are literally $1$.

For the exact transport identities below, we also use a finite function
in an even-conductor row.  If $m(\theta)=2d>1$, choose
$c\cO_E=\p_E^d$ and a representative $\beta$ of the stationary class,
and put
\begin{equation}\label{eq:even-critical-function-constant}
 \Phi_\theta(\bar x)=
 \psi(\beta c\widetilde x/\gamma)
 \theta(1+c\widetilde x)^{-1}\qquad(\bar x\in k_E).
\end{equation}
The even-conductor Lamprecht linearization is valid on $\p_E^d$ and
shows that
\[
 \Phi_\theta(\bar x)=1\qquad(\bar x\in k_E).
\]
Thus in every even row the symbol $\Phi_\theta$ denotes the literal
constant function $1$, and its normalized residue class may be taken
to be $0$.  This convention introduces no Hasse factor; it only makes
the transport identities valid without referring to an undefined odd
critical function.

For an odd-conductor character $\theta$, choose a critical coordinate
$c$ and define its coefficient pair $(A_\theta,B_\theta)$ by the exact
identity
\begin{equation}\label{eq:critical-pair-definition}
 \psi(\beta cx/\gamma)\theta(1+cx)^{-1}
 =\psi_0(A_\theta x^2/2+B_\theta x)
 \qquad(x\in k).
\end{equation}
This identity is simply the finite function occurring in Lamprecht's
formula, written in the unique quadratic-polynomial form.  Multiplying
characters adds the pairs whenever their critical layers coincide.  If
one character has smaller conductor, its second finite difference
vanishes on the deeper critical layer; its effect is then a translation
of the stationary point and has no quadratic coefficient there.  More
precisely, for $x,y\in\cO_E$ one has the exact identity
\[
 \frac{(1+cx)(1+cy)}{1+c(x+y)}
 =1+\frac{c^2xy}{1+c(x+y)}.
\]
If the smaller character has conductor at most $2v_E(c)$, it is trivial
on the unit on the right.  Hence its polar quotient on the deeper
critical line is identically one.  The quotient of its finite function
by a constant is therefore an additive character, which changes only
the affine coefficient and is exactly the stationary-point
translation retained below.

Write $b=m_F(\chi_F)-1$ and keep $t$ for the ramification break.
The next lemma fixes the residue coordinates on the last ramification
line.  This normalization is used for every odd-prime critical sum
below.

\begin{lemma}[Break-line norm, trace, and character normalization]
\label{lem:odd-break-line-normalization}
Assume that $K/F$ is wildly ramified cyclic of odd prime degree $p$,
with lower break $t$.  Choose a uniformizer $\varpi_K$, put
$\varpi_F=N_{K/F}(\varpi_K)$, choose a generator $\sigma$ of
$\Gal(K/F)$, and set
\begin{equation}\label{eq:odd-ramification-coordinate}
 \beta_0=
 \overline{\frac{\sigma(\varpi_K)-\varpi_K}
                  {\varpi_K^{t+1}}}\in k^\times.
\end{equation}
For $z\in k$ and any lift $\widetilde z\in\cO_F$, one has
\begin{equation}\label{eq:odd-critical-norm-polynomial}
 N_{K/F}(1+\varpi_K^t\widetilde z)
 \equiv
 1+\varpi_F^t\widetilde{\bigl(z^p-\beta_0^{p-1}z\bigr)}
 \pmod{\p_F^{t+1}},
\end{equation}
and consequently
\begin{equation}\label{eq:odd-break-trace-coordinate}
 \Tr_{K/F}(\varpi_K^t\widetilde z)
 \equiv
 -\beta_0^{p-1}\varpi_F^t\widetilde z
 \pmod{\p_F^{t+1}}.
\end{equation}
More generally, if $w\in\cO_K^\times$, then
\begin{equation}\label{eq:odd-break-trace-unit}
 \Tr_{K/F}(w\varpi_K^t\widetilde z)
 \equiv
 -\beta_0^{p-1}\varpi_F^t
  \widetilde{\bar w z}
 \pmod{\p_F^{t+1}}.
\end{equation}

After replacing a generator $\tau$ of $S(K/F)$ by a nonzero power,
one may and shall arrange that
\begin{equation}\label{eq:odd-tau-last-layer}
 \tau(1+\varpi_F^t\widetilde z)
   =\psi_0(\beta_0^{-p}z)\qquad(z\in k).
\end{equation}
Put
\begin{equation}\label{eq:odd-eta0-normalization}
 \eta_0=\beta_0^{-p}.
\end{equation}
If $T=t+1$ is odd, then, with the critical element
$c_\tau=\varpi_F^{t/2}$, the quadratic coefficient of the critical
function of $\tau$ is exactly $\eta_0$.
\end{lemma}

\begin{proof}
On the critical quotient $U_K^t/U_K^{t+1}$, the norm is an additive
map.  Its kernel has order $p$ by
Theorem~\ref{thm:graded-norm}(ii).  The ramification map sends
$\sigma^j$ to $j\beta_0$, so the kernel is the line
$\beta_0\F_p$.

We determine the polynomial, including its scalar.  In the exact norm
expansion of $1+\varpi_K^t\widetilde z$, the degree-$p$ endpoint is
\[
 N(\varpi_K^t\widetilde z)=\varpi_F^t\widetilde z^{\,p}.
\]
For $2\le j\le p-1$, Lemma~\ref{lem:wild-symmetric-bound} gives
\[
 v_F\bigl(s_j(\varpi_K^t\widetilde z)\bigr)
 \ge
 \left\lfloor\frac{jt+(p-1)(t+1)}p\right\rfloor
 \ge t+1,
\]
because $(j-1)t\ge1$.  Hence the induced norm polynomial is monic,
additive, and of the form $z^p+cz$.  Its roots are precisely
$\beta_0\F_p$, so
\[
 z^p+cz=\prod_{j\in\F_p}(z-j\beta_0)
       =z^p-\beta_0^{p-1}z.
\]
This proves \eqref{eq:odd-critical-norm-polynomial}; subtracting the
norm endpoint proves \eqref{eq:odd-break-trace-coordinate}.
If $w-\widetilde{\bar w}\in\p_K$, then
$\varpi_K^t(w-\widetilde{\bar w})\in\p_K^{t+1}$, and
Theorem~\ref{thm:norm-filtration}(b) gives
$\Tr(\p_K^{t+1})=\p_F^{t+1}$.  This proves
\eqref{eq:odd-break-trace-unit}.

The image of the polynomial in
\eqref{eq:odd-critical-norm-polynomial} is
\[
 \beta_0^p\ker\bigl(\Tr_{k/\F_p}:k\to\F_p\bigr).
\]
Indeed, after writing $z=\beta_0y$ it becomes
$\beta_0^p(y^p-y)$, and the image of $y\mapsto y^p-y$ is the
trace-zero hyperplane.  Write the restriction of a nontrivial norm
character to the last unit quotient as
\[
 \tau(1+\varpi_F^t\widetilde z)=\psi_0(az).
\]
It is trivial on the norm image, so the trace pairing puts
$a$ in the annihilator
$\beta_0^{-p}\F_p$.  It is nonzero because $\tau$ has conductor
$t+1$.  Thus $a=j\beta_0^{-p}$ for some $j\in\F_p^\times$, and
replacing $\tau$ by $\tau^{j^{-1}}$ proves
\eqref{eq:odd-tau-last-layer}.

If $t$ is even, the polar quotient of the $\tau$-critical function in
the coordinate $c_\tau=\varpi_F^{t/2}$ is
\[
 (x,y)\longmapsto
 \tau(1+c_\tau^2\widetilde x\widetilde y)
 =\psi_0(\beta_0^{-p}xy).
\]
Thus its quadratic coefficient is $\eta_0=\beta_0^{-p}$, as claimed.
\end{proof}

Fix from now on the generator $\tau$, the uniformizers, and the
coefficient $\eta_0$ of Lemma~\ref{lem:odd-break-line-normalization}.
After this replacement of the generator, reapply
Lemma~\ref{lem:norm-stationary-representatives} and the appropriate
high- or low-conductor stationary-class table; hence the symbols
$\alpha,\alpha_1$ and the twist representatives below refer to choices
made for this normalized generator.  Whenever $\chi_F$ has odd
conductor, write
$(\eta,\eta\gamma)$ for its coefficient pair in the critical
coordinate specified in Lemma~\ref{lem:residual-parameter-transport}; write
$(\eta_0,\eta_0\gamma_0)$ for the corresponding pair of $\tau$ when
$T$ is odd.  When $T$ is even, $\eta_0$ is the normalized last-layer
coefficient \eqref{eq:odd-eta0-normalization}; no $\tau$-critical sum
occurs.


\begin{remark}[Intrinsic break-line coefficient]
\label{rem:intrinsic-break-line-coefficient}
The coefficient $\eta_0$ in the break-line formulas is retained as the
coefficient determined by the fixed Galois generator and the chosen
uniformizer.  We do not replace the generator merely to force a cosmetic
normal form such as $\eta_0=\beta_0^{-p}$.  All later annihilator,
Frobenius-preimage, and coefficient identities are proved for this
intrinsic value.  This keeps the source of every critical row visible.
\end{remark}

\begin{lemma}[Critical polynomial, translations, and changes of coordinate]
\label{lem:critical-polynomial-coordinate}
Let $E$ be a local field of odd residue characteristic, let
$m_E(\theta)=2d+1$, and choose $c\in\p_E^d\setminus\p_E^{d+1}$.
For a chosen representative $\beta_\theta$ of the stationary
numerator class, put
\[
 \Phi_\theta(\bar x)=
 \psi_E(\beta_\theta c\widetilde x/\gamma_\theta)
 \theta(1+c\widetilde x)^{-1}\qquad(\bar x\in k_E).
\]
Let
\[
 A_\theta=\kappa_\theta(c)\in k_E^\times
\]
be the polar coefficient defined by
Lemma~\ref{lem:critical-polar-coordinate}.  Then there is a unique
$B_\theta\in k_E$ such that
\[
 \Phi_\theta(x)=
 \psi_0(A_\theta x^2/2+B_\theta x).
\]
Replacing $c$ by $\widetilde\lambda c$, where
$\lambda\in k_E^\times$ and $\widetilde\lambda$ is a unit lift,
replaces $(A_\theta,B_\theta)$ by
$(\lambda^2A_\theta,\lambda B_\theta)$.  Precomposing the finite
function with Frobenius $x\mapsto x^p$ replaces the pair by
$(A_\theta^{1/p},B_\theta^{1/p})$, whereas precomposing with inverse
Frobenius $x\mapsto x^{1/p}$ replaces it by
$(A_\theta^p,B_\theta^p)$.  These operations do not change the
normalized finite sum, because they merely reindex it.

If a permitted representative is replaced by
$\beta_\theta+c_0$ and the elementary stationary factor is translated
at the same time, then $A_\theta$ is unchanged and
$B_\theta$ is replaced by
\[
 B_\theta+\Delta B_\theta(c_0,c),
\]
where $\Delta B_\theta(c_0,c)$ is defined by
\eqref{eq:critical-translation-character}.  Thus a linear coefficient
may never be discarded without simultaneously retaining the
corresponding elementary translation factor.
\end{lemma}

\begin{proof}
By Lemma~\ref{lem:critical-polar-coordinate}, the polar quotient of
$\Phi_\theta$ is $\psi_0(A_\theta xy)$.  Since the residue
characteristic is odd, the function
$x\mapsto\psi_0(A_\theta x^2/2)$ has the same polar quotient.
Their quotient is therefore an additive character of $k_E$, and the
perfect trace pairing gives a unique $B_\theta\in k_E$ for which it is
$\psi_0(B_\theta x)$.  This proves existence and uniqueness.

The coordinate-scaling rule follows from
$\Phi_{\theta,\beta_\theta,\widetilde\lambda c}(x)
 =\Phi_{\theta,\beta_\theta,c}(\lambda x)$ and
\eqref{eq:critical-polar-scaling}.  For Frobenius, write
$A_\theta=A_0^p$ and $B_\theta=B_0^p$.  Since
$\psi_0(z^p)=\psi_0(z)$,
\[
 \psi_0(A_\theta x^{2p}/2+B_\theta x^p)
 =\psi_0(A_0x^2/2+B_0x),
\]
which proves the stated $p$th-root direction; the inverse direction is
equivalent.  Finally,
\eqref{eq:critical-translation-function} shows that changing the
representative multiplies the critical function by
$\psi_0(\Delta B_\theta(c_0,c)x)$, so it changes only the displayed
affine coefficient.  Lemma~\ref{lem:complete-stationary-invariant}
provides the matching elementary translation.
\end{proof}

\begin{lemma}[Higher symmetric terms on the upper critical line]
\label{lem:odd-upper-higher-vanishing}
Assume that $K/F$ is wildly ramified cyclic of odd prime degree $p$,
with break $t$, and put $T=t+1$.  Suppose
$m_F(\chi_F)=2d+1>1$, set $b=2d$ and $a=t-b$, and suppose that the
upper conductor is odd, say $m_K(\chi_K)=2d_K+1$, with
\[
 2d_K=
 \begin{cases}
  b,&b<t,\\
  t,&b=t,\\
  t+p(b-t),&b>t.
 \end{cases}
\]
Let $u\in K^\times$ satisfy $v_K(u)=a$, put $n=N_{K/F}(u)$, and let
$x\in\p_K^{d_K}$.  Then, for every $3\le j\le p-1$,
\begin{equation}\label{eq:odd-upper-higher-terms-vanish}
 s_j(ux)\in\p_F^T,
 \qquad
 n\,s_j(x)\in\p_F^T.
\end{equation}
Moreover
\begin{equation}\label{eq:odd-upper-H-quadratic-reduction}
 \sum_{j=2}^{p}\bigl(s_j(ux)-n\,s_j(x)\bigr)
 \equiv s_2(ux)-n\,s_2(x)\pmod{\p_F^T}.
\end{equation}
Thus the surviving correction on the last ramification quotient is
homogeneous quadratic and has no affine part.
\end{lemma}

\begin{proof}
Since $v_F(n)=v_K(u)=a$, Lemma~\ref{lem:wild-symmetric-bound} gives
\begin{align*}
 v_F(s_j(ux))
 &\ge
 \left\lfloor
  \frac{j(d_K+a)+(p-1)T}{p}\right\rfloor,\\
 v_F(n\,s_j(x))
 &\ge a+
 \left\lfloor
  \frac{jd_K+(p-1)T}{p}\right\rfloor.
\end{align*}
It is therefore enough to prove
\begin{equation}\label{eq:odd-upper-two-numerical-inequalities}
 j(d_K+a)\ge T,
 \qquad
 jd_K+pa\ge T.
\end{equation}
If $p=3$, the range $3\le j\le p-1$ is empty and
$s_3(ux)=N(u)N(x)=n\,s_3(x)$, so both assertions of the lemma follow
immediately.  We may therefore suppose $p\ge5$.  The formulas below
show in particular that $d_K+a>0$; hence both left sides increase with
$j$, and it suffices to take $j=3$.
For $b<t$, put $a=t-b>0$, so $d_K=(t-a)/2$; for $b>t$, put
$q=b-t=-a>0$, so $d_K=(t+pq)/2$.  Direct substitution gives
\[
\begin{array}{c|c|c}
 &3(d_K+a)-T&3d_K+pa-T\\ \hline
 b<t&\dfrac{t+3a-2}{2}&\dfrac{t+(2p-3)a-2}{2}\\[2mm]
 b=t&\dfrac{t-2}{2}&\dfrac{t-2}{2}\\[2mm]
 b>t&\dfrac{t+3(p-2)q-2}{2}&\dfrac{t+pq-2}{2}.
\end{array}
\]
Every entry is nonnegative.  In the middle row, $b=t=2d$ and
$d\ge1$, so $t\ge2$; in the other rows use $t\ge1$, $a\ge1$ or
$q\ge1$, and $p\ge3$.  This proves
\eqref{eq:odd-upper-two-numerical-inequalities} and hence
\eqref{eq:odd-upper-higher-terms-vanish}.

Finally,
\[
 s_p(ux)=N(ux)=N(u)N(x)=n\,s_p(x),
\]
so the degree-$p$ terms cancel exactly.  Removing them and the terms
of degrees $3,\ldots,p-1$ proves
\eqref{eq:odd-upper-H-quadratic-reduction}.  The remaining term is
homogeneous of degree two in $x$.
\end{proof}

\begin{lemma}[Odd-prime critical-coordinate transport]
\label{lem:residual-parameter-transport}
Put
\[
 (u,n)=
 \begin{cases}
 \left(\eps_1\beta_1/\alpha_1,\ \eps\beta/\alpha\right),&b\le t,\\[1mm]
 \left(\beta_1/\alpha_1,\ \beta/\alpha\right),&b>t.
 \end{cases}
\]
Then $n=N(u)$.  Put $a=t-b$, write
\[
 u=\varpi_K^a u_0,\qquad n=\varpi_F^a n_0,
 \qquad u_0,n_0\in\cO^\times,
\]
and set
\begin{equation}\label{eq:odd-zeta-definition}
 \zeta=\bar u_0,\qquad d_0=\zeta^p.
\end{equation}
Then $\bar n_0=d_0$.  When $m_F(\chi_F)=2d+1$ is odd, the lower
critical coordinate $c_\chi=\varpi_F^d$ gives
\begin{equation}\label{eq:eta-normalization}
 \eta=\eta_0d_0.
\end{equation}
The upper conductor has the same parity as $m_F(\chi_F)$.  When this
common parity is odd, the phase-preserving upper critical coordinate
constructed in the proof gives the complete coefficient pair
\begin{equation}\label{eq:odd-transport-direct-pair}
 (A_K,B_K)=
 \begin{cases}
 (\eta,\eta\gamma),&b<t,\\[1mm]
 (-\eta_0(d_0^p-d_0),\ \eta_0d_0(\gamma-\gamma_0)),&b=t,\\[1mm]
 (-d_0^{p-1}\eta,0),&b>t.
 \end{cases}
\end{equation}
When the common parity is even, the corresponding residual factors are
$1$ by \eqref{eq:odd-phase-parity-convention}.  In the odd rows,
\eqref{eq:odd-transport-direct-pair} is an equality of complete
critical functions, including the affine term and the translation of
the stationary origin.
\end{lemma}

\begin{proof}
Use the uniformizers fixed in
Lemma~\ref{lem:odd-break-line-normalization}.  In every range of the
statement,
\[
 v_K(u)=v_F(n)=a=t-b.
\]
Thus the decompositions in the statement satisfy $n_0=N(u_0)$.  Since
the extension is totally ramified, the residue norm is Frobenius, and
therefore
\begin{equation}\label{eq:odd-transport-leading-norm}
 \bar n_0=\bar u_0^{\,p}=\zeta^p=d_0.
\end{equation}
There is no additional residue scalar when $n$ is viewed as an element
of $K$.  Indeed, if $\sigma$ generates $\Gal(K/F)$, then
\[
 \frac{\varpi_F}{\varpi_K^p}
 =\prod_{i=0}^{p-1}\frac{\sigma^i(\varpi_K)}{\varpi_K}
 \equiv1\pmod{\p_K},
\]
because the extension is wildly ramified.  Hence, after division by
$\varpi_K^{pa}$, the initial residue of
$n=\varpi_F^a n_0$ is $d_0$ for every integer $a$, including $a<0$.

Put
\begin{equation}\label{eq:odd-A-definition}
 \mathcal A=
 \begin{cases}
  \alpha/\delta,&b\le t,\\[1mm]
  \alpha/\gamma_F,&b>t.
 \end{cases}
\end{equation}
Here $\delta$ in the first row is the low-conductor admissible element
of \eqref{eq:gammas-low}, not the residue $\zeta$ of
\eqref{eq:odd-zeta-definition}.  Corollary
\ref{cor:stationary-class-under-norm}, together with
Propositions~\ref{prop:high-parameter-table} and
\ref{prop:low-parameter-table}, identifies the three stationary
classes.  Choose the representatives displayed in those tables and,
whenever a conductor is odd, translate the corresponding critical
function as in Lemma~\ref{lem:complete-stationary-invariant}.  For these
compatible choices the following are exact equalities in the fields,
not merely congruences of classes:
\begin{equation}\label{eq:odd-three-stationary-ratios}
 \frac{\beta(\tau)}{\gamma_\tau}=\mathcal A,
 \qquad
 \frac{\beta(\chi_F)}{\gamma_{\chi}}=\mathcal A n,
 \qquad
 \frac{\beta_K}{\gamma_K}=\mathcal A(n-u).
\end{equation}
In the low range the selected ratios are respectively
$\alpha/\delta$, $\varepsilon\beta/\delta$, and
$\alpha(n-u)/\delta$, whereas in the high range they are
$\alpha/\gamma_F$, $\beta/\gamma_F$, and
$\alpha(n-u)/\gamma_F$.  The normalization of the generator $\tau$
in Lemma~\ref{lem:odd-break-line-normalization} is consequently the
identity
\begin{equation}\label{eq:odd-tau-last-layer-normalization}
 \psi_F(\mathcal A\varpi_F^t\widetilde z)
 =\psi_0(\eta_0z)\qquad(z\in k).
\end{equation}

Suppose first that $m_F(\chi_F)=2d+1$.  Take
$c_\chi=\varpi_F^d$.  Since $a+2d=t$, the polar quotient of the
$\chi_F$ critical function is
\begin{align*}
 (X,Y)&\longmapsto
 \psi_F(\mathcal A n c_\chi^2\widetilde X\widetilde Y)\\
 &=\psi_F(\mathcal A\varpi_F^t n_0
                 \widetilde X\widetilde Y)
 =\psi_0(\eta_0\bar n_0XY).
\end{align*}
Together with \eqref{eq:odd-transport-leading-norm}, this proves
\eqref{eq:eta-normalization} in the stated coordinate.

We now calculate the upper quadratic coefficient without suppressing
the trace scalar.  The parity of $m_K(\chi_K)$ is the parity of
$m_F(\chi_F)$: in the low range the conductors are equal, while in the
high range
\[
 m_K(\chi_K)-1=t+p(b-t)\equiv b\pmod2.
\]
Assume this parity is odd and write $m_K(\chi_K)=2d_K+1$.  Put
$r=v_K(n-u)$.  The exact conductor of $\chi_K$ implies $r=0$ at the
boundary; away from the boundary the two terms have distinct
valuations.  Hence
\begin{equation}\label{eq:odd-upper-depth-table}
 \begin{array}{c|c|c}
  &r&2d_K\\ \hline
  b<t&t-b&b\\
  b=t&0&t\\
  b>t&p(t-b)&t+p(b-t).
 \end{array}
\end{equation}
In every row,
\begin{equation}\label{eq:odd-upper-break-depth}
 r+2d_K=t.
\end{equation}
Take $c_K=\varpi_K^{d_K}$ and put
$w_0=\varpi_K^{-r}(n-u)\in\cO_K^\times$.  In the natural upper
residue coordinate $Z$, define
\[
 \Phi_K^{\mathrm{nat}}(Z)=
 \psi_K\!\left(\frac{\beta_Kc_K\widetilde Z}{\gamma_K}\right)
 \chi_K(1+c_K\widetilde Z)^{-1}.
\]
For $Z,W\in k$, Lemma~\ref{lem:critical-polynomial-coordinate} and
\eqref{eq:odd-three-stationary-ratios} give
\begin{align*}
 \frac{\Phi_K^{\mathrm{nat}}(Z+W)}
      {\Phi_K^{\mathrm{nat}}(Z)\Phi_K^{\mathrm{nat}}(W)}
 &=\psi_F\!\left(
   \mathcal A\Tr_{K/F}\bigl((n-u)c_K^2
                  \widetilde Z\widetilde W\bigr)\right).
\end{align*}
By \eqref{eq:odd-upper-break-depth}, the argument of the trace is
$w_0\varpi_K^t\widetilde Z\widetilde W$.  Formula
\eqref{eq:odd-break-trace-unit} and the fact that
$x\mapsto\psi_F(\mathcal A x)$ is trivial on $\p_F^{t+1}$ therefore
give
\begin{align}
 \frac{\Phi_K^{\mathrm{nat}}(Z+W)}
      {\Phi_K^{\mathrm{nat}}(Z)\Phi_K^{\mathrm{nat}}(W)}
 &=\psi_0\!\left(
   -\eta_0\beta_0^{p-1}\bar w_0ZW\right)\notag\\
 &=\psi_0\!\left(
   \beta_0^{-1}\operatorname{in}(u-n)ZW\right),
 \label{eq:odd-upper-natural-polar}
\end{align}
where
$\operatorname{in}(y)=\overline{\varpi_K^{-v_K(y)}y}$.
The last equality uses $\eta_0=\beta_0^{-p}$ and
$\operatorname{in}(u-n)=-\bar w_0$.

The displayed upper coordinate is the inverse-Frobenius reindexing of
the natural one:
\begin{equation}\label{eq:odd-upper-coordinate-reindexing}
 \Phi_{\chi_K}(X)=\Phi_K^{\mathrm{nat}}(X^{1/p}).
\end{equation}
It is phase preserving because Frobenius permutes $k$.  By
Lemma~\ref{lem:critical-polynomial-coordinate}, the quadratic
coefficient in this coordinate is the $p$th power of the coefficient
in \eqref{eq:odd-upper-natural-polar}.  Thus
\begin{equation}\label{eq:upper-A-leading}
 \frac{A_K}{\eta_0}=\operatorname{in}(u-n)^p.
\end{equation}
Equivalently, for $X,Y\in k$,
\begin{equation}\label{eq:odd-upper-displayed-polar}
 \frac{\Phi_{\chi_K}(X+Y)}
      {\Phi_{\chi_K}(X)\Phi_{\chi_K}(Y)}
 =\psi_0\!\left(
   \eta_0\operatorname{in}(u-n)^pXY\right).
\end{equation}
This proves the exact scalar in the upper polar form; no square-class
choice or unrecorded uniformizer factor remains.

The initial term of $u-n$ is determined by the two valuations.  If
$b<t$, then $a>0$ and $v_K(u)=a<pa=v_K(n)$, so it is $\zeta$.  If
$b=t$, both terms are units and it is $\zeta-d_0$.  If $b>t$, then
$a<0$ and $v_K(n)=pa<a=v_K(u)$, so it is $-d_0$.  Consequently
\[
 A_K=\eta_0
 \begin{cases}
  \zeta^p,&b<t,\\
  (\zeta-d_0)^p,&b=t,\\
  (-d_0)^p,&b>t,
 \end{cases}
 =
 \begin{cases}
  \eta,&b<t,\\
  -\eta_0(d_0^p-d_0),&b=t,\\
  -d_0^{p-1}\eta,&b>t.
 \end{cases}
\]
Here $\zeta^p=d_0$, $\eta=\eta_0d_0$, and $p$ is odd.

It remains to determine the affine coefficient.  We do this from an
exact identity of complete critical functions, retaining the
stationary ratios \eqref{eq:odd-three-stationary-ratios}.

Let
\[
 P(y)=N(1+y)-1,
 \qquad
 H(y)=P(y)-\Tr(y)=\sum_{j=2}^{p}s_j(y).
\]
For $X\in k$, let $Z=X^{1/p}$, choose a lift $\widetilde Z$, and put
$x=c_K\widetilde Z$.  Thus $X$ is precisely the displayed upper
coordinate of \eqref{eq:odd-upper-coordinate-reindexing}.  Let
$r_\chi(X)$ be the normalized class of $P(x)$ on the $\chi_F$ critical
line and let $r_\tau(X)$ be the normalized class of $P(ux)$ on the
$\tau$ critical line.  A class is defined to be zero if the
corresponding conductor is even or if the element lies one step beyond
the critical line.  Then
\begin{equation}\label{eq:odd-exact-critical-transport}
 \Phi_{\chi_K}(X)
 =\Phi_{\chi_F}(r_\chi(X))\,
  \Phi_\tau(r_\tau(X))^{-1}\,
  \psi_F\!\left(\mathcal A\{H(ux)-nH(x)\}\right).
\end{equation}
To verify this identity, first note that
$1+P(ux)=N(1+ux)$ and hence
$\tau(1+P(ux))=1$.  Therefore the right-hand side of
\eqref{eq:odd-exact-critical-transport} is
\begin{align*}
 &\psi_F(\mathcal A nP(x))\chi_F(1+P(x))^{-1}
  \psi_F(-\mathcal A P(ux))
  \psi_F\!\left(\mathcal A\{H(ux)-nH(x)\}\right)\\
 &\qquad=
 \psi_F\!\left(\mathcal A\{n\Tr(x)-\Tr(ux)\}\right)
 \chi_F(N(1+x))^{-1}\\
 &\qquad=
 \psi_K\!\left(\frac{\beta_Kx}{\gamma_K}\right)
 \chi_K(1+x)^{-1}
 =\Phi_{\chi_K}(X),
\end{align*}
where the second equality uses $P=\Tr+H$ and the third uses
\eqref{eq:odd-three-stationary-ratios}.

We now compute the two classes in
\eqref{eq:odd-exact-critical-transport}.  Only odd conductors need be
considered.  If $b<t$, write $m=2d+1$, so $b=2d<t$ and the upper
conductor is again $m$.  Here $c_K=\varpi_K^d$ and
$c_\chi=\varpi_F^d$.  For
$x=c_K\widetilde{X^{1/p}}$, the norm term is
\[
 N(x)=c_\chi\widetilde X
 \pmod{\p_F^{d+1}},
\]
whereas Lemma~\ref{lem:wild-symmetric-bound} gives
$v_F(s_j(x))\ge d+1$ for $1\le j<p$.  Thus $r_\chi(X)=X$ exactly in
the coordinate \eqref{eq:odd-upper-coordinate-reindexing}.  If the
conductor $T=t+1$ of $\tau$ is odd, put $d_\tau=t/2$.  Since
$v_K(u)=t-b$,
\[
 v_K(ux)=d+t-b=t-d\ge d_\tau+1,
\]
and the same symmetric-function estimate puts every term of $P(ux)$
in $\p_F^{d_\tau+1}$.  If $T$ is even there is no $\tau$ critical
line.  In either case $r_\tau(X)=0$.

If $b=t$, then necessarily $m=T=2d+1$ in the odd row and
$v_K(u)=0$.  In this case
\[
 c_K=\varpi_K^d,
 \qquad c_\chi=c_\tau=\varpi_F^d.
\]
Modulo the next critical ideal,
\[
 P(x)\equiv N(x)=c_\chi\widetilde X,
 \qquad
 P(ux)\equiv N(u)N(x)=c_\tau\widetilde{d_0X}.
\]
Hence $r_\chi(X)=X$ and $r_\tau(X)=d_0X$ with no further coordinate
choice.

Finally suppose $b>t$ and write again $m=2d+1$.  The upper critical
depth is
\[
 d_K=\frac{t+p(b-t)}2.
\]
The norm term of $P(x)$ lies beyond the $\chi_F$ critical line because
\[
 d_K-d=\frac{(p-1)(b-t)}2\ge1.
\]
For the remaining symmetric terms, the smallest lower bound occurs at
$j=1$, and
\[
 2\{d_K+(p-1)T-p(d+1)\}=(p-1)t-2\ge0.
\]
Hence Lemma~\ref{lem:wild-symmetric-bound} gives
$P(x)\in\p_F^{d+1}$, so $r_\chi(X)=0$.  If $T$ is odd, write
$t=2d_\tau$.  Then
\[
 v_K(ux)=d_K+t-b
 =d_\tau+(p-2)(d-d_\tau)\ge d_\tau+1,
\]
and again all terms of $P(ux)$ lie beyond the $\tau$ critical line.
Thus $r_\tau(X)=0$; the assertion is automatic when $T$ is even.
We have proved
\begin{equation}\label{eq:odd-critical-class-table}
 \begin{array}{c|c|c}
  &r_\chi(X)&r_\tau(X)\\ \hline
  b<t&X&0\\
  b=t&X&d_0X\\
  b>t&0&0.
 \end{array}
\end{equation}

The polynomial $H(ux)-nH(x)$ has no terms of degree zero or one.
Lemma~\ref{lem:odd-upper-higher-vanishing} shows that, modulo
$\p_F^T$, its degree-$p$ terms cancel exactly and every term of degree
$3,\ldots,p-1$ vanishes.  Since
$z\mapsto\psi_F(\mathcal A z)$ is trivial on $\p_F^T$, the correction
factor in \eqref{eq:odd-exact-critical-transport} is therefore
homogeneous quadratic and contributes no affine character.  The total
upper quadratic coefficient has already been computed independently in
\eqref{eq:upper-A-leading}; no identification of the correction factor
alone with that coefficient is required.  Consequently the affine
coefficient of \eqref{eq:odd-exact-critical-transport} is obtained only
from the two lower critical functions in
\eqref{eq:odd-critical-class-table}.  Their affine parts give
\begin{equation}\label{eq:odd-transport-affine-functions}
 L_{\chi_K}(X)=
 \begin{cases}
  \psi_0(\eta\gamma X),&b<t,\\[1mm]
  \psi_0\!\left((\eta\gamma-d_0\eta_0\gamma_0)X\right)
   =\psi_0\!\left(\eta_0d_0(\gamma-\gamma_0)X\right),&b=t,\\[1mm]
  1,&b>t.
 \end{cases}
\end{equation}
Together with the quadratic coefficients already computed, this proves
\eqref{eq:odd-transport-direct-pair}.  Since
\eqref{eq:odd-exact-critical-transport} is an identity of the complete
finite functions, the translation of the elementary stationary value
has been retained throughout.
\end{proof}


\begin{lemma}[Successor-half-depth Teichmuller congruence]
\label{lem:odd-teichmuller-successor}
Let $K/F$ be wildly ramified cyclic of odd prime degree $p$, with
positive lower break $t$, and put $T=t+1$.  Suppose that $T>1$ is odd.
For $j\in\F_p$, let $[j]\in\cO_F$ be its Teichmuller lift and let the
same symbol $j$ on the right denote its natural prime-field lift.  Then
\begin{equation}\label{eq:odd-teichmuller-successor}
 [j]-j\in\p_F^{(T+1)/2}.
\end{equation}
\end{lemma}

\begin{proof}
In equal characteristic the two lifts agree.  In mixed characteristic,
$[j]-j$ is divisible by $p$ in the maximal unramified subfield of $F$.
Since $p=\Tr_{K/F}(1)$ and the different exponent is $(p-1)T$, the
trace-ideal formula gives
\[
 v_F(p)\ge \left\lfloor\frac{(p-1)T}{p}\right\rfloor.
\]
Write $T=2r+1$.  Since $T>1$, one has $r\ge1$, and for $p\ge3$,
\[
 p(r+1)\le(p-1)(2r+1).
\]
Hence
\[
 \left\lfloor\frac{(p-1)T}{p}\right\rfloor
 \ge r+1=\frac{T+1}{2},
\]
which proves \eqref{eq:odd-teichmuller-successor}.
\end{proof}

\begin{proposition}[Vanishing of the actual wild-odd affine translations]
\label{prop:wild-odd-translation-vanishing}
Use the actual high- and low-conductor representatives from
Propositions~\ref{prop:high-parameter-table} and
\ref{prop:low-parameter-table}, and compare them with the literal
natural-power representatives used in the finite coefficient table.
Whenever the corresponding critical conductor is odd, let $d_*$ be its
critical half-depth and let
\[
 c_{0,j}^{H},\qquad c_{0,j}^{R},\qquad c_{0,j}^{T}
\]
denote the high norm-row, low norm-row, and low boundary-twist
representative differences.  Then
\begin{equation}\label{eq:wild-odd-translation-differences-deep}
 c_{0,j}^{H},\ c_{0,j}^{R},\ c_{0,j}^{T}
 \in\p_F^{d_*+1},
\end{equation}
and hence all corresponding affine translation coefficients vanish:
\begin{equation}\label{eq:wild-odd-translation-coefficients-zero}
 \Delta_j^H=\Delta_j^R=\Delta_j^T=0.
\end{equation}
Thus the actual source rows and the normalized rows used in the
coefficient table have identical complete critical functions, not merely
the same polar coefficients.
\end{proposition}

\begin{proof}
In every row where this comparison occurs, the odd norm-character
conductor is $T=2d_*+1$.  The actual scalar is the Teichmuller
representative $[j]$, whereas the literal power row uses the natural
scalar $j$.  Lemma~\ref{lem:odd-teichmuller-successor} gives
\[
 [j]-j\in\p_F^{d_*+1}.
\]
Multiplication by the selected stationary source coefficient preserves
this extra layer and proves the high inclusion in
\eqref{eq:wild-odd-translation-differences-deep}.  The low norm and
boundary-twist rows use the same Teichmuller-to-natural comparison after
the common low stationary scalar is inserted, so the identical argument
proves the other two inclusions.  Corollary
\ref{cor:deep-representative-change} now gives
\eqref{eq:wild-odd-translation-coefficients-zero} and the equality of
complete critical functions.
\end{proof}

\begin{proposition}[Complete odd-prime coefficient table]
\label{prop:odd-coefficient-table}
The following coefficient pairs apply whenever the corresponding
conductor is odd; for even conductor the phase is $1$ by
\eqref{eq:odd-phase-parity-convention}.  If $b<t$, then
\begin{equation}\label{eq:lower-coeff-table-below}
 (A_j,B_j)=(j\eta_0,j\eta_0\gamma_0)
 \quad(j\in\F_p^\times),
 \qquad
 (A_0,B_0)=(\eta,\eta\gamma).
\end{equation}
If $b=t$, then for every $j\in\F_p$,
\begin{equation}\label{eq:lower-coeff-table-boundary}
 (A_j,B_j)=
 (\eta_0(j+d_0),\eta_0(j\gamma_0+d_0\gamma)).
\end{equation}
If $b>t$, then for every $j\in\F_p$,
\begin{equation}\label{eq:lower-coeff-table-above}
 (A_j,B_j)=(\eta,\eta\gamma).
\end{equation}
The upper character $\chi_K$ has the pair
\begin{equation}\label{eq:upper-coeff-table}
 (A_K,B_K)=
 \begin{cases}
 (\eta,\eta\gamma),&b<t,\\[1mm]
 (-\eta_0(d_0^p-d_0),
   \eta_0d_0(\gamma-\gamma_0)),&b=t,\\[1mm]
 (-d_0^{p-1}\eta,0),&b>t.
 \end{cases}
\end{equation}
\end{proposition}

\begin{proof}
Proposition~\ref{prop:wild-odd-translation-vanishing} first identifies
the actual source rows with the normalized rows: every affine
representative-change coefficient is zero.  We may therefore compute
with the displayed normalized representatives without dropping a
translation term.

We next prove equality of the complete lower critical functions, not
only equality of their polar quotients.  On the dominant critical
layer, write all linearized restrictions with one additive denominator
$\Gamma$.  Let $\beta_\tau$ and $\beta_\chi$ denote coefficients for
the restrictions of $\tau$ and $\chi_F$ relative to $\Gamma$; if a
character is already trivial on that layer, take the corresponding
coefficient to be $0$.  By
Lemma~\ref{lem:stationary-class-calculus}(c), the coefficient
$j\beta_\tau+\beta_\chi$ represents the restriction of
$\tau^j\chi_F$ there.

Suppose $b<t$ and $j\ne0$.  A coefficient pair is needed only when
$T$ is odd.  Put $d_\tau=(T-1)/2$ and choose
$c_\tau\in\p_F^{d_\tau}\setminus\p_F^{d_\tau+1}$.  Since
$m_F(\chi_F)\le T-1$,
\[
 v_F(c_\tau)=\frac{T-1}{2}
 \ge\left\lceil\frac{m_F(\chi_F)}2\right\rceil.
\]
Thus the stationary linearization of $\chi_F$ is valid at
$c_\tau\widetilde x$.  With the sum representative just chosen,
\begin{align*}
 \frac{\Phi_{\tau^j\chi_F}(x)}{\Phi_{\tau^j}(x)}
 &=\psi_F\!\left(
   \frac{\beta_\chi c_\tau\widetilde x}{\Gamma}
  \right)
   \chi_F(1+c_\tau\widetilde x)^{-1}\\
 &=1.
\end{align*}
Hence the two complete critical functions agree.  Since
$\Phi_{\tau^j}=\Phi_\tau^j$, its pair is
$(j\eta_0,j\eta_0\gamma_0)$.  The case $j=0$ is $\chi_F$ itself and
has pair $(\eta,\eta\gamma)$.

Suppose $b>t$.  A coefficient pair is needed only when
$m_F(\chi_F)$ is odd.  Put
$d_\chi=(m_F(\chi_F)-1)/2$ and choose
$c_\chi\in\p_F^{d_\chi}\setminus\p_F^{d_\chi+1}$.  Since
$m_F(\chi_F)>T$,
\[
 v_F(c_\chi)=\frac{m_F(\chi_F)-1}{2}
 \ge\left\lceil\frac T2\right\rceil.
\]
The stationary linearization of every $\tau^j$ is therefore valid at
$c_\chi\widetilde x$, and for every $j\in\F_p$,
\begin{align*}
 \frac{\Phi_{\tau^j\chi_F}(x)}{\Phi_{\chi_F}(x)}
 &=\psi_F\!\left(
   \frac{j\beta_\tau c_\chi\widetilde x}{\Gamma}
  \right)
   \tau^j(1+c_\chi\widetilde x)^{-1}\\
 &=1.
\end{align*}
Thus every complete critical function has the pair
$(\eta,\eta\gamma)$.

At $b=t$, the two characters have the common conductor $T$.  Lemma
\ref{lem:stationary-class-calculus}(c) identifies the stationary class
of $\tau^j\chi_F$ with the sum of the two classes, and
Lemma~\ref{lem:minimal-orbit-stationary} shows that its leading class
does not cancel.  We may therefore choose the sum representative at the
same conductor and translate the critical functions consistently.  The
complete critical function of the product is the pointwise product of
the two complete critical functions, so the coefficient pairs add:
\[
 j(\eta_0,\eta_0\gamma_0)
 +(\eta_0d_0,\eta_0d_0\gamma).
\]
This proves the entire lower table, including its affine coefficients,
without reusing a cancelled representative at a smaller conductor.
The upper table is exactly \eqref{eq:odd-transport-direct-pair} of
Lemma~\ref{lem:residual-parameter-transport}.
\end{proof}


\begin{proposition}[Odd-prime finite phase cancellation]
\label{prop:odd-phase-cancel}
Let $R_j$ be the critical phase of $\tau^j$, $T_j$ that of
$\tau^j\chi_F$, and $U$ the upper critical phase.  Then
\begin{equation}\label{eq:Rodd}
 \mathcal R_{\mathrm{odd}}
 :=U\frac{\prod_{j\ne0}R_j}{\prod_jT_j}=1.
\end{equation}
\end{proposition}

\begin{proof}
All phases in this proof use the parity convention
\eqref{eq:odd-phase-parity-convention}.

If $b<t$, \eqref{eq:lower-coeff-table-below} gives $T_j=R_j$ for
$j\ne0$, while the $j=0$ pair and the upper pair agree, so $T_0=U$.

Suppose $b>t$.  If $m_F(\chi_F)$ is even, then every $T_j$ and $U$ is
$1$.  If it is odd, all $T_j$ are the same phase
$T=Q_{\psi_0}(\eta,\eta\gamma)$.  In the explicit formula for
$Q_{\psi_0}$, the additive factor has $p$th power $1$, because its
value lies in the group of $p$th roots of unity.  Hence Lemma~\ref{lem:q0-powers} gives
\[
 T^p=q_0^p\nu(\eta)^p
 =q_0\nu(-1)\nu(\eta)
 =q_0\nu(-\eta)=U,
\]
because $d_0^{p-1}$ is a square and the upper affine coefficient is
zero.  Independently, if the conductor $T=t+1$ of $\tau$ is even then
all $R_j$ are $1$.  If it is odd, then
\begin{align*}
 \prod_{j\ne0}R_j
 &=q_0^{p-1}\nu(\eta_0)^{p-1}
   \nu\!\left(\prod_{j\ne0}j\right)
   \psi_0\!\left(-\frac{\eta_0\gamma_0^2}{2}
                         \sum_{j\ne0}j\right)\\
 &=q_0^{p-1}\nu(-1)=1.
\end{align*}
Here $\sum_{j\ne0}j=0$, $\prod_{j\ne0}j=-1$, and
$q_0^{p-1}=\nu(-1)$ by Lemma~\ref{lem:q0-powers}.  Hence
$U\prod R_j=\prod T_j$ in every parity
pattern above the break.

Finally suppose $b=t$.  All four kinds of character have conductor
$T$.  If $T$ is even, every residual phase is $1$.  If $T$ is odd, the
actual boundary coefficient pairs carry one common nonzero scalar.
Proposition~\ref{prop:affine-pencil-homogeneous} shows that removing
this scalar does not change the complete affine-pencil quotient, while
Proposition~\ref{prop:wild-odd-translation-vanishing} shows that no
affine representative-change term remains.  We may therefore use the
normalized pairs of Proposition~\ref{prop:odd-coefficient-table}.
Put $d_0=\zeta^p$ as in \eqref{eq:odd-zeta-definition}.  Here
$d_0=\bar n$, and Lemma~\ref{lem:minimal-orbit-stationary} together
with \eqref{eq:stationary-leading-cancellation} gives
$j+d_0\ne0$ for every $j\in\F_p$.
The calculation of $\prod_{j\ne0}R_j$ above depends only on
$\tau$ and therefore also gives
\begin{equation}\label{eq:boundary-R-product-one}
 \prod_{j\ne0}R_j=1.
\end{equation}
For the Gaussian prefactors in $\prod_jT_j$, put
\[
 A_{\mathrm{prod}}=
 \prod_{j\in\F_p}\eta_0(j+d_0)
 =\eta_0^p(d_0^p-d_0),
 \qquad
 A_K=-\eta_0(d_0^p-d_0).
\]
Then
\[
 \frac{A_{\mathrm{prod}}}{A_K}=-\eta_0^{p-1}.
\]
Since $\eta_0^{p-1}$ is a square,
\[
 \nu(A_{\mathrm{prod}})=\nu(-1)\nu(A_K).
\]
Together with $q_0^p=q_0\nu(-1)$ from
Lemma~\ref{lem:q0-powers}, this gives
\begin{equation}\label{eq:boundary-Gaussian-prefactors}
 q_0^p\nu(A_{\mathrm{prod}})
 =q_0\nu(-1)^2\nu(A_K)
 =q_0\nu(A_K).
\end{equation}
Thus the complete Gaussian prefactors agree, although their quadratic
character parts alone differ by $\nu(-1)$.  Put
$e=\gamma-\gamma_0$.  The additive part is
\begin{align*}
 \sum_{j\in\F_p}\frac{B_j^2}{A_j}
 &=\eta_0\sum_j
   \frac{(\gamma_0(j+d_0)+d_0e)^2}{j+d_0}\\
 &=\eta_0d_0^2e^2\sum_j\frac1{j+d_0}
 =\frac{\eta_0d_0^2e^2}{d_0-d_0^p}.
\end{align*}
The first two expanded sums vanish because
$\sum_j(j+d_0)=0$ and $p=0$ in $k$; the last equality is the logarithmic
derivative of $\prod_j(X+j)=X^p-X$.  For the upper pair $(A,B)$,
\[
 \frac{B^2}{A}=\frac{\eta_0d_0^2e^2}{d_0-d_0^p}.
\]
Thus $U=\prod_jT_j$.  Together with
\eqref{eq:boundary-R-product-one}, this proves \eqref{eq:Rodd} in every
parity case.
\end{proof}

\subsubsection{Exact stationary-factor assembly}

We now assemble the complete stationary factors.  Let
$b=m_F(\chi_F)-1$ and put
\begin{equation}\label{eq:a-def}
 a=t-b.
\end{equation}
In the low range $b\le t$, use the notation of
\cref{prop:low-parameter-table} and set
\begin{equation}\label{eq:A-u-low}
 A=\alpha/\delta,
 \qquad u=\eps_1\beta_1/\alpha_1,
 \qquad n=N(u)=\eps\beta/\alpha.
\end{equation}
In the high range $b>t$, use \cref{prop:high-parameter-table} and set
\begin{equation}\label{eq:A-u-high}
 A=\alpha/\gamma_F,
 \qquad u=\beta_1/\alpha_1,
 \qquad n=N(u)=\beta/\alpha.
\end{equation}
In both cases
\begin{equation}\label{eq:a-valuations}
 v_K(u)=v_F(n)=a.
\end{equation}
The character linearizations needed for the correction units will be
proved only after their exact depths have been established.

Represent $\F_p^\times$ by its Teichmuller lifts.  Define
\begin{equation}\label{eq:correction-units}
 u^\vee=-n/u,
 \qquad
 z_0=\frac{N(1+u^\vee)}{1+N(u^\vee)},
 \qquad
 z_j=\frac{N(u+j)}{n+j}\quad(j\ne0).
\end{equation}
If $a\ne0$, the denominators are nonzero by valuation.  If $a=0$,
then $m_F(\chi_F)=T$.  On the last layer the stationary class of
$\tau^j\chi_F$ is represented, up to the common nonzero scalar, by
$n+j$.  Lemma~\ref{lem:minimal-orbit-stationary}, equivalently the
noncancellation criterion
\eqref{eq:stationary-leading-cancellation}, therefore gives
$\bar n+j\ne0$ for every $j\in\F_p$ (for $j=0$ this also follows from
$n\in\cO_F^\times$).  Thus $\bar n\notin-\F_p$; in particular every
$n+j$ and $1-n^{p-1}$ is a unit.  Put $x_0=z_0-1$ and $x_j=z_j-1$.

\begin{lemma}[Parameter-ratio identities]\label{lem:parameter-ratios}
One has
\begin{equation}\label{eq:parameter-ratios}
 N(u+j)=(n+j)z_j,
 \qquad
 N(n-u)=\left(\prod_{j\in\F_p}(n+j)\right)z_0.
\end{equation}
\end{lemma}

\begin{proof}
The first identity is the definition.  Since $p$ is odd,
$N(u^\vee)=-n^{p-1}$ and
\[
 N(1+u^\vee)=-\frac{N(n-u)}n.
\]
Also $\prod_j(n+j)=n^p-n$.  Substitution in the definition of $z_0$
gives the second identity.
\end{proof}

\begin{lemma}[Depth of the correction units]\label{lem:correction-depth}
Let
\[
 R(c)=\floor{\frac{c+(p-1)T}{p}}\qquad(c\ge0).
\]
Then
\begin{equation}\label{eq:zj-depth}
 z_j-1\in\p_F^{R(|a|)}\subseteq\p_F^{s_0}
 \qquad(j\ne0),
\end{equation}
and
\begin{equation}\label{eq:z0-depth}
 z_0-1\in\p_F^{R((p-1)|a|)}
 \subseteq\p_F^{\ceil{m/2}}.
\end{equation}
Thus $z_j-1$ lies in the linearization domain of $\tau$ and
$z_0-1$ lies in the linearization domain of $\chi_F$.
\end{lemma}

\begin{proof}
For $w\in K$ with $1+N(w)\ne0$, put
$Z(w)=N(1+w)/(1+N(w))$.  If $v_K(w)=c>0$, then
\[
 Z(w)-1=\frac{s_1(w)+\cdots+s_{p-1}(w)}{1+N(w)},
\]
so \cref{lem:wild-symmetric-bound} gives
$v_F(Z(w)-1)\ge R(c)$.  If $c<0$, the exact identity
$Z(w)=Z(w^{-1})$ reduces to the positive case.  If $c=0$, the
minimal-conductor condition used above says that $1+N(w)$ is a unit in
each occurrence of $Z(w)$: for $w=u/j$ this is equivalent to
$\bar n+j\ne0$, and for $w=u^\vee$ to $1-\bar n^{p-1}\ne0$.
The strong symmetric estimate then gives
\[
 v_F\bigl(s_i(w)\bigr)\ge
 \floor{\frac{(p-1)T}{p}}=R(0)
 \qquad(1\le i\le p-1),
\]
so the same conclusion holds when $c=0$.

Now $z_j=Z(u/j)$, and $z_0=Z(u^\vee)$ with
$v_K(u^\vee)=(p-1)a$.  This proves the first inclusions.

We use the elementary integer inequality
\begin{equation}\label{eq:odd-half-depth-integer}
 \floor{\frac{(p-1)M}{p}}\ge\ceil{\frac M2}
 \qquad(p\ge3,\ M\ge2).
\end{equation}
If $M$ is even, the real inequality
$(p-1)M/p\ge M/2$ already has an integral right side.  If $M$ is odd,
then $M\ge3$, and
\[
 \frac{(p-1)M}{p}\ge\frac{M+1}{2}
 \quad\Longleftrightarrow\quad
 (p-2)M\ge p,
\]
which follows from $M\ge3$ and $p\ge3$.

Since $R(|a|)\ge R(0)$, applying
\eqref{eq:odd-half-depth-integer} with $M=T$ gives
$R(|a|)\ge\ceil{T/2}=s_0$.  For $z_0$, if $a\ge0$, then
$m=T-a$ and
\[
 R((p-1)a)
 =\floor{\frac{(p-1)(T+a)}p}
 \ge\ceil{\frac{T+a}{2}}
 \ge\ceil{\frac{T-a}{2}}
 =\ceil{\frac m2}.
\]
If $a<0$, put $c=-a$, so $m=T+c$; applying
\eqref{eq:odd-half-depth-integer} with $M=T+c$ gives
\[
 R((p-1)c)
 =\floor{\frac{(p-1)(T+c)}p}
 \ge\ceil{\frac{T+c}{2}}
 =\ceil{\frac m2}.
\]
This proves \eqref{eq:z0-depth}.
\end{proof}


\begin{proposition}[Correction-depth character linearization]
\label{prop:correction-linearization}
With the notation of
\eqref{eq:A-u-low}--\eqref{eq:correction-units}, one has
\begin{equation}\label{eq:two-linearizations}
 \chi_F(z_0)=\psi_F\bigl(An(z_0-1)\bigr),
 \qquad
 \tau^j(z_j)=\psi_F\bigl(Aj(z_j-1)\bigr)
 \quad(j\ne0).
\end{equation}
\end{proposition}

\begin{proof}
The stationary linearization of $\tau$ is valid on $\p_F^{s_0}$, and
that of $\chi_F$ is valid on
$\p_F^{\lceil m/2\rceil}$.  Lemma~\ref{lem:correction-depth} puts
$z_j-1$ and $z_0-1$ in precisely these two domains.  The coefficients of
the two linearizations are $A$ and $An$, respectively.  Applying the
linearizations and taking the $j$th power of the first identity gives
\eqref{eq:two-linearizations}.
\end{proof}

\begin{proposition}[Exact assembly of the complete stationary quotient]
\label{prop:exact-odd-assembly}
With the orientation \eqref{eq:error-term}, substitution of the complete
Lamprecht factors gives
\begin{equation}\label{eq:exact-assembly-odd}
 \Err_{K/F}(\chi_F)
 =\mathcal R_{\mathrm{odd}}\,
 \left[
 \psi_F(A \Tr_{K/F}(u))\chi_F(z_0)
 \prod_{j\ne0}\tau^j(z_j)
 \right]^{-1}.
\end{equation}
Consequently
\begin{equation}\label{eq:error-psiX}
 \Err_{K/F}(\chi_F)
 =\mathcal R_{\mathrm{odd}}\,\psi_F(-AX),
\end{equation}
where
\begin{equation}\label{eq:X-def}
 X=\Tr_{K/F}(u)+n(z_0-1)+\sum_{j\ne0}j(z_j-1).
\end{equation}
\end{proposition}

\begin{proof}
Write every local constant by
\eqref{eq:Delta-factorization-new}.  Corollary
\ref{cor:stationary-class-under-norm} identifies the pullback class,
Lemma~\ref{lem:stationary-class-calculus}(c) identifies every twist
class, and Propositions~\ref{prop:high-parameter-table} and
\ref{prop:low-parameter-table} give the displayed representatives of
those classes.  At the boundary $m=T$, Lemma
\ref{lem:minimal-orbit-stationary} proves that no leading twist class
cancels, so every representative below is used at its actual conductor.
Choose the normalized representatives displayed below.  Whenever the
conductor is odd, translate the associated critical function as
prescribed by Lemma~\ref{lem:complete-stationary-invariant}.  Thus the
calculation uses no canonical field element and discards neither an
elementary factor nor a Hasse correction.

In the high range $b>t$, the admissible-element and stationary-parameter
pairs are
\[
 \begin{array}{c|c}
 \chi_K&(\gamma_F,\ \alpha(n-u)),\\
 \tau^j&(\gamma_F/\alpha,\ j)\quad(j\ne0),\\
 \tau^j\chi_F&(\gamma_F,\ \alpha(n+j))\quad(j\in\F_p).
 \end{array}
\]
In the low range $b\le t$, they are
\[
 \begin{array}{c|c}
 \chi_K&(\delta/\eps_1,\ (\alpha/\eps_1)(n-u)),\\
 \tau^j&(\delta,\ \alpha j)\quad(j\ne0),\\
 \chi_F&(\delta/\eps,\ (\alpha/\eps)n),\\
 \tau^j\chi_F&(\delta,\ \alpha(n+j))\quad(j\ne0).
 \end{array}
\]
Each ordered pair consists of an admissible element and the selected
representative of the corresponding stationary numerator class.  The
tables are precisely the high- and low-conductor stationary-class tables
with their common scalar made visible; their use is legitimate because
the complete stationary factor is representative-independent.

The factors contributed by the admissible elements cancel exactly.  In
the high range their quotient is
\[
 \frac{\chi_K(\gamma_F)
       \prod_{j\ne0}\tau^j(\gamma_F/\alpha)}
      {\prod_j(\tau^j\chi_F)(\gamma_F)}=1,
\]
and in the low range it is
\[
 \frac{\chi_K(\delta/\eps_1)\prod_{j\ne0}\tau^j(\delta)}
      {\chi_F(\delta/\eps)
       \prod_{j\ne0}(\tau^j\chi_F)(\delta)}=1.
\]
Here $N(\eps_1)=\eps$ and $\alpha$ is a norm.

After the common scalar is cancelled, the multiplicative stationary
values contribute
\begin{equation}\label{eq:H-mult}
 H=\chi_F\!\left(\frac{\prod_j(n+j)}{N(n-u)}\right)
   \prod_{j\ne0}\tau^j\!\left(\frac{n+j}{j}\right).
\end{equation}
By Lemma~\ref{lem:parameter-ratios}, the first ratio is $z_0^{-1}$.
Moreover
\[
 \frac{n+j}{j}=N\!\left(\frac{u+j}{j}\right)z_j^{-1},
\]
so the norm character kills the norm and
\begin{equation}\label{eq:H-exact}
 H=\chi_F(z_0)^{-1}\prod_{j\ne0}\tau^j(z_j)^{-1}.
\end{equation}

The additive stationary values have signed sum
\[
 \Tr_{K/F}(n-u)+\sum_{j\ne0}j-
 \sum_{j\in\F_p}(n+j)=-\Tr_{K/F}(u),
\]
using the Teichmuller sum in
Lemma~\ref{lem:teichmueller-stationary}.  Their additive factor is
therefore $\psi_F(-A\Tr_{K/F}(u))$.  Together with
\eqref{eq:H-exact} and the complete residual quotient
$\mathcal R_{\mathrm{odd}}$, this is precisely
\eqref{eq:exact-assembly-odd}.

Finally Proposition~\ref{prop:correction-linearization} evaluates all
correction characters at their now-verified depths.  Substitution gives
\eqref{eq:error-psiX} and \eqref{eq:X-def}.
\end{proof}

\subsubsection{The norm--trace formula and its valuation}

Let $s_r=s_r(u)$ be the elementary symmetric functions of the
conjugates of $u$, with $s_1=\Tr_{K/F}(u)$ and $s_p=n$.  Put
\begin{equation}\label{eq:B-def}
 B=s_1+\sum_{r=1}^{p-2}(-n)^r s_{p-r}.
\end{equation}

\begin{proposition}[Closed norm--trace formula]\label{prop:X-closed}
The element $X$ in \eqref{eq:X-def} satisfies the exact identity
\begin{equation}\label{eq:X-closed}
 X=\frac{pB}{1-n^{p-1}}.
\end{equation}
\end{proposition}

\begin{proof}
For $j\ne0$,
\[
 x_j=\frac{N(u+j)-(n+j)}{n+j}
 =\frac{\sum_{r=1}^{p-1}j^{p-r}s_r}{n+j}.
\]
Because $j^{p-1}=1$,
\begin{equation}\label{eq:inverse-nj}
 \frac1{n+j}=
 \frac{\sum_{q=0}^{p-2}(-n)^qj^{p-2-q}}
      {1-n^{p-1}}.
\end{equation}
Indeed, multiplying the numerator by $n+j$ telescopes to
$1-n^{p-1}$.  Summing over the $(p-1)$ Teichmuller roots and using
\[
 \sum_{j\ne0}j^r=
 \begin{cases}p-1,&p-1\mid r,\\0,&p-1\nmid r,
 \end{cases}
\]
shows that only the indices $r+q=1$ and $r+q=p$ survive.  Hence
\begin{equation}\label{eq:sum-jxj}
 \sum_{j\ne0}jx_j=\frac{(p-1)B}{1-n^{p-1}}.
\end{equation}

On the other hand,
\[
 N(n-u)-(n^p-n)=\sum_{r=1}^{p-1}(-1)^rs_rn^{p-r}.
\]
Since $z_0=N(n-u)/(n^p-n)$,
\begin{align*}
 nx_0
 &=-\frac{\sum_{r=1}^{p-1}(-1)^rs_rn^{p-r}}
          {1-n^{p-1}}\\
 &=\frac{s_1n^{p-1}+\sum_{q=1}^{p-2}(-n)^qs_{p-q}}
          {1-n^{p-1}}\\
 &=\frac{B}{1-n^{p-1}}-s_1.
\end{align*}
Adding $s_1$, this identity, and \eqref{eq:sum-jxj} proves
\eqref{eq:X-closed}.  This is an identity in the field $F$ itself.  In
equal characteristic the same finite sums have $p=0$ and $p-1=-1$ in
$F$, so no lifting to characteristic zero is being used.
\end{proof}

\begin{lemma}[The mixed-characteristic $p$-valuation forced by ramification]
\label{lem:p-valuation}
Assume $\operatorname{char}F=0$.  Put $e=v_F(p)$ and
\[
 e_0=\floor{\frac{(p-1)T}{p}}=T-\ceil{T/p}.
\]
Then $e\ge e_0$ and, for odd $p$, $2e_0\ge T$.
\end{lemma}

\begin{proof}
Since $p=\Tr_{K/F}(1)$, the trace-ideal formula in
Theorem~\ref{thm:norm-filtration}(b) gives $e\ge e_0$.
For $p\ge3$ and $T\ge2$, $\ceil{T/p}\le T/2$, hence
$2e_0=2T-2\ceil{T/p}\ge T$.
\end{proof}

\begin{proposition}[Vanishing or depth of the norm--trace error]
\label{prop:X-depth}
One has
\begin{equation}\label{eq:X-depth-final}
 X\in\p_F^T,
 \qquad \psi_F(A\p_F^T)=1.
\end{equation}
More precisely, if $\operatorname{char}F=p$, then $X=0$; if
$\operatorname{char}F=0$, the first inclusion follows from the
ramification valuation of $p$.  In either case
$\psi_F(-AX)=1$.
\end{proposition}

\begin{proof}
The denominator $1-n^{p-1}$ in \eqref{eq:X-closed} is nonzero.  If
$v_F(n)\ne0$, this follows from valuations; if $v_F(n)=0$, Lemma~\ref{lem:minimal-orbit-stationary} and
\eqref{eq:stationary-leading-cancellation} show that its residue is
nonzero, so it is a unit.

Suppose first that $\operatorname{char}F=p$.  In the exact identity
\eqref{eq:X-closed}, the integer $p$ is zero in $F$.  Therefore
\[
 X=\frac{pB}{1-n^{p-1}}=0,
\]
which proves the first assertion in equal characteristic.

Now assume $\operatorname{char}F=0$.  Let
$a=v_K(u)=v_F(n)$ and use the notation $e,e_0$ of
Lemma~\ref{lem:p-valuation}.

If $a>0$, the denominator in \eqref{eq:X-closed} is a unit.  By
\cref{lem:wild-symmetric-bound}, every $s_j(u)$ has valuation at least
$e_0$, and the powers of $n$ only increase valuation.  Hence
$v_F(pB)\ge e+e_0\ge2e_0\ge T$.

If $a=0$, the denominator is again a unit by
Lemma~\ref{lem:minimal-orbit-stationary} and
\eqref{eq:stationary-leading-cancellation}, and the same estimate
applies.

If $a<0$, put $c=-a>0$, $y=u^{-1}$, and for
$1\le r\le p-1$ set
\begin{equation}\label{eq:Lrc}
 L_r(c)=\floor{\frac{rc+(p-1)T}{p}}.
\end{equation}
Then $v_F(1-n^{p-1})=(p-1)a$, so division by the denominator is
multiplication, up to a unit, by $n^{1-p}$.  The reciprocal
identity is
\begin{equation}\label{eq:reciprocal-symmetric}
 s_j(u)=n\,s_{p-j}(y)\qquad(1\le j\le p-1).
\end{equation}
For the $s_1$ term,
\[
 v_F\bigl(pn^{1-p}s_1(u)\bigr)
 =v_F\bigl(pn^{2-p}s_{p-1}(y)\bigr)
 \ge e+(p-2)c+L_{p-1}(c)
 \ge e+e_0\ge T.
\]
For a term $n^{p-j}s_j(u)$ with $2\le j\le p-1$,
\[
 v_F\bigl(pn^{1-j}s_j(u)\bigr)
 =v_F\bigl(pn^{2-j}s_{p-j}(y)\bigr)
 \ge e+(j-2)c+L_{p-j}(c)
 \ge e+e_0\ge T.
\]
Thus $X\in\p_F^T$ also in mixed characteristic.

Finally $\tau$ has conductor $T$ and
$\tau(1+x)=\psi_F(Ax)$ on a layer containing $\p_F^T$.  Since
$\tau$ is trivial on $1+\p_F^T$, one has
$\psi_F(A\p_F^T)=1$.  The conclusion follows.
\end{proof}

\begin{corollary}[Wild odd-prime non-stable case]
\label{cor:wild-odd}
For odd $p$, every non-stable wildly ramified case satisfies
$\Err_{K/F}(\chi_F)=1$.
\end{corollary}

\begin{proof}
Combine \cref{prop:odd-phase-cancel,prop:exact-odd-assembly,prop:X-depth}.
\end{proof}

\subsection{Wild quadratic extensions}\label{sec:wild-quadratic}

Assume $p=2$ and $K/F$ is wildly ramified quadratic.  Let
$S(K/F)=\{1,\tau\}$.  The stationary calculation in this subsection is
made under the standing hypotheses that $\chi_F$ has minimal conductor
in its $S(K/F)$-orbit and that $m_F(\chi_F)>1$.  Proposition
\ref{prop:wild-quadratic-cancellation} removes these restrictions by
twist invariance and the endpoint calculation.  Under these standing
hypotheses, keep
\[
 b=m_F(\chi_F)-1,
 \qquad a=t-b=T-m_F(\chi_F),
 \qquad T=t+1.
\]
If $b\le t$, use the low-parameter notation and put
\[
 A=\frac{\alpha}{\delta},\qquad
 u=\frac{\eps_1\beta_1}{\alpha_1},\qquad
 n=N(u)=\frac{\eps\beta}{\alpha}.
\]
If $b>t$, use Proposition~\ref{prop:wild-quadratic-high-parameters} and put
\[
 A=\frac{\alpha\delta_h}{\gamma_F},\qquad
 u=\frac{\beta_1}{\alpha_1},\qquad
 n=N(u)=\frac{\beta}{\alpha\delta_h}.
\]
In both ranges $v_K(u)=v_F(n)=a$, and on every unit layer used below
\[
 \tau(1+z)=\psi_F(Az),\qquad
 \chi_F(1+z)=\psi_F(Anz).
\]
Let $u^\sigma$ be the nontrivial conjugate of $u$ and put
\begin{equation}\label{eq:quad-corrections}
 s=\Tr_{K/F}(u),\qquad u^\vee=-n/u=-u^\sigma,
\end{equation}
\[
 z_0=\frac{N(1+u^\vee)}{1+n},
 \qquad z_1=\frac{N(1+u)}{1+n},
 \qquad x=z_0-1,\quad y=z_1-1.
\]
If $a\ne0$, then $1+n\ne0$ by valuation.  If $a=0$, then
$m_F(\chi_F)=T$, and on the last layer the stationary class of
$\tau\chi_F$ is represented, up to the common nonzero scalar, by
$1+n$.  Lemma~\ref{lem:minimal-orbit-stationary} and
\eqref{eq:stationary-leading-cancellation} exclude its cancellation.
Thus $1+n$ is a unit.

\begin{lemma}[Exact quadratic norm correction]
\label{lem:quad-X}
One has
\begin{equation}\label{eq:quad-zxy}
 x=-\frac{s}{1+n},
 \qquad y=\frac{s}{1+n},
 \qquad
 X:=s+nx+y=\frac{2s}{1+n}.
\end{equation}
\end{lemma}

\begin{proof}
The exact quadratic norm identity gives
\[
 N(1+u)=1+s+n,
 \qquad N(1+u^\vee)=1-s+n.
\]
Division by $1+n$ gives the first two formulae.  Substituting them in
$s+nx+y$ gives $2s/(1+n)$.
\end{proof}

\begin{lemma}[Quadratic graded trace]
\label{lem:quadratic-graded-trace}
Let $K/F$ be wildly ramified quadratic and let
$T=v_K(\mathfrak D_{K/F})=t+1$.  If $r\ge0$ and $r+T$ is odd, put
$q=(r+T-1)/2$.  Then trace induces a $k$-linear isomorphism
\begin{equation}\label{eq:quadratic-graded-trace}
 \Tr_{K/F}:\p_K^r/\p_K^{r+1}
 \xrightarrow{\ \sim\ }
 \p_F^q/\p_F^{q+1}.
\end{equation}
In particular:
\begin{enumerate}[label=(\alph*)]
\item If $T$ is odd and $d_\tau=(T-1)/2$, then
$v_F(2)=d_\tau$.  For every $w\in\cO_K^\times$ and every lift
$\widetilde{\bar w}\in\cO_F$ of its residue,
\begin{equation}\label{eq:quadratic-unit-trace-class}
 \Tr(w)\equiv2\widetilde{\bar w}\pmod{\p_F^{d_\tau+1}}.
\end{equation}
\item If $m=2d+1>T$ and $r=m-T$, then trace identifies
$\p_K^r/\p_K^{r+1}$ with $\p_F^d/\p_F^{d+1}$.
\end{enumerate}
\end{lemma}

\begin{proof}
For a quadratic extension, \eqref{eq:trace-ideal} reads
\[
 \Tr(\p_K^j)=\p_F^{\floor{(j+T)/2}}.
\]
If $r+T=2q+1$, then
\[
 \Tr(\p_K^r)=\p_F^q,
 \qquad
 \Tr(\p_K^{r+1})=\p_F^{q+1}.
\]
Thus trace induces a surjection in
\eqref{eq:quadratic-graded-trace}.  Both source and target are
one-dimensional $k$-vector spaces, and the map is $k$-linear because
the residue extension is trivial.  It is therefore an isomorphism.

If $T$ is odd, apply the isomorphism with $r=0$.  The class of $1$ is
nonzero and maps to the class of $\Tr(1)=2$, proving
$v_F(2)=d_\tau$.  Since $w-\widetilde{\bar w}\in\p_K$, the trace-ideal formula
at $j=1$ gives
\[
 \Tr(w)-2\widetilde{\bar w}\in\Tr(\p_K)
 =\p_F^{(T+1)/2}=\p_F^{d_\tau+1},
\]
which proves (a).  In particular, a wildly ramified quadratic
extension with $T$ odd is necessarily of mixed characteristic; in
equal characteristic two the equality $v_F(2)=d_\tau$ is impossible.
If $m=2d+1$ and $r=m-T$, then $r+T=m$ and $q=d$, proving (b).
\end{proof}

\subsection{A normalized quadratic refinement in characteristic two}

Put
\[
 \psi_0(z)=(-1)^{\Tr_{k/\F_2}(z)}.
\]
An odd critical layer supplies a nowhere-zero unimodular function
$h:k\to\C^\times$ satisfying
\begin{equation}\label{eq:h-law}
 h(r+s)=h(r)h(s)\psi_0(rs).
\end{equation}
If no odd critical layer occurs, there is no Hasse factor and this
subsection is unnecessary.

\begin{lemma}[Normalized refinement]
\label{lem:normalized-refinement}
After multiplying $h$ by an additive character, one obtains a
unimodular function $\mathfrak q:k\to\C^\times$ such that
\begin{equation}\label{eq:q-refinement}
 \mathfrak q(r+s)=\mathfrak q(r)\mathfrak q(s)\psi_0(rs),
 \qquad
 \mathfrak q(r^2)=\mathfrak q(r).
\end{equation}
Every other function with the same polar character is uniquely
\begin{equation}\label{eq:q-linear-family}
 h_\lambda(z)=\mathfrak q(z)\psi_0(\lambda z)
 \qquad(\lambda\in k).
\end{equation}
If
\[
 H(\lambda)=\ph\!\left(\sum_{z\in k}h_\lambda(z)\right),
 \qquad H_0=H(0),
\]
then
\begin{equation}\label{eq:q-phase-translation}
 H(\lambda)=H_0\overline{\mathfrak q(\lambda)},
 \qquad H_0^2=\mathfrak q(1),
 \qquad \mathfrak q(r)^2=\psi_0(r).
\end{equation}
\end{lemma}

\begin{proof}
The quotient
\[
 \ell(z)=\frac{h(z^2)}{h(z)}
\]
is additive.  Indeed, using \eqref{eq:h-law} and Frobenius invariance
of $\psi_0$,
\[
 \ell(r+s)
 =\frac{h((r+s)^2)}{h(r+s)}
 =\frac{h(r^2)}{h(r)}\frac{h(s^2)}{h(s)}
 =\ell(r)\ell(s).
\]
The trace pairing on $k$ is perfect, so there is a unique $c\in k$
such that $\ell(z)=\psi_0(cz)$.  Let $f=[k:\F_2]$.  For every
$z\in k$, telescoping around the Frobenius orbit gives
\begin{align*}
 1
 &=\frac{h(z^{2^f})}{h(z)}
   =\prod_{i=0}^{f-1}\frac{h(z^{2^{i+1}})}{h(z^{2^i})}\\
 &=\psi_0\!\left(c\sum_{i=0}^{f-1}z^{2^i}\right)
   =\psi_0\!\left(c\Tr_{k/\F_2}(z)\right).
\end{align*}
Choose $z$ of absolute trace $1$.  Then $\psi_0(c)=1$, or equivalently
$\Tr_{k/\F_2}(c)=0$.

Consider the $\F_2$-linear map
\[
 L:k\longrightarrow k,
 \qquad L(d)=d+d^{2^{f-1}}.
\]
Its image is contained in the trace-zero hyperplane.  Its kernel is
$\F_2$: indeed, $d^{2^{f-1}}=d$ means that $d$ is fixed by
$\operatorname{Frob}^{f-1}$, whose fixed field in $\F_{2^f}$ is
$\F_{2^{\gcd(f,f-1)}}=\F_2$.  Thus $|\operatorname{im}L|=|k|/2$,
which is the size of the trace-zero hyperplane.  Hence there is
$d\in k$ such that
\begin{equation}\label{eq:quadratic-refinement-coboundary}
 c=d+d^{2^{f-1}}.
\end{equation}
Replace $h(z)$ by
\[
 \mathfrak q(z)=h(z)\psi_0(dz).
\]
The polar character is unchanged.  Moreover
\begin{align*}
 \frac{\mathfrak q(z^2)}{\mathfrak q(z)}
 &=\psi_0(cz)\psi_0(dz^2+dz)\\
 &=\psi_0\!\left((c+d^{2^{f-1}}+d)z\right)=1,
\end{align*}
where
$\psi_0(dz^2)=\psi_0(d^{2^{f-1}}z)$ and
\eqref{eq:quadratic-refinement-coboundary} were used.  This proves
$\mathfrak q(z^2)=\mathfrak q(z)$.  If $h'$ has the same polar
character as $\mathfrak q$, then $h'/\mathfrak q$ is an additive
character of $k$, hence is uniquely $z\mapsto\psi_0(\lambda z)$.
This proves \eqref{eq:q-linear-family}.

From \eqref{eq:q-refinement},
\[
 \mathfrak q(z)\psi_0(\lambda z)
 =\mathfrak q(z+\lambda)\mathfrak q(\lambda)^{-1},
\]
so translation of the finite sum gives the first formula in
\eqref{eq:q-phase-translation}.  Setting $r=s$ in
\eqref{eq:q-refinement} gives
\[
 \mathfrak q(r)^2=\psi_0(r^2)=\psi_0(r),
\]
and therefore
$\overline{\mathfrak q(r)}=\mathfrak q(r)\psi_0(r)$.

Put $S_0=\sum_{r\in k}\mathfrak q(r)$.  Its polar bicharacter is
$(r,s)\mapsto\psi_0(rs)$, which is nondegenerate by the perfect trace
pairing.  Lemma~\ref{lem:finite-quadratic-magnitude} therefore gives
\[
 |S_0|^2=|k|,
\]
so $S_0\ne0$.  Finally, translation by $1$ gives
\[
 \overline{S_0}
 =\sum_r\mathfrak q(r)\psi_0(r)
 =\frac1{\mathfrak q(1)}\sum_r\mathfrak q(r+1)
 =\frac{S_0}{\mathfrak q(1)}.
\]
Taking phases yields
$H_0^{-1}=H_0/\mathfrak q(1)$, and hence
$H_0^2=\mathfrak q(1)$.
\end{proof}


\begin{lemma}[Correction on an odd critical line]
\label{lem:quadratic-correction-rule}
Let $m_E(\theta)=2d+1$, choose
$c_\theta\in\p_E^d\setminus\p_E^{d+1}$, and write the corresponding
critical function in normalized characteristic-two form as
\[
 h_\lambda(z)=\mathfrak q(z)\psi_0(\lambda z).
\]
If $z_\theta=1+w$ with $w\in\p_E^d$ and
$c=\overline{w/c_\theta}\in k_E$, then
\begin{equation}\label{eq:quadratic-correction-rule}
 \theta(z_\theta)^{-1}
 =\psi_E\!\left(-\frac{\beta_\theta w}{\gamma_\theta}\right)
  h_\lambda(c).
\end{equation}
Thus the linear approximation and the complete value of the critical
function must be translated together.
\end{lemma}

\begin{proof}
Choose a lift $\widetilde c\in\cO_E$ of $c$ and write
\[
 w=c_\theta\widetilde c+r,
 \qquad r\in\p_E^{d+1}.
\]
Then
\[
 1+w=(1+c_\theta\widetilde c)(1+y),
 \qquad
 y=\frac{r}{1+c_\theta\widetilde c}\in\p_E^{d+1}.
\]
The odd-conductor stationary linearization gives
$\theta(1+y)=\psi_E(\beta_\theta y/\gamma_\theta)$.  By the definition
of the critical function,
\[
 \theta(1+c_\theta\widetilde c)^{-1}
 =\psi_E\!\left(-\frac{\beta_\theta c_\theta
                              \widetilde c}{\gamma_\theta}\right)
  h_\lambda(c).
\]
Consequently
\[
 \theta(1+w)^{-1}
 =\psi_E\!\left(-\frac{\beta_\theta
              (c_\theta\widetilde c+y)}{\gamma_\theta}\right)
  h_\lambda(c).
\]
Finally
\[
 c_\theta\widetilde c+y-w
 =-c_\theta\widetilde c\,y\in\p_E^{2d+1}.
\]
Since $\gamma_\theta\cO_E=\p_E^{2d+1+n_E(\psi_E)}$ and
$\beta_\theta$ is a unit, the additive character is trivial on the
resulting difference.  Replacing
$c_\theta\widetilde c+y$ by $w$ proves
\eqref{eq:quadratic-correction-rule}.
\end{proof}


\begin{proposition}[Branch matrix for characteristic-two critical refinements]
\label{prop:quadratic-refinement-branch-matrix}
Put $m=m_F(\chi_F)=b+1$ and $T=t+1$.  A normalized quadratic
refinement is needed exactly when an actual positive-polar critical
function occurs.  In particular, no refinement is chosen in precisely
the following three all-even rows:
\begin{equation}\label{eq:quadratic-all-even-rows}
 \begin{array}{c|c|c}
  \text{range}&m\pmod2&T\pmod2\\ \hline
  b<t&0&0\\
  b=t&0&0\\
  b>t&0&0.
 \end{array}
\end{equation}
In every other row, a refinement is attached only to the odd critical
source or sources that actually occur.  At the boundary the
noncancellation assertion $1+\rho\ne0$ is required only in the odd row
$b=t$, $T$ odd.
\end{proposition}

\begin{proof}
For even conductor, the finite stationary-layer function is the constant
$1$ by the even Lamprecht linearization and has no positive polar part.
For odd conductor, the critical function has a nondegenerate polar
character and is described by a normalized characteristic-two
refinement.  Checking the three possible ranges $b<t$, $b=t$, and
$b>t$ gives exactly the rows in
\eqref{eq:quadratic-all-even-rows}; all other parity rows retain at least
one odd lower or upper critical source.  At the boundary $m=T$, so the
only odd boundary is the common odd row.  There minimality excludes loss
of the final conductor layer and gives $1+\rho\ne0$.
\end{proof}

\subsection{The four critical functions}

By Proposition~\ref{prop:quadratic-refinement-branch-matrix}, no
quadratic refinement is present in an all-even row.  Whenever the
corresponding conductor is odd, let
\[
 \gamma_0,\qquad \gamma,\qquad \gamma_1,\qquad \gamma'
\]
be the residual linear coefficients in \eqref{eq:q-linear-family} for
$\tau,\chi_F,\tau\chi_F,\chi_K$, respectively.  If a conductor is
even, its finite stationary-layer function is the constant $1$ by the
even Lamprecht linearization, and the corresponding linear coefficient
is omitted.  Thus every formula below involving one of these
coefficients is asserted only in the parity row in which that
coefficient exists.  At the boundary
$b=t$, the element $u$ in \eqref{eq:quad-corrections} is a unit.  Put
\begin{equation}\label{eq:quadratic-rho-definition}
 \rho=\bar u\in k^\times.
\end{equation}
Then $n=N(u)$ has residue $\bar n=\rho^2$.  On the common
critical line, the exact stationary ratios are $A$ for $\tau$ and
$An$ for $\chi_F$.  Hence
Lemma~\ref{lem:critical-polar-coordinate}, applied in one common
critical coordinate, shows that their polar coefficients have ratio
$\bar n=\rho^2$.  We fix a lift $\widetilde\rho\in\cO_F^\times$
and take the $\chi_F$ critical element to be
$c_\chi=c_\tau/\widetilde\rho$ whenever the two lower functions
are compared.  Thus the normalized
$\chi_F$ variable is $z\mapsto\rho z$.  Put $r=\sqrt\rho$; the square
root is unique because Frobenius is an automorphism of $k$.

At the boundary, the critical coefficient of $\tau\chi_F$ is a
nonzero scalar times $1+\rho^2$.  It is nonzero because minimality
prevents $\tau\chi_F$ from losing its last conductor layer.  Since
$1+\rho^2=(1+\rho)^2$ in characteristic two, one has
$1+\rho\ne0$.

\begin{proposition}[Wild-quadratic critical functions]
\label{prop:quadratic-critical-functions}
There are phase-preserving coordinates for every critical function
that occurs.  The lower assertions below are made whenever the relevant
lower conductors are odd, and each assertion about the upper function
is made whenever the upper conductor is odd.  In an even-conductor row
the corresponding critical function is the constant $1$, and any
displayed relation involving its linear coefficient is omitted.
\begin{enumerate}[label=(\roman*)]
\item If $b<t$, the lower product has $\gamma_1=\gamma_0$, and before
the final Frobenius reindexing the upper critical function is
\[
 \mathfrak q(z)\psi_0(\gamma z).
\]
\item If $b>t$, the lower product has $\gamma_1=\gamma$, and the raw
upper function is
\[
 \mathfrak q(z)\psi_0((1+\gamma_0)z).
\]
\item If $b=t$, the two lower functions in the common coordinate are
\[
 \mathfrak q(z)\psi_0(\gamma_0z),
 \qquad
 \mathfrak q(\rho z)\psi_0(\rho\gamma z),
\]
and the raw upper function is
\begin{equation}\label{eq:raw-upper-quadratic-identities}
 \mathfrak q(z)\psi_0(\gamma z)\,
 \mathfrak q(\rho z)\psi_0(\rho\gamma_0z)\,
 \psi_0(\rho z).
\end{equation}
\end{enumerate}
The upper normalized coordinate is obtained from the displayed raw
coordinate by precomposition with Frobenius $z\mapsto z^2$.
\end{proposition}

\begin{proof}
Put $m=b+1$.  Corollary~\ref{cor:stationary-class-under-norm},
Proposition~\ref{prop:low-parameter-table}, and
Proposition~\ref{prop:wild-quadratic-high-parameters} identify the
stationary classes in the two ranges.  Choose the representatives
displayed in those propositions and translate every odd critical
function as in Lemma~\ref{lem:complete-stationary-invariant}.  With the
element $A$ fixed at the beginning of this subsection, the resulting
ratios are exact equalities in the fields:
\begin{equation}\label{eq:quadratic-unified-stationary-ratios}
 \frac{\beta(\chi_F)}{\gamma_F}=An,
 \qquad
 \frac{\beta_K}{\gamma_K}=A(n-u),
 \qquad n=N(u).
\end{equation}
For the selected representatives, in the low range these equalities are
\[
 \frac{\beta}{\delta/\varepsilon}
   =\frac{\alpha}{\delta}\frac{\varepsilon\beta}{\alpha},
 \qquad
 \frac{(\alpha/\varepsilon_1)(n-u)}{\delta/\varepsilon_1}
   =\frac{\alpha}{\delta}(n-u),
\]
and in the high range they are
\[
 \frac{\beta}{\gamma_F}
   =\frac{\alpha\delta_h}{\gamma_F}
      \frac{\beta}{\alpha\delta_h},
 \qquad
 \frac{\alpha\delta_h(n-u)}{\gamma_F}=A(n-u).
\]

We first establish the lower-function assertions, independently of
whether an upper critical function occurs.  Suppose $b<t$ and $T$ is
odd.  Put $d_\tau=(T-1)/2$ and choose
$c_\tau\in\p_F^{d_\tau}\setminus\p_F^{d_\tau+1}$.  Since
$m=b+1<T$,
\[
 d_\tau\ge\left\lceil\frac m2\right\rceil,
\]
so the stationary linearization of $\chi_F$ is valid on the
$\tau$-critical line.  Choose the sum representative for the
stationary class of $\tau\chi_F$.  Then, for $z\in k$,
\[
 \frac{\Phi_{\tau\chi_F}(z)}{\Phi_\tau(z)}
 =\psi_F(An c_\tau\widetilde z)
   \chi_F(1+c_\tau\widetilde z)^{-1}=1.
\]
Thus the complete lower critical functions agree and
$\gamma_1=\gamma_0$.

Suppose $b>t$ and $m$ is odd.  Put $d_\chi=(m-1)/2$ and choose
$c_\chi\in\p_F^{d_\chi}\setminus\p_F^{d_\chi+1}$.  Since $m>T$,
\[
 d_\chi\ge\left\lceil\frac T2\right\rceil,
\]
so the stationary linearization of $\tau$ is valid on the
$\chi_F$-critical line.  With the sum representative for
$\tau\chi_F$ one obtains
\[
 \frac{\Phi_{\tau\chi_F}(z)}{\Phi_{\chi_F}(z)}
 =\psi_F(Ac_\chi\widetilde z)
   \tau(1+c_\chi\widetilde z)^{-1}=1.
\]
Hence the complete lower critical functions agree and
$\gamma_1=\gamma$.

At the boundary $b=t$, all relevant lower conductors have the same
parity.  In the odd row, the common-coordinate convention
$c_\chi=c_\tau/\widetilde\rho$ gives the two lower functions
\[
 \mathfrak q(z)\psi_0(\gamma_0z),
 \qquad
 \mathfrak q(\rho z)\psi_0(\rho\gamma z),
\]
as asserted.

We now assume that the upper conductor is odd and compute its critical
function.  Let $c_K$ be an upper critical element.  For $Z\in k$, choose a
lift $\widetilde Z\in\cO_F$, put
\[
 x_Z=c_K\widetilde Z,
 \qquad P(y)=N(1+y)-1=\Tr(y)+N(y),
\]
and let $p_{\rm nat}(Z)$ and $q_{\rm nat}(Z)$ be the classes of
$P(x_Z)$ and $P(ux_Z)$ on the $\chi_F$- and $\tau$-critical lines,
respectively.  Thus, whenever the corresponding lower conductor is
odd,
\[
 P(x_Z)\equiv c_\chi\widetilde{p_{\rm nat}(Z)},
 \qquad
 P(ux_Z)\equiv c_\tau\widetilde{q_{\rm nat}(Z)}
\]
modulo the next critical ideal.  If an element lies in that next ideal,
its class is $0$; if a lower conductor is even, the corresponding
critical function is $1$ and its class is immaterial.

Let
\[
 e_\chi\equiv m_F(\chi_F)\pmod2,
 \qquad e_\tau\equiv T\pmod2,
 \qquad e_\chi,e_\tau\in\{0,1\},
\]
and define the lower finite functions by
\[
 \Phi_{\chi_F}(z)=
 \begin{cases}
  h_\gamma(z),&e_\chi=1,\\
  1,&e_\chi=0,
 \end{cases}
 \qquad
 \Phi_\tau(z)=
 \begin{cases}
  h_{\gamma_0}(z),&e_\tau=1,\\
  1,&e_\tau=0.
 \end{cases}
\]
In an even row these are literal constant functions: the even
stationary linearization cancels the additive and multiplicative
factors on the stationary layer.  Thus no undefined characteristic-two
critical function is being used.

Write $\Phi_K^{\rm nat}(Z)$ for the upper critical function in this
natural coordinate.  The defining upper function and
\eqref{eq:quadratic-unified-stationary-ratios} give
\begin{align*}
 \Phi_K^{\rm nat}(Z)
 &=\psi_F\!\left(A\Tr((n-u)x_Z)\right)
   \chi_F(1+P(x_Z))^{-1}\\
 &=\Phi_{\chi_F}(p_{\rm nat}(Z))
   \psi_F\!\left(-A\{\Tr(ux_Z)+nN(x_Z)\}\right)\\
 &=\Phi_{\chi_F}(p_{\rm nat}(Z))\psi_F(-AP(ux_Z)).
\end{align*}
For the second equality we used the exact stationary-layer identity
\[
 \Phi_{\chi_F}(p_{\rm nat}(Z))
 =\psi_F(AnP(x_Z))\chi_F(1+P(x_Z))^{-1}
\]
and cancelled $n\Tr(x_Z)$; the last equality uses
$nN(x_Z)=N(ux_Z)$.  Since
$1+P(ux_Z)=N(1+ux_Z)$ and $\tau$ is trivial on norms, the corresponding
exact stationary-layer identity is
\[
 \Phi_\tau(q_{\rm nat}(Z))=\psi_F(AP(ux_Z)).
\]
If $e_\tau=1$, Lemma~\ref{lem:normalized-refinement} gives
$\Phi_\tau(q)^2=\psi_0(q)$; if $e_\tau=0$, then
$\Phi_\tau\equiv1$.  Hence in both cases
\[
 \Phi_\tau(q)^{-1}
 =\Phi_\tau(q)\psi_0(q)^{e_\tau}.
\]
Therefore
\begin{equation}\label{eq:quadratic-exact-transport-natural}
 \Phi_K^{\rm nat}(Z)
 =\Phi_{\chi_F}(p_{\rm nat}(Z))
  \Phi_\tau(q_{\rm nat}(Z))
  \psi_0(q_{\rm nat}(Z))^{e_\tau}.
\end{equation}

We now calculate the two classes, including their scalars.  Write
$m_K=2d_K+1$ whenever the upper conductor is odd and take
$v_K(c_K)=d_K$.  The trace-ideal formula gives
\begin{equation}\label{eq:quadratic-critical-trace-bound}
 v_F(\Tr(c_K\widetilde Z))
 \ge \left\lfloor\frac{d_K+T}{2}\right\rfloor,
 \qquad
 v_F(N(c_K\widetilde Z))=d_K.
\end{equation}
Thus, in each case in which $P(c_K\widetilde Z)$ reaches a lower
critical line, the trace term lies in the next ideal and the initial
term is $N(c_K)\widetilde Z^{\,2}$.  Consequently its natural class is
\begin{equation}\label{eq:quadratic-natural-norm-classes}
 p_{\rm nat}(Z)=\lambda_\chi Z^2,
 \qquad
 \lambda_\chi=\overline{\frac{N(c_K)}{c_\chi}},
\end{equation}
whenever the $\chi_F$ line is reached.  Similarly, whenever the
$\tau$ line is reached and the trace term is one step deeper,
\begin{equation}\label{eq:quadratic-natural-tau-class}
 q_{\rm nat}(Z)=\lambda_\tau Z^2,
 \qquad
 \lambda_\tau=\overline{\frac{nN(c_K)}{c_\tau}}.
\end{equation}
Both displayed scalars are nonzero when defined.  Replacing $c_K$ by
$\xi c_K$, with $\xi\in\cO_K^\times$, multiplies both scalars by
$\bar\xi^{\,2}$.  Since Frobenius is an automorphism of $k$, we may
therefore normalize either nonzero scalar to $1$ without changing any
valuation or the ratio of the two scalars.

Put $a=T-m=v_K(u)$.  Suppose first that $b<t$.  Then $m_K=m$ and, in
the odd row, $d_K=d_\chi=(m-1)/2$.  Formula
\eqref{eq:quadratic-critical-trace-bound} and
\eqref{eq:quadratic-natural-norm-classes} show that
$p_{\rm nat}(Z)=\lambda_\chi Z^2$.  Scale $c_K$ so that
$\lambda_\chi=1$.  If $T$ is odd, put $d_\tau=(T-1)/2$.  Since the
odd conductors $m<T$ differ by at least two,
\[
 v_K(uc_K)=d_K+a=d_\tau+\frac a2\ge d_\tau+1.
\]
The norm and trace terms of $P(uc_K\widetilde Z)$ then lie in
$\p_F^{d_\tau+1}$, so $q_{\rm nat}(Z)=0$; if $T$ is even there is no
$\tau$-critical factor.  Define the raw variable by $z=Z^2$.  Then
\begin{equation}\label{eq:quadratic-raw-classes-below}
 (p(z),q(z))=(z,0)\qquad(b<t).
\end{equation}

Suppose next that $b>t$ and that the upper critical line occurs.  Then
$T$ is odd.  With $d_\tau=(T-1)/2$ and
$m_K=2m-T$, one has
\[
 d_K=m-d_\tau-1,
 \qquad
 v_K(uc_K)=a+d_K=d_\tau.
\]
The trace of $uc_K\widetilde Z$ lies in
$\p_F^{d_\tau+1}$, while its norm has valuation $d_\tau$; hence
\eqref{eq:quadratic-natural-tau-class} applies.  Scale $c_K$ so that
$\lambda_\tau=1$.  If $m$ is odd and $d_\chi=(m-1)/2$, then
\[
 d_K-d_\chi=\frac{m-T}{2}\ge1,
\]
so $P(c_K\widetilde Z)$ is beyond the $\chi_F$ critical line; if $m$
is even there is no such critical factor.  Again putting $z=Z^2$
gives
\begin{equation}\label{eq:quadratic-raw-classes-above}
 (p(z),q(z))=(0,z)\qquad(b>t).
\end{equation}

Finally suppose $b=t$.  Then $a=0$, $m_K=m=T$, and in the odd row
$d_K=d_\chi=d_\tau$.  Both initial terms are norm terms.  Equations
\eqref{eq:quadratic-natural-norm-classes} and
\eqref{eq:quadratic-natural-tau-class} give
\[
 p_{\rm nat}(Z)=\lambda_\chi Z^2,
 \qquad
 q_{\rm nat}(Z)=\lambda_\tau Z^2.
\]
The common lower coordinates were fixed by
$c_\chi=c_\tau/\widetilde\rho$, and $\bar n=\rho^2$.  Therefore
\begin{equation}\label{eq:quadratic-boundary-lambda-ratio}
 \frac{\lambda_\tau}{\lambda_\chi}
 =\bar n\,\overline{\frac{c_\chi}{c_\tau}}
 =\rho^2\rho^{-1}=\rho.
\end{equation}
Scale $c_K$ so that $\lambda_\chi=1$.  This single scaling preserves
\eqref{eq:quadratic-boundary-lambda-ratio}, and hence simultaneously
normalizes the two maps to
$p_{\rm nat}(Z)=Z^2$ and $q_{\rm nat}(Z)=\rho Z^2$.  With the raw
variable $z=Z^2$ one obtains
\begin{equation}\label{eq:quadratic-raw-classes-boundary}
 (p(z),q(z))=(z,\rho z)\qquad(b=t).
\end{equation}

Combining the three rows gives
\begin{equation}\label{eq:quadratic-residue-transport-table}
\begin{array}{c|c|c}
 &p(z)&q(z)\\ \hline
 b<t&z&0\\
 b>t&0&z\\
 b=t&z&\rho z.
\end{array}
\end{equation}
Define
$\Phi_{\chi_K}^{\rm raw}(z)=\Phi_K^{\rm nat}(z^{1/2})$; this is
well defined because Frobenius is bijective.  Substitution of
\eqref{eq:quadratic-residue-transport-table} into
\eqref{eq:quadratic-exact-transport-natural} gives
\begin{equation}\label{eq:quadratic-raw-transport}
 \Phi_{\chi_K}^{\rm raw}(z)=
 \begin{cases}
  \Phi_{\chi_F}(z),&b<t,\\[1mm]
  \Phi_\tau(z)\psi_0(z),&b>t,\\[1mm]
  \Phi_{\chi_F}(z)\Phi_\tau(\rho z)\psi_0(\rho z),&b=t.
 \end{cases}
\end{equation}
Indeed, in the first row $q(z)=0$, while in the second and third rows
an upper critical function can occur only when $T$ is odd, so
$e_\tau=1$.
The lower assertions were proved before the upper transport
calculation, and the upper formulae follow by inserting
$h_\lambda(z)=\mathfrak q(z)\psi_0(\lambda z)$ in
\eqref{eq:quadratic-raw-transport}.  Since the chosen raw variable is
$z=Z^2$, the normalized upper function is
\[
 \Phi_{\chi_K}(Z)=\Phi_{\chi_K}^{\rm raw}(Z^2),
\]
which is exactly precomposition of the raw function with Frobenius.
\end{proof}

\begin{proposition}[Quadratic coefficient and linear-residue table]
\label{prop:quadratic-gamma-table}
In each parity row in which the indicated coefficients exist, the
complete critical functions satisfy
\begin{equation}\label{eq:quadratic-gamma-table}
\begin{array}{c|c}
 b<t&\gamma_1=\gamma_0,
       \qquad (\gamma')^2=\gamma,\\[2mm]
 b>t&\gamma_1=\gamma,
       \qquad (\gamma')^2=1+\gamma_0,\\[2mm]
 b=t&\displaystyle
 \gamma_1=\frac{\gamma_0+\rho\gamma+r}{1+\rho},
 \qquad
 (\gamma')^2=
 \frac{\gamma+\rho\gamma_0+\rho+r}{1+\rho}.
\end{array}
\end{equation}
\end{proposition}

\begin{proof}
Only the boundary algebra requires calculation.  Multiplying the two
lower functions in Proposition~\ref{prop:quadratic-critical-functions} and using
\eqref{eq:q-refinement} gives
\[
 \mathfrak q((1+\rho)z)
 \psi_0\bigl((\gamma_0+\rho\gamma)z+\rho z^2\bigr).
\]
Since $\psi_0(\rho z^2)=\psi_0(rz)$, normalization by the coordinate
$(1+\rho)z$ gives the displayed formula for $\gamma_1$.

For the upper boundary function,
\eqref{eq:raw-upper-quadratic-identities} becomes
\[
 \mathfrak q((1+\rho)z)
 \psi_0\bigl((\gamma+\rho\gamma_0+\rho)z+\rho z^2\bigr).
\]
Replacing $\rho z^2$ by $rz$ and normalizing by $(1+\rho)z$ gives the
raw affine coefficient
\[
 L=\frac{\gamma+\rho\gamma_0+\rho+r}{1+\rho}.
\]
The upper normalized function is obtained by precomposition with
Frobenius.  If $L_0^2=L$, then
\[
 \mathfrak q(z^2)\psi_0(Lz^2)
 =\mathfrak q(z)\psi_0(L_0z),
\]
by \eqref{eq:q-refinement} and Frobenius invariance of the absolute
trace.  Thus Frobenius replaces $L$ by its unique square root
$\gamma'$, so $(\gamma')^2=L$.  The two nonboundary rows are the same
argument with $L=\gamma$ and $L=1+\gamma_0$.
\end{proof}

\subsection{Stationary-factor assembly}

The elementary stationary value and the odd critical phase are canonical
only as a product.  We therefore assemble the four complete stationary
factors before simplifying their residual phases.

\begin{proposition}[Wild-quadratic stationary assembly]
\label{prop:quadratic-complete-assembly}
Use the compatible stationary-class representatives of
Proposition~\ref{prop:wild-quadratic-high-parameters} in the high range
and of Proposition~\ref{prop:low-parameter-table} in the low range.  Use the
critical coordinates of Proposition~\ref{prop:quadratic-critical-functions},
and keep each elementary stationary value together with the critical
phase translated with it.  Then the complete error term has the
following form.

If $b<t$, or if $b+1$ is even, or if $b>t$ and $b+1$ is odd, then
\begin{equation}\label{eq:quadratic-complete-nonboundary}
 \Err_{K/F}(\chi_F)=\psi_F(-AX).
\end{equation}

At the only exceptional boundary, namely $b=t$ and $T=t+1$ odd, put
\begin{equation}\label{eq:quadratic-ABC}
 A_\rho=\frac{\rho}{1+\rho},\qquad
 B=\frac{\rho^2}{1+\rho^2},\qquad
 C=\frac{\rho}{1+\rho^2}.
\end{equation}
Then the complete elementary-plus-Hasse quotient is
\begin{equation}\label{eq:quadratic-boundary-error-formula}
 \Err_{K/F}(\chi_F)
 =\mathfrak q(A_\rho)\mathfrak q(B)\mathfrak q(C)\,\psi_0(a_1),
\end{equation}
where
\begin{equation}\label{eq:quadratic-a1}
 a_1=\gamma_1A_\rho+\gamma B+\gamma_0C+C+A_\rho.
\end{equation}
Here $C$ is the elementary correction
$\psi_F(-AX)=\psi_0(C)$, and the last term $A_\rho$ is the affine term
which occurs when the raw upper and lower Hasse quotients are compared.
\end{proposition}

\begin{proof}
We first assemble the elementary stationary factors without making any
linear approximation.  Corollary~\ref{cor:stationary-class-under-norm}
identifies the pullback class, Lemma
\ref{lem:stationary-class-calculus}(c) identifies the twist class, and
Propositions~\ref{prop:wild-quadratic-high-parameters} and
\ref{prop:low-parameter-table} provide simultaneous representatives.
At $m=T$, Lemma~\ref{lem:minimal-orbit-stationary} excludes cancellation
of the twist's leading class.  Choose these representatives and retain
the translated critical functions required by
Lemma~\ref{lem:complete-stationary-invariant}.  In the high range put
$c=\alpha\delta_h$; then $A=c/\gamma_F$, $\beta=cn$, and the four pairs
(admissible element, selected stationary representative) are
\begin{equation}\label{eq:quadratic-high-four-pairs}
 \begin{array}{c|c}
 \chi_K&(\gamma_F,\ c(n-u)),\\
 \tau&(\gamma_F/c,\ 1),\\
 \chi_F&(\gamma_F,\ cn),\\
 \tau\chi_F&(\gamma_F,\ c(n+1)).
 \end{array}
\end{equation}
In the low range the selected pairs are
\begin{equation}\label{eq:quadratic-low-four-pairs}
 \begin{array}{c|c}
 \chi_K&(\delta/\varepsilon_1,\ (\alpha/\varepsilon_1)(n-u)),\\
 \tau&(\delta/\alpha,\ 1),\\
 \chi_F&(\delta/\varepsilon,\ (\alpha/\varepsilon)n),\\
 \tau\chi_F&(\delta,\ \alpha(n+1)).
 \end{array}
\end{equation}
These are the representatives of the stationary classes already proved,
with their common scalar factors displayed.  They are not asserted to
be canonical field elements; the accompanying translated critical
functions make each complete stationary factor independent of the
selection.

The scalar and admissible-element contributions cancel exactly.  For
example, in the high range the quotient of the character values at the
admissible elements is $\tau(c)^{-1}$, whereas the quotient of the
scalar parts of the four selected stationary representatives is
$\tau(c)$.  In the
low range the $\chi_F$-parts cancel because
$N(\delta/\varepsilon_1)=\delta^2/\varepsilon$ and
$N(\alpha/\varepsilon_1)=\alpha^2/\varepsilon$, and the remaining
$\tau$-parts are trivial because $\alpha$ is a norm.

After these cancellations the multiplicative stationary values are
\begin{equation}\label{eq:quadratic-multiplicative-exact}
 \chi_F(z_0)^{-1}\tau(z_1)^{-1}.
\end{equation}
Indeed,
\[
 n-u=-u(1+u^\vee),\qquad
 N(n-u)=nN(1+u^\vee)=n(1+n)z_0,
\]
so the $\chi_F$-quotient is $\chi_F(z_0)^{-1}$.  Also
$N(1+u)=(1+n)z_1$ and $\tau$ is trivial on norms, whence
$\tau(1+n)=\tau(z_1)^{-1}$.  The signed sum of the additive contributions of the selected stationary
representatives in either \eqref{eq:quadratic-high-four-pairs} or
\eqref{eq:quadratic-low-four-pairs} is
\[
 \Tr_{K/F}(n-u)+1-n-(n+1)=-s.
\]
Consequently, before the odd critical corrections are inserted, the
complete quotient is
\begin{equation}\label{eq:quadratic-before-critical-corrections}
 R_H\,\psi_F(-As)\chi_F(1+x)^{-1}\tau(1+y)^{-1},
\end{equation}
where
\begin{equation}\label{eq:quadratic-raw-Hasse-quotient}
 R_H=
 \frac{H(\gamma')^{e_K}H(\gamma_0)^{e_\tau}}
      {H(\gamma)^{e_\chi}H(\gamma_1)^{e_{\tau\chi}}}
\end{equation}
and $e_\theta$ is $1$ or $0$ according as the conductor of $\theta$
is odd or even.

Suppose an odd critical factor
$h_\lambda(z)=\mathfrak q(z)\psi_0(\lambda z)$ receives a correction
$1+w$ whose normalized critical class is $c$.  Lemma
\ref{lem:quadratic-correction-rule} gives
\begin{equation}\label{eq:quadratic-residual-correction}
 \theta(1+w)^{-1}
 =\psi_E\!\left(-\frac{\beta_\theta w}{\gamma_\theta}\right)
  \mathfrak q(c)\psi_0(\lambda c).
\end{equation}
Thus the correction factor is $\mathfrak q(c)$, rather than its complex
conjugate; this orientation follows from the occurrence of
$\theta(\beta)^{-1}$ in Lamprecht's formula.  Inserting the two linear
parts in \eqref{eq:quadratic-before-critical-corrections} gives
\begin{equation}\label{eq:quadratic-elementary-assembly}
 \psi_F\bigl(-A(s+nx+y)\bigr)=\psi_F(-AX)
\end{equation}
by Lemma~\ref{lem:quad-X}.  It remains to evaluate $R_H$ and the two
finite correction values.

Let $\epsilon(r)\in\{0,1\}$ denote the parity of the integer $r$ and
put $m=b+1$.  The conductor formulas give
\begin{equation}\label{eq:quadratic-parity-table}
 \begin{array}{c|c|c|c|c}
 &e_K&e_\tau&e_\chi&e_{\tau\chi}\\ \hline
 b<t&\epsilon(m)&\epsilon(T)&\epsilon(m)&\epsilon(T)\\
 b=t&\epsilon(T)&\epsilon(T)&\epsilon(T)&\epsilon(T)\\
 b>t&\epsilon(T)&\epsilon(T)&\epsilon(m)&\epsilon(m).
 \end{array}
\end{equation}
Here the last row uses $m_K=2m-T$ and
$m_F(\tau\chi_F)=m$ from
Lemma~\ref{lem:minimal-orbit-stationary}.  In the middle row $b=t$,
\eqref{eq:minimal-critical} gives $m_K=T$, while
Lemma~\ref{lem:minimal-orbit-stationary} gives
$m_F(\tau\chi_F)=T$.

If $b<t$, Proposition~\ref{prop:quadratic-critical-functions} identifies
the upper critical function with that of $\chi_F$ and the lower product
with that of $\tau$; hence $R_H=1$.  Moreover the trace-ideal formula
gives
\[
 v_F(s)\ge \left\lfloor\frac{T-m+T}{2}\right\rfloor.
\]
This is at least one step beyond every odd critical line of $\chi_F$ and
$\tau$, because $m<T$.  Thus neither $x$ nor $y$ supplies a critical
correction, and the residual quotient is $1$.

Assume $b>t$.  If $m$ is even, the $\chi_F$ and $\tau\chi_F$ factors
are even.  If $T$ is also even, $R_H=1$.  If $T$ is odd, Frobenius
invariance of the finite sum and
$(\gamma')^2=1+\gamma_0$ give
\[
 R_H=H(1+\gamma_0)H(\gamma_0)=1.
\]
Indeed,
\[
 H(1+\gamma_0)H(\gamma_0)
 =H_0^2\overline{\mathfrak q(1+\gamma_0)}
       \overline{\mathfrak q(\gamma_0)}=1
\]
by \eqref{eq:q-refinement} and
$H_0^2=\mathfrak q(1)$.  The same trace bound, now using
$v_F(1+n)=T-m$, places $y$ at least one step beyond the $\tau$ critical
line, so there is no correction.  Thus the residual quotient is again
$1$.

Now let $b>t$ and $m$ be odd.  In both parities of $T$ one obtains
\begin{equation}\label{eq:quadratic-high-odd-raw}
 R_H=\frac1{\mathfrak q(1)\psi_0(\gamma)}.
\end{equation}
For $T$ even this is simply $H(\gamma)^{-2}$; for $T$ odd, the numerator
in \eqref{eq:quadratic-raw-Hasse-quotient} is
$H(1+\gamma_0)H(\gamma_0)=1$, and the denominator is again
$H(\gamma)^2=\mathfrak q(1)\psi_0(\gamma)$.  Put $c=m-T>0$ and $w=u^{-1}$.  Then $v_K(w)=c$ and
$c+T=m$ is odd.  Lemma~\ref{lem:quadratic-graded-trace}(b) shows that
\[
 c_\chi:=-\Tr(w)\in\p_F^{(m-1)/2}\setminus\p_F^{(m+1)/2}
\]
is a valid critical element for $\chi_F$.  We must also check that this
critical element is phase preserving.  Since
\[
 v_F(nc_\chi^2)=(T-m)+(m-1)=T-1,
\]
the map
\[
 \lambda:k\longrightarrow\C^\times,\qquad
 \lambda(z)=\psi_F(An c_\chi^2\widetilde z)
\]
is the nontrivial character induced by $\tau$ on
$U_F^{T-1}/U_F^T$.  Indeed, $T\ge2$ and hence
$T-1\ge\lceil T/2\rceil$, so the stationary linearization of $\tau$
is valid on $U_F^{T-1}$; nontriviality follows from the exact conductor
$T$ of $\tau$.

Put
\[
 r=\frac{w}{\Tr(w)}.
\]
Then $\Tr(r)=1$ and, since $N(w)=n^{-1}$,
\[
 N(r)=\frac{N(w)}{\Tr(w)^2}
     =\frac1{n c_\chi^2},
 \qquad
 N(r)^{-1}=n c_\chi^2.
\]
Moreover,
\[
 v_K(r)=v_K(w)-v_K(\Tr(w))
       =(m-T)-(m-1)=1-T.
\]
For every $z\in k$ and every lift $\widetilde z\in\cO_F$, the element
$1+(\widetilde z/N(r))r$ is a unit of $K$: indeed,
\[
 v_K\!\left(\frac{\widetilde z}{N(r)}r\right)
 \ge 2(T-1)+(1-T)=T-1>0.
\]
The exact quadratic norm identity therefore gives
\[
 \begin{aligned}
 N_{K/F}\!\left(1+\frac{\widetilde z}{N(r)}r\right)
 &=1+\frac{\widetilde z}{N(r)}\Tr(r)
      +\frac{\widetilde z^{\,2}}{N(r)^2}N(r)\\
 &=1+n c_\chi^2\bigl(\widetilde z+\widetilde z^{\,2}\bigr).
 \end{aligned}
\]
Since $\tau$ is trivial on norms, its stationary linearization gives
\[
 \lambda(z+z^2)=1 \qquad(z\in k).
\]
The image of the Artin--Schreier map $z\mapsto z^2+z$ is
$\ker\Tr_{k/\F_2}$.  Thus $\lambda$ is a nontrivial additive character
trivial on the trace-zero hyperplane.  There is exactly one such
character, namely
$z\mapsto(-1)^{\Tr_{k/\F_2}(z)}=\psi_0(z)$.  Consequently
\begin{equation}\label{eq:quadratic-chi-polar-normalization-nonboundary}
 \psi_F(An c_\chi^2\widetilde z)=\psi_0(z),
\end{equation}
so $c_\chi=-\Tr(w)$ is phase preserving and may be used in the
normalized critical-function calculation.  Since
\[
 \frac{s}{n}=\Tr(u^{-1})=\Tr(w)
\]
and $x=-s/(1+n)$, one has the exact identity
\[
 \frac{x}{c_\chi}=\frac{n}{1+n}
 =1-\frac1{1+n}.
\]
Here $v_F(n)=T-m=-c$, so $v_F((1+n)^{-1})=c>0$.  Therefore the
normalized correction class is
\begin{equation}\label{eq:quadratic-correction-residue-nonboundary}
 x'=\overline{x/c_\chi}=1.
\end{equation}
The correction in \eqref{eq:quadratic-residual-correction} is then
$\mathfrak q(1)\psi_0(\gamma)$, which is the inverse of
\eqref{eq:quadratic-high-odd-raw}.  Hence the residual quotient is $1$
in every nonboundary case.  Combining this with
\eqref{eq:quadratic-elementary-assembly} proves
\eqref{eq:quadratic-complete-nonboundary}.

It remains to assume $b=t$ and $T$ odd.  All four critical factors are
present.  Put
\[
 L=\frac{\gamma+\rho\gamma_0+\rho+r}{1+\rho},
 \qquad A_\rho=\frac{\rho}{1+\rho}.
\]
Since $(\gamma')^2=L$ and the finite sum is Frobenius-invariant,
$H(\gamma')=H(L)$.  Hence
\begin{align*}
 R_H
 &=\frac{\overline{\mathfrak q(L)}
          \overline{\mathfrak q(\gamma_0)}}
         {\overline{\mathfrak q(\gamma)}
          \overline{\mathfrak q(\gamma_1)}}\\
 &=\mathfrak q(\gamma)\mathfrak q(\gamma_1)
   \overline{\mathfrak q(L)}\overline{\mathfrak q(\gamma_0)}.
\end{align*}
Set $X_0=\gamma+\gamma_1$ and $Y_0=L+\gamma_0$.  In characteristic
two, substitution of the boundary formulae gives
\begin{align*}
 X_0+Y_0
 &=\gamma+\gamma_0
   +\frac{\gamma_0+\rho\gamma+r
          +\gamma+\rho\gamma_0+\rho+r}{1+\rho}\\
 &=\frac{\rho}{1+\rho}=A_\rho.
\end{align*}
Thus $X_0=Y_0+A_\rho$.  The refinement law therefore gives
\begin{equation}\label{eq:quadratic-boundary-raw-intermediate}
 R_H=\mathfrak q(A_\rho)\psi_0(E),
 \qquad
 E=Y_0A_\rho+\gamma\gamma_1+L\gamma_0.
\end{equation}
Let $a_\rho=r/(1+r)$, so $a_\rho^2=A_\rho$.  A direct expansion in
characteristic two gives
\begin{equation}\label{eq:quadratic-boundary-AS-difference}
 E+\gamma_1A_\rho+A_\rho
 =\bigl(a_\rho(\gamma+\gamma_0)+A_\rho\bigr)^2
  +\bigl(a_\rho(\gamma+\gamma_0)+A_\rho\bigr).
\end{equation}
The right side has absolute trace zero.  Thus
\begin{equation}\label{eq:quadratic-boundary-raw-Hasse}
 R_H=\mathfrak q(A_\rho)
      \psi_0(\gamma_1A_\rho+A_\rho).
\end{equation}

We now determine the two correction classes directly from the graded
trace.  Put $d_\tau=(T-1)/2$.  By
Lemma~\ref{lem:quadratic-graded-trace}(a), $v_F(2)=d_\tau$, so we may
choose
\begin{equation}\label{eq:quadratic-boundary-critical-elements}
 c_\tau=2,
 \qquad c_\chi=2/\widetilde\rho.
\end{equation}
Since $u-\widetilde\rho\in\p_K$, the same lemma gives
\[
 s-2\widetilde\rho=\Tr(u-\widetilde\rho)\in\p_F^{d_\tau+1}.
\]
Consequently
\begin{equation}\label{eq:quadratic-boundary-trace-class}
 \overline{s/c_\tau}=\rho.
\end{equation}
Moreover, our choice $c_\tau=2$ gives
\begin{equation}\label{eq:quadratic-boundary-two-class}
 \overline{2/c_\tau}=1.
\end{equation}
Together with $\bar n=\rho^2$, Lemma~\ref{lem:quad-X} gives
\begin{equation}\label{eq:quadratic-boundary-correction-classes}
 x'=\overline{x/c_\chi}=\frac{\rho^2}{1+\rho^2}=B,
 \qquad
 y'=\overline{y/c_\tau}=\frac{\rho}{1+\rho^2}=C.
\end{equation}
We next prove the exact normalization of the polar character; the
valuation of $c_\tau$ alone does not determine its residual scalar.
Since this is the odd boundary and the endpoint conductors have already
been removed, $T\ge3$.  Hence
$T-1\ge\lceil T/2\rceil$, so the stationary linearization of $\tau$
is valid on $4\cO_F$.  Since $c_\tau=2$ and $v_F(4)=T-1$, the map
\[
 \lambda:k\longrightarrow\C^\times,
 \qquad
 \lambda(z)=\psi_F(Ac_\tau^2\widetilde z)
            =\tau(1+4\widetilde z)
\]
is the character induced by $\tau$ on
$U_F^{T-1}/U_F^T$.  It is nontrivial because $\tau$ has exact
conductor $T$.  For every $z\in k$, the element $1+2\widetilde z$
lies in $F^\times\subset K^\times$, and hence
\[
 N_{K/F}(1+2\widetilde z)
 =(1+2\widetilde z)^2
 =1+4(\widetilde z+\widetilde z^{\,2}).
\]
The character $\tau$ is trivial on norms, so its stationary
linearization gives
\[
 \lambda(z+z^2)=1.
\]
The image of the Artin--Schreier map $z\mapsto z^2+z$ is exactly
$\ker\Tr_{k/\F_2}$: it is contained there, its kernel is $\F_2$, and
therefore its image has $|k|/2$ elements.  Thus $\lambda$ is a
nontrivial additive character trivial on the trace-zero hyperplane.
There is only one such character, namely
$z\mapsto(-1)^{\Tr_{k/\F_2}(z)}=\psi_0(z)$.  Consequently
\begin{equation}\label{eq:quadratic-tau-polar-normalization}
 \psi_F(Ac_\tau^2\widetilde z)=\psi_0(z)
 \qquad(z\in k).
\end{equation}
At the boundary $X=2y=c_\tau y$.  Since
$y/c_\tau\equiv C\pmod{\p_F}$, one has
$X\equiv c_\tau^2\widetilde C\pmod{\p_F^T}$ for any lift
$\widetilde C\in\cO_F$ of $C$; the character
$x\mapsto\psi_F(Ax)$ is trivial on $\p_F^T$.  Thus
\begin{equation}\label{eq:quadratic-boundary-elementary-class}
 \psi_F(-AX)=\psi_0(C),
\end{equation}
where the sign disappears after reduction to the residue field of
characteristic two.

The two factors from \eqref{eq:quadratic-residual-correction} are
$\mathfrak q(B)\psi_0(\gamma B)$ and
$\mathfrak q(C)\psi_0(\gamma_0C)$.  Multiplying them with
\eqref{eq:quadratic-boundary-raw-Hasse} and
\eqref{eq:quadratic-boundary-elementary-class} yields
\[
 \Err_{K/F}(\chi_F)
 =\mathfrak q(A_\rho)\mathfrak q(B)\mathfrak q(C)
  \psi_0(\gamma_1A_\rho+\gamma B+\gamma_0C+C+A_\rho),
\]
which is \eqref{eq:quadratic-boundary-error-formula} and
\eqref{eq:quadratic-a1}.
\end{proof}

\subsection{Final quadratic cancellation}

\begin{lemma}[Artin--Schreier parity of the quadratic break]
\label{lem:AS-break-parity}
Assume $\operatorname{char}F=2$ and let $K/F$ be ramified quadratic.
Then its unique lower ramification break $t$ is odd.  Consequently
$T=t+1$ is even.
\end{lemma}

\begin{proof}
Because $K/F$ is separable quadratic in characteristic two, it has an
Artin--Schreier presentation
\[
 K=F(y),\qquad y^2-y=a.
\]
Replacing $y$ by $y-c$ replaces $a$ by $a-(c^2-c)$.  Choose a
uniformizer $\pi_F$.  If the leading pole order
$r=-v_F(a)>0$ is even, write its leading residue coefficient as
$\bar d^{\,2}$, possible because the finite residue field is perfect,
and take a lift $d\in\cO_F$.  Subtracting
\[
 (d\pi_F^{-r/2})^2-d\pi_F^{-r/2}
\]
cancels the term of pole order $r$ and introduces only terms of
strictly smaller pole order.  Repeating this finite process gives an
equivalent equation for which either $v_F(a)\ge0$ or
\[
 v_F(a)=-t<0\quad\text{with }t\text{ odd}.
\]
The first alternative cannot give a ramified quadratic extension: its
reduction is either irreducible, in which case the extension is
unramified, or has a root, in which case Hensel's lemma (the derivative
is $-1$) removes the extension.  Thus the second alternative holds.

Normalize $v_K$ by $v_K(K^\times)=\mathbf Z$.  Since
$v_K|_F=2v_F$ and $v_K(y)<0$, the equation gives
\[
 2v_K(y)=v_K(a)=-2t,
 \qquad v_K(y)=-t.
\]
Hence
\[
 \pi_K=\pi_F^{(t+1)/2}y
\]
is a uniformizer of $K$.  For the nontrivial automorphism
$\sigma(y)=y+1$,
\[
 \sigma(\pi_K)-\pi_K=\pi_F^{(t+1)/2},
\]
whose $K$-valuation is $t+1$.  Therefore the lower break is exactly
$t$, and it is odd.
\end{proof}

\begin{lemma}[Depth outside the odd boundary]
\label{lem:quadratic-depth}
If $\operatorname{char}F=2$, then $X=0$.  If
$\operatorname{char}F=0$ and $a\ne0$, then $X\in\p_F^T$; the same is
true when $a=0$ and $T$ is even.  Hence $\psi_F(AX)=1$ in all these
cases.
\end{lemma}

\begin{proof}
In equal characteristic two, the exact identity
\eqref{eq:quad-zxy} gives
\[
 X=\frac{2s}{1+n}=0;
\]
the denominator is nonzero by the minimality condition stated before
Lemma~\ref{lem:quad-X}.

Assume now $\operatorname{char}F=0$ and put $e=v_F(2)$.  Since
$2=\Tr_{K/F}(1)$, the trace-ideal formula gives
$e\ge\floor{T/2}$.  Also
$v_F(s)\ge\floor{(a+T)/2}$.  If $a>0$, $1+n$ is a unit and
\[
 v_F(X)\ge e+\floor{\frac{a+T}{2}}\ge T.
\]
If $a<0$, put $c=-a$; then $v_F(1+n)=a=-c$ and
\[
 v_F(X)\ge e+\floor{\frac{T-c}{2}}+c
          =e+\floor{\frac{T+c}{2}}\ge T.
\]
If $a=0$, the bound is $e+\floor{T/2}$, which is at least $T$ when
$T$ is even.  Finally $\psi_F(A\p_F^T)=1$ because $\tau$ has conductor
$T$ and is linearized by $A$.
\end{proof}

\begin{proposition}[Wild quadratic cancellation]
\label{prop:wild-quadratic-cancellation}
For every wildly ramified quadratic extension $K/F$ of nonarchimedean
local fields and every quasi-character $\chi_F$ of $F^\times$,
\[
 \Err_{K/F}(\chi_F)=1.
\]
\end{proposition}

\begin{proof}
By Lemma~\ref{lem:twist-invariance}, replace $\chi_F$ by a member of
minimal conductor in its $S(K/F)$-orbit.  The cases
$m_F(\chi_F)=0$ and $m_F(\chi_F)=1$ were proved in
Subsection~\ref{sec:wild-endpoints}.  Hence assume
$m_F(\chi_F)>1$; all stationary parameters and assembly formulas of
this subsection are now defined.

Suppose first that $\operatorname{char}F=2$.  By
Lemma~\ref{lem:AS-break-parity}, $T$ is even, so the exceptional boundary
$b=t$, $T$ odd cannot occur.  Equation
\eqref{eq:quadratic-complete-nonboundary} therefore applies, and
Lemma~\ref{lem:quadratic-depth} gives $X=0$.  Hence
$\Err_{K/F}(\chi_F)=1$.

Assume now $\operatorname{char}F=0$.  Outside the boundary $b=t$, $T$
odd, equation \eqref{eq:quadratic-complete-nonboundary} and
Lemma~\ref{lem:quadratic-depth} give
$\Err_{K/F}(\chi_F)=\psi_F(-AX)=1$.

It remains to take $b=t$ and $T$ odd.  Use the notation
\eqref{eq:quadratic-ABC}.  Then $B+C=A_\rho$.  By
\eqref{eq:quadratic-boundary-error-formula},
\begin{equation}\label{eq:quadratic-boundary-error}
 \Err_{K/F}(\chi_F)
 =\mathfrak q(A_\rho)\mathfrak q(B)
  \mathfrak q(C)\psi_0(a_1).
\end{equation}
Since $B+C=A_\rho$, the refinement identity gives
\[
 \mathfrak q(B)\mathfrak q(C)
 =\mathfrak q(A_\rho)\psi_0(BC),
 \qquad BC=\frac{\rho^3}{1+\rho^4}.
\]
Since $\mathfrak q(A_\rho)^2=\psi_0(A_\rho)$,
\eqref{eq:quadratic-boundary-error} is $\psi_0(a_2)$ with
\[
 a_2=a_1+A_\rho+\frac{\rho^3}{1+\rho^4}.
\]
Substituting
$\gamma_1=(\gamma_0+\rho\gamma+r)/(1+\rho)$ cancels the terms in
$\gamma$ and $\gamma_0$ and gives
\begin{equation}\label{eq:quadratic-a2}
 a_2=
 \frac{r^7+r^3+r^2}{1+r^8}
 =\frac{\rho^3}{(1+\rho^2)^2}
  +\frac{\rho(1+r)}{(1+\rho)^2}.
\end{equation}
The first term on the right is the square of
$r^3/(1+\rho^2)$.  Frobenius invariance of $\psi_0$ therefore reduces
its phase to that of
\[
 a_4=\frac{r^3}{1+\rho^2}
     +\frac{\rho(1+r)}{(1+\rho)^2}
     =\frac{\rho}{1+\rho^2}
     =\frac{\rho}{(1+\rho)^2}.
\]
With $c=(1+\rho)^{-1}$, $a_4=c+c^2$, so
$\Tr_{k/\F_2}(a_4)=0$ and $\psi_0(a_4)=1$.  This proves the boundary
case and hence the proposition.
\end{proof}


\section{Completion of the First Main Lemma}\label{sec:first-main-completion}

\begin{proof}[Proof of Theorem~\ref{thm:first-main}]
Let $F$ be a nonarchimedean local field and let $K/F$ be cyclic of
prime degree $\ell$.  If $K/F$ is unramified, apply
Proposition~\ref{prop:unramified-new}.  If it is ramified and tame,
then it is totally ramified; the odd-degree and quadratic cases are
Propositions~\ref{prop:tame-odd-new} and
\ref{prop:tame-two-new}.

Suppose that $K/F$ is wild.  Then $\ell=p$, the residue
characteristic, and $K/F$ is totally ramified with one lower break
$t$.  By Lemma~\ref{lem:twist-invariance}, replace $\chi_F$ by a
member of minimal conductor in its $S(K/F)$-orbit.  Conductors $0$ and
$1$ are treated in Subsection~\ref{sec:wild-endpoints}.  For odd $p$,
the stable range is Proposition~\ref{prop:wild-stable-new}, and the
remaining non-stable range is Corollary~\ref{cor:wild-odd}.  For
$p=2$, the wild quadratic calculation is
Proposition~\ref{prop:wild-quadratic-cancellation}.  These cases
exhaust all cyclic extensions of prime degree.  Hence
$\Err_{K/F}(\chi_F)=1$, which is exactly
\eqref{eq:first-main}.
\end{proof}


\begin{thebibliography}{99}
\bibitem{DavenportHasse}
H.~Davenport and H.~Hasse,
\emph{Die Nullstellen der Kongruenzzetafunktionen in gewissen
zyklischen F\"allen},
J. Reine Angew. Math. \textbf{172} (1935), 151--182.

\bibitem{Serre}
J.-P.~Serre,
\emph{Local Fields},
Graduate Texts in Mathematics \textbf{67}, Springer, 1979.

\bibitem{Langlands}
R.~P.~Langlands,
\emph{On the Functional Equation of the Artin $L$-functions},
Yale University preprint (1970), IAS electronic edition.

\bibitem{Dwork}
B.~M.~Dwork,
\emph{On the Root Number in the Functional Equation of the Artin--Weil
$L$-Series}, Ph.D. thesis, Columbia University, 1954.
\end{thebibliography}

\end{document}
